% Exact arithmetic calculations for the intrinsic class-two quotient.
% Statements A.1--A.19 and A.24 onward below are owned by arithmetic_sections.
\subsection{The rank-forty-four residue complex}

We put $K_0=\mathbb Q(\zeta)$, reserving $\mathcal K$ for the
rank-forty-four commutator lattice. Set $S=\{3,\infty\}$ and
$X=\operatorname{Spec}\mathbb Z[1/3]$. Arithmetic Frobenius is used
throughout. For each $q\ne3$ choose a decomposition group and write
\[
 A_q=\mathcal K(-1),\qquad D_q=F_q-1\in\operatorname{End}_{\mathbb Z_3}(A_q).
\]
The subscript records the local Galois action; it does not change the
underlying lattice. Define
\[
 \mathcal R(A)=\left\{(x_q)_{q\ne3}: x_q\in A_q,
 \quad\forall n\quad x_q\in3^nA_q\text{ for almost all }q\right\},
 \qquad D(x)_q=D_qx_q.
\]
All quotients by $D\mathcal R(A)$ below use the actual image.

\begin{statement}{A.1}{Proved}{The intrinsic coefficient module and its residue complex}
The finite locally constant sheaves $\mathcal K_n=\mathcal K/3^n$
on $X$ satisfy
\[
 H^i(X,\mathcal K_n)=0\quad(i\ge2),\qquad
 |H^1(X,\mathcal K_n)|=3^{23n}|H^0(X,\mathcal K_n)|.
 \tag{A.1.1}
\]
Their degree-one cohomology groups are finite. For every finite set
$J$ of primes different from three there is a compatible distinguished
triangle in the derived category of abelian groups
\[
 R\Gamma(X,\mathcal K_n)\longrightarrow
 R\Gamma(X\setminus J,\mathcal K_n)\longrightarrow
 \bigoplus_{q\in J}R\Gamma(\mathbb F_q,\mathcal K_n(-1))[-1]
 \longrightarrow R\Gamma(X,\mathcal K_n)[1].
 \tag{A.1.2}
\]
Every $D_q$ is injective and has finite cokernel. Localization and
continuous cochains induce isomorphisms
\[
 H^2_{\rm cts}(G_{\mathbb Q},\mathcal K)
       \simeq\mathcal R(A)/D\mathcal R(A),
 \qquad
 H^1_{\rm cts}(G_{\mathbb Q},\mathcal K)
       \simeq H^1_{\rm cts}(G_S,\mathcal K).
 \tag{A.1.3}
\]
The latter group is finitely generated of rank $23$; its torsion is
$H^0(G_S,(\mathcal K\otimes\mathbb Q_3)/\mathcal K)$.
This is an exact residue calculation, including the contribution over
$\mathbb Z[1/3]$ in the degrees relevant to lifting sections.
\par\smallskip\noindent\emph{Manuscript comparison: \S51.16,
folios 733v--736v, scans 172--178, formulas (371)--(379).}
\end{statement}
\begin{proof}
Every coordinate cusp of the Fermat cubic is a flex. Its tangent line
cuts out three times that cusp; the ratio of two tangent lines has
divisor three times their difference. The classifying homomorphism
from degree-zero cusp divisors to $E$ therefore has image in $E[3]$.
The extension $0\to W\to V\to T\to0$ is killed by three, and
\[
 V\otimes\mathbb Q_3\simeq
 (W\otimes\mathbb Q_3)\oplus(T\otimes\mathbb Q_3)
 \tag{A.1.4}
\]
as continuous Galois representations. The comparison is the splitting
of the pullback along $[3]:E\to E$ followed by division by three on
the Tate module. Thus the integral lattice is not being replaced
without its explicit isogeny comparison.

Over $K_0$ all nine points of $E[3]$ are the nine cusps. Hence both
$W/3$ and $T/3$ have trivial action over $K_0$. The action on $V/3$
is an extension of trivial modules, so its finite image is a
three-group. The congruence kernel in $\operatorname{GL}(V)$ is
pro-three. Consequently the image on every $\mathcal K_n$ over
$G_{K_0,S}$ is a finite three-group. Any finite module for a finite
three-group has a composition series whose factors are trivial
$\mathbb F_3$ modules: a nonzero invariant exists on every nonzero
$\mathbb F_3$ representation, and induction applies to its quotient.

The ring $\mathbb Z[\zeta]$ is norm-Euclidean. Its units are $\mu_6$
and its unique prime above three is $\lambda=1-\zeta$. Kummer theory
therefore gives
\[
 H^1(G_{K_0,S},\mathbb F_3)
  =\langle\zeta,\lambda\rangle\simeq\mathbb F_3^2.
\]
Indeed valuations outside $\lambda$ of an allowed Kummer element are
multiples of three, and unique factorization removes them. The two
remaining classes are independent, first by valuation and then by
the unit group. Milne's Euler characteristic, I.5.1 \cite{MilneADT}, is $1/3$ for
this imaginary quadratic field and the trivial module $\mathbb F_3$.
Since $H^0$ has order three and $H^1$ order nine, $H^2$ has order one.
Devissage along the composition series just proved gives
$H^2(G_{K_0,S},\mathcal K_n)=0$. The order-two quotient
$\operatorname{Gal}(K_0/\mathbb Q)$ has exact invariants on
three-primary modules, so Hochschild--Serre gives the same vanishing
over $G_S$. Milne I.4.10(c) gives vanishing in degrees at least three,
and I.4.15 gives finiteness. Proposition II.2.9 identifies these
Galois groups with the cohomology of the locally constant sheaves
on $X$. In particular no inference from the cohomological dimension
of a maximal pro-three quotient is needed here.

Complex conjugation has signs $(4,6)$ on $V$: its permutation signs
on $\mathbb Z_3[D]/\mathbb Z_3\delta$ are $(5,3)$, the toric twist
interchanges them, and its signs on $T$ are $(1,1)$. The signs on
$\bigwedge^2V$ are $(21,24)$ and the quotient
$\bigwedge^2T=\mathbb Z_3(1)$ removes one minus sign. Thus the signs
on $\mathcal K$ are $(21,23)$. Milne I.5.1 over $\mathbb Q$ now gives
(A.1.1). The same computation, passage to the inverse limit, and
$H^0(G_S,\mathcal K\otimes\mathbb Q_3)=0$ give the asserted rank and
torsion description. Finite generation follows, for example, from
the finite number of generators of the relative pro-three Galois
group used in LM7 E.34; its cocycles are a closed submodule of a
finite power of $\mathcal K$. Alternatively the finite groups in
(A.1.1), the vanishing of $H^2$, and completeness give the same
conclusion by successive lifting of a finite basis modulo three.

The localization sequence and purity give (A.1.2). At a strict
Henselian trait of residue characteristic $q\ne3$, the wild inertia
is pro-$q$ and tame three-inertia is $\mathbb Z_3(1)$; its degree-one
cohomology is the twist $\mathcal K_n(-1)$ and its higher cohomology
vanishes. This computes the purity shift and the residue-field
action, rather than merely naming a local support group.

Take the direct limit over finite $J$. A continuous cochain with
finite coefficients factors through a finite quotient and is
unramified outside finitely many primes. The generic term is
therefore $R\Gamma(G_{\mathbb Q},\mathcal K_n)$. At each prime the
residue complex is
\[
 [A_q/3^n\xrightarrow{D_q}A_q/3^n]
\]
in degrees zero and one. The transitions of both direct-sum terms
are surjective. Their derived inverse limit is consequently the
literal complex $[\mathcal R(A)\xrightarrow D\mathcal R(A)]$.
Jannsen (1.6) and (2.2) \cite{Jannsen1988} identify the other derived limits with
continuous cochains; the finite groups $H^1(X,\mathcal K_n)$ have
zero first derived limit. Thus the contribution of $X$ in degrees
two and three vanishes.

For injectivity of $D_q$, write $P_D=\mathbb Z_3[D]/\mathbb Z_3\delta$.
The rational graded pieces of $\mathcal K(-1)$ are
\[
 \mathbb Q_3(1)\otimes\bigwedge^2(P_D\otimes\mathbb Q_3),
 \qquad (T\otimes\mathbb Q_3)\otimes(P_D\otimes\mathbb Q_3).
\]
The first has Frobenius eigenvalues of complex absolute value $q$;
the second has absolute value $\sqrt q$. The cusp permutation has
finite order and the elliptic Frobenius polynomial is
$X^2-a_qX+q$. Neither list contains one. An integral endomorphism
invertible over $\mathbb Q_3$ is injective with finite cokernel.
The degree-two and degree-one parts of the limiting triangle now
give (A.1.3). The torsion statement follows from the exact sequence
$0\to\mathcal K\to\mathcal K\otimes\mathbb Q_3\to
(\mathcal K\otimes\mathbb Q_3)/\mathcal K\to0$.
\end{proof}

\begin{statement}{A.2}{Proved}{The exact first derived limit and Lead 5}
Jannsen's degree-two sequence for $\mathcal K$ is
\[
 0\longrightarrow
 \frac{\mathcal R(A)\cap\prod_{q\ne3}D_qA_q}{D\mathcal R(A)}
 \longrightarrow\frac{\mathcal R(A)}{D\mathcal R(A)}
 \longrightarrow
 \varprojlim_n\bigoplus_{q\ne3}\operatorname{coker}(D_q\bmod3^n)
 \longrightarrow0.
 \tag{A.2.1}
\]
Its left term is $\varprojlim{}^1_n H^1(G_{\mathbb Q},\mathcal K_n)$
and is nonzero. It is torsion-free as a $\mathbb Z_3$-module.
In particular Lead 5 is true.
\par\smallskip\noindent\emph{Manuscript comparison: \S51.16,
folios 733v--736v, scans 172--178, formulas (371)--(379).}
\end{statement}
\begin{proof}
The finite-level localization sequence from A.1 gives an exact
sequence
\[
0\to H^1(G_S,\mathcal K_n)\to H^1(G_{\mathbb Q},\mathcal K_n)
 \to\bigoplus_q\ker(D_q\bmod3^n)\to0.
\]
The first groups are finite. Taking first derived limits identifies
the desired group with the first derived limit of the last tower.
Apply inverse limits to the kernels, images, and cokernels of the
surjective-term complexes used in A.1. The image towers have
surjective transitions. Their limits are precisely
$\mathcal R(A)\cap\prod_qD_qA_q$, which proves (A.2.1).

Choose for each $j$ a new prime $q_j$ splitting in the finite Galois
extension cut out by $\mathcal K(-1)/3^j$. Chebotarev \cite[V.3.23]{MilneCFT} supplies
infinitely many choices. Then $D_{q_j}A_{q_j}\subset3^jA_{q_j}$.
Choose a primitive $v_j\in A_{q_j}$ and let $z_{q_j}=D_{q_j}v_j$,
with all other entries zero. The family $z$ is restricted and
pointwise in the image of $D$. Its only preimage is $v_j$ at the
chosen primes, since each $D_q$ is injective; this preimage is not
restricted. Hence $[z]\ne0$ in the left term. For torsion-freeness,
write a left-term representative as $z=Dv$ with $v_q\in A_q$.
A nonzero scalar $m\in\mathbb Z_3$ annihilates its class only when
$mv$ is restricted. Multiplication by a fixed nonzero scalar changes
all valuations by a fixed integer, so then $v$ is restricted.
\end{proof}

\subsection{An actual intrinsic obstruction invisible at all finite levels}

Let $c_d$ denote the central inertia generators in
$A=I/[N,I]=\mathbb Z_3(1)[D]$, and let $e_d$ be their images in $W$.
The rational cusps $O$ and $P$ define a continuous equivariant
homomorphism
\[
 \ell:\mathbb Z_3(1)\longrightarrow N^{\rm cc},\qquad
 z\longmapsto c_O^zc_P^z.
 \tag{A.3.1}
\]
Its image is central. Pull back the central quotient
$N/\gamma_3N\to N^{\rm cc}$ along $\ell$ and call the resulting
geometric group $E_w$, where $w=e_O+e_P$.

\begin{statement}{A.3}{Proved}{A linear family inside the nonabelian lifting problem}
The pullback is an exact sequence of continuous Galois lattices
\[
 0\longrightarrow\mathcal K\longrightarrow E_w
 \longrightarrow\mathbb Z_3(1)\longrightarrow0.
 \tag{A.3.2}
\]
Choose an integral lift $l$ of the standard generator, and after
twisting put
\[
 u_q=F_q(l\otimes1)-l\otimes1\in A_q.
 \tag{A.3.3}
\]
For $a=(a_p)_p\in\mathcal R_{\mathbb P}(\mathbb Z_3)$, the section
$\ell\circ\kappa(a)$ of $N^{\rm cc}\rtimes G_{\mathbb Q}$ has
obstruction represented, under (A.1.3), by
\[
 (a_qu_q)_{q\ne3}.
 \tag{A.3.4}
\]
Its unique possible local corrections are $-a_qD_q^{-1}u_q$.
Thus it lifts precisely when these corrections are integral and
form a restricted family. Changing $l$ changes (A.3.4) by $D$ of
a restricted family.
For the families used below, $a_3=0$. In E24's pro-three coordinates
and the E36 generator $k=-\delta$, their actual torsor coordinates
are
\[
 \alpha:\quad(r_1,r_2;A,B,C)=(0,-\kappa(a);
                  -\kappa(a),-\kappa(a),-\kappa(a)),
 \qquad a_{\rm E36}=-a.
 \tag{A.3.5}
\]
The three $K_0$ coordinates use the inclusion of rational Kummer
classes; the other torsion, framed-$\lambda$, and divisible
coordinates are zero. This specifies the pairs themselves in the
requested convention.
\par\smallskip\noindent\emph{Manuscript comparison: \S51.16,
folios 733v--736v, scans 172--178, formulas (371)--(379).}
\end{statement}
\begin{proof}
The abelian image of a lift of $z$ is $zw$. The commutator of two
lifts is therefore $(zw)\wedge(z'w)=0$ in $\bigwedge^2V$.
The kernel $\mathcal K$ is central, so $E_w$ is abelian. It is
compact, torsion-free, and finitely generated over $\mathbb Z_3$;
the quotient is free of rank one. This proves (A.3.2) and supplies
an integral lift. All actions factor through $G_S$ by smooth
reduction of the pointed curve outside three.

The obstruction is now the ordinary connecting homomorphism
$H^1(G_{\mathbb Q},\mathbb Z_3(1))\to
H^2(G_{\mathbb Q},\mathcal K)$ for (A.3.2). On local tame inertia,
$\kappa(a)(\tau_q)=a_q$. In $E_w(-1)$ the failure of its chosen
lift to be Frobenius invariant is $a_qu_q$. Equivalently, apply the
connecting morphism to the two-term residue complexes of A.1;
the degree-zero element $a_q$ in the trivial quotient has boundary
$a_qu_q$ in the degree-one kernel. This proves (A.3.4) at the
cochain level and hence before passage from the restricted quotient
to a product of local groups. The bounded integral vectors $u_q$
and restricted scalars $a_q$ give a restricted representative.
The final assertion follows from $u_q\mapsto u_q+D_qv$ for a
fixed integral change of lift, with $(a_qv)_q$ restricted.
For (A.3.5), in the non-$O$ cusp basis the vector $w$ has
$P$-coordinate zero and all other coordinates minus one. Formula
(E24.4) sends it to $(0,-1)$; (E24.5) gives the three equal negative
rational Kummer classes. Its lift with $O$-coordinate zero differs
from $c_O+c_P$ by $\delta$. Since $k=-\delta$, the E36 parameter
is $-a$, as asserted.
\end{proof}

\begin{statement}{A.4}{Proved}{A nonzero intrinsic obstruction in the first derived limit}
There is a continuous section of $N^{\rm cc}\rtimes G_{\mathbb Q}$
whose obstruction to $N/\gamma_3N$ is nonzero in
$\varprojlim{}^1_nH^1(G_{\mathbb Q},\mathcal K_n)$.
Its restriction to every local field is zero, and its image in every
finite-coefficient obstruction group is zero. The construction
below specifies such a section as $\ell\circ\kappa(a)$, hence
specifies an actual pair of the torsor coordinates of R5 after the
chosen origins are fixed. Its obstruction has infinite additive
order and survives restriction to every number field.
\par\smallskip\noindent\emph{Manuscript comparison: \S51.16,
folios 733v--736v, scans 172--178, formulas (371)--(379).}
\end{statement}
\begin{proof}
We first prove that the extension (A.3.2) has a nonzero unipotent
action in the kernel of the action on $\mathcal K(-1)$. Let $s_P$
be the rational tangential section at $P$, with parameter $x/z$
and tangent one. Write $s_P=b_Ps_O$ and also $b_P$ for its abelian
cocycle in $V$. For an inertia generator transported along the
chosen path,
\[
 g(c_P)=b_P(g)^{-1}c_P^{\chi(g)}b_P(g),\qquad
 g(c_O)=c_O^{\chi(g)}.
 \tag{A.4.1}
\]
The abelian tangential path is the Kummer class of the framed line
bundle $z_P=\mathcal O(P-O)$ whose frames are $(x/z)^{-1}$ at $P$,
$z/x$ at $O$, and the canonical section at the other cusps.
Indeed the Abel--Jacobi divisor $u-u_0$, as both points approach
the indicated unit tangents, has exactly these two regularized
frames. The Kummer comparison follows by pulling back multiplication
by $3^n$ on the generalized Jacobian, which realizes the
geometrically abelian covering with group $V/3^n$. This identifies
the based path cocycle, not just its projection to the proper curve.

For completeness the needed framed cube is calculated directly.
The function $f=(y+z)/(x+y)$ has divisor $3P-3O$.
Its regularized values are $1/3$ at $P$, $3$ at $O$, and $-1$ at
$Q_c$. Evaluating at the three quadratic cusp pairs and dividing
by the value at $O$ gives
\[
 3z_P=(3^{-2},-3^{-1},\lambda^{-1},\lambda^{-2},
                    \zeta^2\lambda^{-3}).
 \tag{A.4.2}
\]
Here the entries correspond to $P,Q_c$ and the pairs denoted
$A,\bar A$, $B,\bar B$, $C,\bar C$ in LM7 E.24.
The identity $\lambda^2=-3\zeta$ verifies the values and their
conjugates. These are the origin terms needed in this proof.

The action on $V\otimes\mathbb Q_3$ over $K_0$ is abelian.
The toric action is cyclotomic; on the elliptic factor the
endomorphism induced by a coordinate multiplication by $\zeta$
makes $T_3E$ a rank-one module over the discrete valuation ring
$\mathbb Z_3[\zeta]$. The Galois action commutes with this
endomorphism and thus lies in its commutative unit group.
The splitting (A.1.4) proves the assertion for $V$; since the
integral lattice embeds in its rationalization, commutators act
trivially on the integral lattice too.

Choose $\tau\in\ker\chi_3$ with $\kappa_3(\tau)$ a unit. Such an
element exists: $\mathbb Q(\sqrt[3]{3})$ is not contained in
$\mathbb Q(\mu_{3^\infty})$, since every finite subfield of this
abelian extension is Galois over $\mathbb Q$, whereas the real
irreducible pure cubic is not normal. Choose $\sigma\in G_{K_0}$
with $\chi_3(\sigma)=4$, and put $g=[\sigma,\tau]$. The Kummer
homomorphism into $\mathbb Z_3\rtimes\mathbb Z_3^\times$ gives
\[
 \chi_3(g)=1,\qquad\kappa_3(g)=3\kappa_3(\tau).
 \tag{A.4.3}
\]
Thus $g$ fixes $V$ and $\mathcal K(-1)$. On this subgroup,
$\kappa_\zeta(g)=0$ and
$\kappa_\lambda(g)=\kappa_3(g)/2$; the same holds for the conjugate
of $\lambda$. Evaluating (A.4.2) on $g$, where coboundaries in $V$
vanish, gives
\[
 3b_P(g)=\kappa_3(g)v_0,
 \qquad
 v_0=-2e_P-e_{Q_c}
 -\tfrac12(e_A+e_{\bar A})-(e_B+e_{\bar B})
 -\tfrac32(e_C+e_{\bar C}).
 \tag{A.4.4}
\]
The notation is in $W\otimes\mathbb Q_3$; all denominators in
$v_0$ are units in $\mathbb Z_3$.
Take $l=\log c_O+\log c_P$ in the class-two Lie lattice, whose
Baker--Campbell--Hausdorff denominator two is a three-adic unit.
Equation (A.4.1) gives
\[
 u_g=-b_P(g)\wedge e_P
     =-\kappa_3(\tau)v_0\wedge e_P.
 \tag{A.4.5}
\]
The coefficient of $e_P\wedge e_{Q_c}$ in $v_0\wedge e_P$ is one.
The eight non-$O$ cusp classes form a basis of $W$, and $W$ is a
primitive sublattice of $V$. Therefore (A.4.5) is primitive in
$\mathcal K$. A change of based path contributes $(g-1)v=0$;
a change of tangent at $P$ adds a multiple of $e_P$ to $b_P$ and
has zero wedge with $e_P$. This verifies the exact ambiguity.

Apply Chebotarev \cite[V.3.23]{MilneCFT} to the finite representation $E_w(-1)/3^j$.
Choose recursively a new prime $q_j$ whose Frobenius belongs to
the conjugacy class of $g$ in this representation. Then
\[
 D_{q_j}\equiv0\pmod{3^j},\qquad u_{q_j}\notin3A_{q_j}.
 \tag{A.4.6}
\]
Conjugacy preserves these properties: the restriction to the
kernel is the identity modulo $3^j$, and conjugation carries the
off-diagonal vector through an invertible integral matrix.
Define the finite integer
\[
 e_j=\min\{e\ge0:3^eD_{q_j}^{-1}u_{q_j}\in A_{q_j}\}.
 \tag{A.4.7}
\]
It exists by A.1, and (A.4.6) implies $e_j\ge j$. Set
$a_{q_j}=3^{e_j}$ and set every other $a_p$ to zero. This is a
restricted Kummer family. Its residues are pointwise in $D_qA_q$;
their unique corrections
$-3^{e_j}D_{q_j}^{-1}u_{q_j}$ are primitive by the minimality in
(A.4.7). They do not tend to zero. A.3 and A.2 prove that its
obstruction is a nonzero first-derived-limit class.

At each coefficient level only finitely many residue entries
survive. Each is a boundary, so the finite-coefficient obstruction
is zero. At every local field, the finite degree-one cohomology
groups satisfy Mittag--Leffler; continuous degree-two cohomology
therefore injects into the limit of finite degree-two groups.
This proves local vanishing also at three. Real three-primary
cohomology vanishes by averaging. Finally A.2 proves infinite
additive order. Restriction followed by corestriction over a finite
extension $L/\mathbb Q$ is multiplication by $[L:\mathbb Q]$ on
continuous cohomology, by the finite coset cochain formula.
Hence no restriction can kill this obstruction.
\end{proof}

\subsection{Two ramification assertions disproved}

\begin{statement}{A.5}{Refuted}{Lead 2: mixed tame inertia already occurs at the prime two}
Lead 2 is false. The section
\[
 s=\ell\circ(3^{23}\kappa_2)
       :G_{\mathbb Q}\longrightarrow N^{\rm cc}\rtimes G_{\mathbb Q}
 \tag{A.5.1}
\]
lifts to $N/\gamma_3N$. Its abelian tame-inertia vector at two is
$3^{23}(e_O+e_P)$, which is not a multiple of any cusp vector.
More generally, for each rational prime $q\ne3$, every multiple
of $3^{v_3\det D_q}\kappa_q$ in the line (A.3.1) lifts.
\par\smallskip\noindent\emph{Manuscript comparison: \S51.16,
folios 733v--736v, scans 172--178, formulas (371)--(379).}
\end{statement}
\begin{proof}
For a fixed $q$ the only nonzero residue of this family is a scalar
multiple of $u_q$. The adjugate identity shows that
$3^{v_3\det D_q}D_q^{-1}$ is integral. Its local correction has
finite support and is therefore restricted. A.3 proves the general
assertion.

At two the cusp permutation has signs $(5,3)$ on $P_D$. Thus its
exterior square has signs $(13,15)$, and the determinant of
$F_2-1$ on $\bigwedge^2W(-1)$ has three-adic valuation $15$.
The cubic has the three projective points
$[0:1:1],[1:0:1],[1:1:0]$ over $\mathbb F_2$, so its elliptic
Frobenius polynomial is $X^2+2$. On each of the eight copies of
$T$, twisted by a cusp permutation sign, the determinant of
$F_2-1$ is three. The remaining contribution is eight. Consequently
\[
 v_3\det D_2=15+8=23.
\]
This determinant calculation uses the rational filtration and is
valid on the actual integral lattice because determinants are
unchanged by rational change of basis.

To check the cusp claim integrally, eliminate $e_O$ by the relation
$\sum_de_d=0$. In the basis of the other eight cusps,
$e_O+e_P=(0,-1,-1,-1,-1,-1,-1,-1)$. It is proportional neither
to a basis vector nor to $e_O=(-1,-1,-1,-1,-1,-1,-1,-1)$.
Multiplication by a nonzero three-adic scalar preserves this fact.
The exact finite calculations are independently executed in
\texttt{code/a\_residue\_checks.py}.
\end{proof}

\begin{statement}{A.6}{Refuted}{Lead 3: infinitely ramified noncuspidal class-two sections}
Lead 3 is false. Enumerate the rational primes different from three
as $q_1,q_2,\ldots$, put $m_j=v_3\det D_{q_j}$, and set
\[
 a_{q_j}=3^{m_j+j},\qquad a_3=0.
 \tag{A.6.1}
\]
The section $\ell\circ\kappa(a)$ lifts to $N/\gamma_3N$ and every
such lift is ramified at all the $q_j$. None is cuspidal.
The same conclusion holds after replacing these coefficients by
any nonzero coefficients satisfying
$v_3(a_{q_j})-m_j\longrightarrow+\infty$.
\par\smallskip\noindent\emph{Manuscript comparison: \S51.16,
folios 733v--736v, scans 172--178, formulas (371)--(379).}
\end{statement}
\begin{proof}
The family is restricted. By the adjugate identity,
$a_{q_j}D_{q_j}^{-1}u_{q_j}\in3^jA_{q_j}$, so the unique corrections
are integral and restricted. A.3 proves existence of a continuous
lift. Its image in $V$ has nonzero tame-inertia vector
$a_{q_j}(e_O+e_P)$ at every $q_j$; thus geometric conjugation cannot
remove the ramification. A cuspidal section at a rational cusp
has all abelian tame-inertia vectors in the line generated by
that cusp. A.5 proves that no one of the displayed vectors lies
in any such line. The nonrational cusps have decomposition groups
whose image in $G_{\mathbb Q}$ has index two, so cannot support
a section over $\mathbb Q$. This proves noncuspidality and the
more general assertion with the stated convergence bound.
\end{proof}

\subsection{The lift torsor and a formula retaining all path origins}

\begin{statement}{A.7}{Proved}{The integral group acting on every nonempty fibre}
There is an isomorphism of finitely generated $\mathbb Z_3$-modules
\[
 H^1_{\rm cts}(G_{\mathbb Q},\mathcal K)
   \simeq\mathbb Z_3^{23}\oplus(\mathbb Z/3)^9
                         \oplus(\mathbb Z/9)^5.
 \tag{A.7.1}
\]
Every nonempty fibre of
$\operatorname{Sec}(N/\gamma_3N)\to\operatorname{Sec}(N^{\rm cc})$
is a torsor under this group. The ambiguity in (A.7.1) is the choice
of a Smith basis and of a free direct summand. No further geometric
conjugations identify members of the torsor.
\par\smallskip\noindent\emph{Manuscript comparison: \S51.16,
folios 733v--736v, scans 172--178, formulas (371)--(379).}
\end{statement}
\begin{proof}
We give the integral lattice and the finite calculation determining
its torsion. Write $\mathcal O_3=\mathbb Z_3[\zeta]$ and identify
$T$ with $\mathcal O_3$ so that $P$ is the class of $1$ in $T/3T$.
The automorphism $[x:y:z]\mapsto[x:y:\zeta z]$ fixes $O$ and
acts on its tangent by $\zeta$. It is the chosen complex
multiplication. Its orbit of $P$ is $P,A_{\!\rm bar},A$.
The eight nonzero cusp labels in the basis $(P,\zeta P)$ are
\[
\begin{array}{c|rrrrrrrr}
 d&P&Q_c&A&\bar A&B&\bar B&C&\bar C\\ \hline
 a_d&1&-1&-1&0&1&0&-1&1\\
 b_d&0&0&-1&1&1&-1&1&-1
\end{array}
\tag{A.7.2}
\]
For example the line through $P$ and $-\zeta P$ has third
intersection $[1:-\zeta:0]=C$, so
$C=-(P-\zeta P)$. This determines the only labels not obtained
immediately by multiplication by $\zeta$ and negation.

Let $\Psi$ be the $8\times2$ integer matrix with row $(b_d,-a_d)$.
In the rational representation $W_{\mathbb Q_3}\oplus T_{\mathbb Q_3}$
the lattice $V$ is
\[
 L_\Psi=\{(w+\Psi t/3,t):w\in W,\ t\in T\}.
 \tag{A.7.3}
\]
Here is the comparison proving this formula. Pullback of the
semiabelian extension $J_{E,D}$ along $[3]:E\to E$ is split, since
its class is killed by three. The splitting is unique because
$\operatorname{Hom}(E,T_D)=0$. Thus it is defined over $\mathbb Q$.
The resulting isogeny $T_D\times E\to J_{E,D}$ has kernel the graph
of a homomorphism $E[3]\to T_D[3]$; its projection to $E$ is $[3]$.
Evaluation of the graph on a degree-zero cusp character gives the
Weil pairing with that cusp difference, by the defining Poincare
biextension of the generalized Jacobian. In the bases in (A.7.2),
this pairing is $t\mapsto\Psi t$ modulo three. One overall sign
changes $\Psi$ to $-\Psi$ and conjugates all the matrices by
$(-1)_W\oplus1_T$, so it does not affect their invariant factors.
The Tate lattice of the quotient is obtained by adding
$(\Psi t/3,t/3)$ to $W\oplus T$ and then using $[3]$ on the second
coordinate. This is precisely (A.7.3), including its integral
comparison with the geometric Tate module.

For a Galois element with matrices $R$ on $W$ and $T_g$ on $T$,
the matrix in the basis in (A.7.3) is therefore
\[
 G_V=\begin{pmatrix}
 R&(R\Psi-\Psi T_g)/3\\0&T_g
 \end{pmatrix}.
 \tag{A.7.4}
\]
The numerator is divisible by three by equivariance of the Weil
pairing. Over $K_0$, $R=\chi(g)I_8$. Complex conjugation has
$R=-P_c$, where $P_c$ fixes $P,Q_c$ and interchanges the three
quadratic cusp pairs, and
\[
 T_c=\begin{pmatrix}1&-1\\0&-1\end{pmatrix}.
\]

The exact Tate image over $K_0$ is $1+3\mathcal O_3$.
Containment follows from rationality of $E[3]$ and complex
multiplication. To prove equality, count the Fermat cubic over
$\mathbb F_7$: it has nine points, so its Frobenius polynomial is
$X^2+X+7$. Its roots in $\mathcal O_3$ are
$\alpha=1+3\zeta$ and $\bar\alpha=-2-3\zeta$. Frobenius at a prime
above seven gives one root; conjugation gives the other. The
three-adic logarithm identifies $1+3\mathcal O_3$ with
$3\mathcal O_3$. The vectors
$\log(\alpha)/3$ and $\log(\bar\alpha)/3$ reduce to
$\zeta$ and $-1-\zeta$ modulo three, a basis of
$\mathcal O_3/3\mathcal O_3$. Their closed $\mathbb Z_3$-span is
therefore all $\mathcal O_3$. This proves equality without an
unspecified index from a complex multiplication reciprocity theorem.
The matrices on $T$ are
\[
 T_\alpha=\begin{pmatrix}1&-3\\3&-2\end{pmatrix},\qquad
 T_{\bar\alpha}=\begin{pmatrix}-2&3\\-3&1\end{pmatrix}.
 \tag{A.7.5}
\]
The determinants are seven. Equations (A.7.4)--(A.7.5) and $c$
therefore topologically generate the actual image on $V$.

Take exterior squares in (A.7.4) and omit the basis wedge of the
two $T$ lifts. The remaining forty-four wedges are a basis of
$\mathcal K$, since that omitted wedge maps to a generator of
$\bigwedge^2T$. Let the three resulting matrices be
$G_\alpha,G_{\bar\alpha},G_c$. The integer matrix
\[
 B=\begin{pmatrix}G_\alpha-1\\G_{\bar\alpha}-1\\G_c-1\end{pmatrix}
       :\mathbb Z_3^{44}\longrightarrow\mathbb Z_3^{132}
 \tag{A.7.6}
\]
has three-adic Smith diagonal
\[
 \underbrace{1,\ldots,1}_{30},\quad
 \underbrace{3,\ldots,3}_{9},\quad
 \underbrace{9,\ldots,9}_{5}.
 \tag{A.7.7}
\]
The program \texttt{code/a\_k\_lattice.py} constructs every entry
from (A.7.2)--(A.7.6), checks $c^2=1$ and
$cG_\alpha c=G_{\bar\alpha}$, and performs elementary row and
column operations over $\mathbb Z/3^8$. Each operation is
invertible and every pivot has valuation at most two. Thus no
factor divisible by $3^8$ remains, and this calculation determines
the exact three-adic Smith exponents. The full matrices and the
output are included. As an independent consistency check, the
same calculation on $V$ gives seven unit factors and three factors
three, recovering $J_{E,D}(\mathbb Q)[3^\infty]=(\mathbb Z/3)^3$.

An invariant in $(\mathcal K\otimes\mathbb Q_3)/\mathcal K$ is
exactly an $x$ with $Bx$ integral, modulo integral $x$. A Smith
factor $3^e$ contributes $\mathbb Z/3^e$. A.1 and (A.7.7) prove
(A.7.1). Finally fix a cocycle representing a section of
$N^{\rm cc}$. Central lifts differ by $\mathcal K$-valued cocycles
and central conjugation gives coboundaries. A geometric element
of $N^{\rm cc}$ stabilizing that section has invariant image in
$V$, which is zero. It is then in the central kernel
$\mathbb Z_3(1)$ and is invariant there, hence zero again.
There is no additional stabilizer quotient. This proves the exact
fibre statement.
\end{proof}

\begin{statement}{A.8}{Proved}{The full class-two residue criterion in based cusp-path coordinates}
Choose geometric paths from the specified tangent at $O$ to a
Galois-equivariant family of tangents at the nine cusps. Let $c_d$
be the corresponding inertia elements in $N/\gamma_3N$ and
$b_{g,d}\in V$ their abelian comparison cochains, with convention
\[
 g(\log c_d)=\chi(g)\bigl(\log c_{gd}
                         -b_{g,d}\wedge e_{gd}\bigr).
 \tag{A.8.1}
\]
For a continuous representative $s$ of a class in
$\operatorname{Sec}(N^{\rm cc})$, write
\[
 \log s(\tau_q)=\sum_{d\in D}t_{q,d}c_d,\quad
 v_q=\alpha(\phi_q),\quad
 w_q=\sum_dt_{q,d}e_d\in W.
 \tag{A.8.2}
\]
Here $s(\tau_q)$ means its geometric cocycle value, the chosen
arithmetic action is unramified at $q$, and $\phi_q\tau_q\phi_q^{-1}
=\tau_q^q$. Then its full intrinsic obstruction is represented by
the restricted family
\[
 r_q=v_q\wedge w_q-
            \sum_dt_{q,d}\,b_{\phi_q,d}\wedge e_{\phi_qd}
       \in\mathcal K(-1).
 \tag{A.8.3}
\]
The section lifts exactly when
\[
 -D_q^{-1}r_q\in A_q\quad\text{for all }q\ne3,
 \qquad (-D_q^{-1}r_q)_q\in\mathcal R(A).
 \tag{A.8.4}
\]
The nonempty fibres are exactly the torsors in A.7.
\par\smallskip\noindent\emph{Manuscript comparison: \S51.16,
folios 733v--736v, scans 172--178, formulas (371)--(379).}
\end{statement}
\begin{proof}
The comparison cochains exist by transport of the inertia loops
along the chosen paths. Composition gives
$b_{gh,d}=b_{g,hd}+g b_{h,d}$ in $V$. Formula (A.8.1) is the
class-two expansion of conjugation and includes its sign convention.
It specifies the action on the cusp lifts themselves, including
the path origins.

The cocycle restriction to local inertia is a homomorphism,
since the original action is unramified there. Wild inertia maps
to zero and the prime-to-three tame factors also map to zero.
The elliptic projection $t_T$ satisfies $F_qt_T=qt_T$.
The Frobenius eigenvalues on $T$ have absolute value $\sqrt q$,
so $t_T=0$. Hence the logarithm of tame inertia lies in the central
lattice $A=\mathbb Z_3(1)[D]$ in $N^{\rm cc}$, proving (A.8.2).
The relation also gives $t_{q,\phi_qd}=t_{q,d}$.

Lift this inertia logarithm to
$L_q=\sum_dt_{q,d}\log c_d+k_q$ in the full class-two lattice,
where $k_q\in\mathcal K$. Lift the Frobenius cocycle arbitrarily;
its abelian coordinate is $v_q$. In a class-two group conjugation
by that lift changes a logarithm of abelian coordinate $qw_q$
by $qv_q\wedge w_q$. The logarithm of the tame-relation defect,
divided by $q$ to identify inertia homomorphisms with the twist
$\mathcal K(-1)$, is therefore
\[
 D_qk_q+v_q\wedge w_q-
       \sum_dt_{q,d}b_{\phi_q,d}\wedge e_{\phi_qd}.
\]
The central coordinate of the Frobenius lift cancels. This is
exactly the cochain-level localization boundary: it measures the
failure of a chosen inertia splitting to be invariant under
Frobenius. At every finite coefficient level the action and the
section are unramified at almost all primes, so these choices give
a direct-sum residue cochain. The inverse-limit family is
restricted; concretely every term contains one of the restricted
coefficients $t_{q,d}$ and all other vectors are integral. A.1 now
identifies the global obstruction with $[r_q]$, before any passage
to a product of local groups. Injectivity of every $D_q$ proves
(A.8.4), with the stated convergence condition. A.7 supplies the
geometric-conjugacy fibre.
\end{proof}

\paragraph{Coordinates of the tame coefficients.}
The formulas can be connected to the chosen E24 coordinates without
suppressing the central parameter. Separate those coordinates into
the entries supported at three, including $\mu_3^3$, the framed
point $z'$, and the noncanonical divisible $\mathbb Z_3$ parameter,
and the prime valuations outside three. The elliptic projection is
always unramified outside three: the same weight calculation used
in A.8 makes every local integral elliptic tame vector zero.
The surjectivity $H^1(G_S,V)\to H^1(G_S,T)$ follows from
$H^2(G_S,W)=0$, proved by the same composition-series argument as
A.1. Thus origins can be chosen by a full unramified lift for the
first set of parameters, as constructed in LM7 E.34, and by the
central splitting $W\to A$ with $O$-coordinate zero for the
remaining toric coordinates. We make this choice agree with the E36 convention as follows.
For each elliptic class $\beta$, choose an unramified section $c_\beta$
of $C$; take $c_0$ to be the unit tangent at $O$, and take the
classes of $P-O$ and $Q_c-O$ to the images of their unit tangents.
The reference over every $\alpha$ is $(c_\beta,\alpha)$ in
$C\times_TV$. The prescribed unramified reference itself has a
full unramified lift by the free pro-three embedding argument.
The E36 central generator is $k=[a,b]=-\delta$, so the exact
coordinate comparison is $a_\delta=-a_{\rm E36}$. In the formulas
below $a$ always means $a_{\rm E36}$. This fixes both the origin
and its sign; it agrees with the packets in B.3.

Write $r_{1,q},r_{2,q}$ for the valuations of the two rational
toric coordinates, and $\nu_{A,\mathfrak q}$,
$\nu_{B,\mathfrak q}$, $\nu_{C,\mathfrak q}$ for those of the
three $K_0$ toric coordinates. Put
$N\nu_{A,q}=\sum_{\mathfrak q\mid q}f_{\mathfrak q/q}
\nu_{A,\mathfrak q}$, and likewise for $B,C$. In this chosen
origin, for $q\ne3$,
\[
\begin{aligned}
 t_{q,O}&=-a_q,\\
 t_{q,P}&=r_{1,q}+N\nu_{C,q}-N\nu_{A,q}-a_q,\\
 t_{q,Q_c}&=r_{2,q}+N\nu_{C,q}-N\nu_{B,q}-a_q,\\
 t_{q,d}&=\nu_{A,\mathfrak q(d)}-a_q
                  &&(d\in\{A,\bar A\}),
\end{aligned}
\tag{A.8.5}
\]
with identical last formulas for the $B$ and $C$ pairs. At an inert
prime the two valuations in a pair coincide and their norm
contribution is twice that valuation. At a split prime each has
residue degree one. These formulas are the valuation form of the
actual embedding (E24.5), followed by minus the central diagonal
$\mathbb Z_3(1)\delta$; the diagonal is the kernel of $A\to W$.

The cochains $b_{g,d}$ in (A.8.3) are abelian data. Let $z_d$ be
the framed tangential divisor $d-O$ in $J_{E,D}$ and choose
compatible division points $z_{d,n}$ with $3^nz_{d,n}=z_d$.
Then $g z_{d,n}-z_{gd,n}\in J_{E,D}[3^n]$ gives their reduction
modulo $3^n$, with the chosen path convention. Their origin
relations are obtained by the tangent-line functions $f_d$ with
$\operatorname{div}f_d=3(d-O)$, regularized at the indicated
tangents. Formula (A.4.2) is the complete $P$ instance and is the
one needed for the proved counterexamples. Statements A.9 and A.10 below supply the complete regularized
origin table and fix the extra unramified generator. Evaluating a
Frobenius in that free pro-three Galois group is part of the
arithmetic input; a finite-query program on arbitrary completed
families cannot decide the final convergence condition.


\begin{statement}{A.9}{Proved}{All nine tangential origin data, including roots of unity}
Order the non-$O$ cusps as $P,Q_c,A,\bar A,B,\bar B,C,\bar C$.
Use $x/z$ at the first and third/fourth cusps, $y/z$ at the second
and fifth/sixth cusps, and $z/x$ at the last two; all tangents are
one. For $z_d=\mathcal O(d-O)$ with these regularized frames,
write $3z_d=h_d\in T_D(K_0)$. Its $O$-frame is one. The other
entries are $\pm\zeta^{Z_{d,e}}\lambda^{L_{d,e}}$, with matrices
\[
 L=\begin{pmatrix}
-4 & -2 & -1 & -1 & -2 & -2 & -3 & -3 \\
-2 & -4 & -2 & -2 & -1 & -1 & -3 & -3 \\
-1 & -2 & -4 & -1 & -2 & -2 & -3 & -3 \\
-1 & -2 & -1 & -4 & -2 & -2 & -3 & -3 \\
-2 & -1 & -2 & -2 & -4 & -1 & -3 & -3 \\
-2 & -1 & -2 & -2 & -1 & -4 & -3 & -3 \\
-3 & -3 & -3 & -3 & -3 & -3 & -6 & -3 \\
-3 & -3 & -3 & -3 & -3 & -3 & -3 & -6
 \end{pmatrix},\qquad
 Z=\begin{pmatrix}
2 & 1 & 0 & 1 & 0 & 2 & 2 & 1 \\
1 & 2 & 0 & 2 & 0 & 1 & 1 & 2 \\
1 & 2 & 2 & 0 & 1 & 0 & 2 & 1 \\
0 & 0 & 1 & 2 & 2 & 1 & 2 & 1 \\
2 & 1 & 1 & 0 & 2 & 0 & 1 & 2 \\
0 & 0 & 2 & 1 & 1 & 2 & 1 & 2 \\
1 & 2 & 1 & 1 & 2 & 2 & 0 & 0 \\
2 & 1 & 2 & 2 & 1 & 1 & 0 & 0
 \end{pmatrix}.
 \tag{A.9.1}
\]
The signs, recorded in the exact program output, have zero pro-three
Kummer cocycle, using the rational compatible odd-power roots of
$-1$. No $\zeta$ term may be discarded. The row at $O$ is the
identity. These data are equivariant under complex conjugation.
\par\smallskip\noindent\emph{Manuscript comparison: \S51.16,
folios 733v--736v, scans 172--178, formulas (371)--(379).}
\end{statement}
\begin{proof}
For $d=[x_d:y_d:z_d]$, put
$L_d=x_d^2x+y_d^2y+z_d^2z$ and $f_d=L_d/(x+y)$.
The flex property proves $\operatorname{div}f_d=3d-3O$.
At $O$, write $s=z/x$ and
$y/x=-(1+s^3)^{1/3}$. The regularized value of $f_d$ is
\[
 c_O(d)=\lim_{s\to0}f_ds^3=-3(x_d^2-y_d^2).
\]
For $d=[0:-\eta:1]$ or $[-\eta:0:1]$, with $\eta^3=1$,
the regularized value at $d$ is $1/(3\eta)$, by expanding the
nonzero coordinate as $-\eta(1+s^3)^{1/3}$.
For $d=[1:-\eta:0]\ne O$ it is
$-1/(3(1-\eta))$. At every other cusp $e$ the ordinary value
$f_d(e)$ applies. Division of all these values by $c_O(d)$ gives
$h_{d,e}$. This proves the claimed framed relation over $K_0$.
All arithmetic in the evaluation is in $\mathbb Q(\zeta)$;
\texttt{code/a\_cusp\_origins.py} evaluates the formulas as rational
pairs in the basis $1,\zeta$, factors each nonzero value uniquely
as $\pm\zeta^a\lambda^b$, and verifies conjugation entry by entry.
It also verifies the complete $P$ row against (A.4.2).
This gives (A.9.1), not a numerical approximation to its entries.
\end{proof}

\begin{statement}{A.10}{Proved}{An explicit choice of the extra unramified cohomology generator}
Let $P_3$ be the maximal pro-three quotient of $G_{K_0,S}$ and
let $\Gamma=P_3\rtimes\langle c\rangle$ be the relative quotient
of $G_S$. Choose $u$ to be Frobenius at the prime
$(7,\zeta-2)$, with a decomposition group fixed once and for all,
and put $v=cuc$. Then $u,v$ are a free pro-three basis of $P_3$,
and $c$ interchanges them. In the representation (A.7.3) let
$\widetilde t_1=(\Psi(1,0)/3,(1,0))\in V$. The prescriptions
\[
 \alpha_\infty(c)=0,\qquad
 \alpha_\infty(u)=\widetilde t_1
 \tag{A.10.1}
\]
define a continuous $V$-valued cocycle on $\Gamma$ and fix a
splitting of the final $\mathbb Z_3$ coordinate in (E24.1).
Together with the specified prime generators and the framed points,
this fixes the formerly noncanonical cohomology coordinate.
\par\smallskip\noindent\emph{Manuscript comparison: \S51.16,
folios 733v--736v, scans 172--178, formulas (371)--(379).}
\end{statement}
\begin{proof}
The vanishing $H^2(G_{K_0,S},\mathbb F_3)=0$ was proved in A.1.
Inflation injects $H^2(P_3,\mathbb F_3)$ into this group: an invariant
nonzero character of the kernel of $G_{K_0,S}\to P_3$ would give
a central extension of $P_3$ by $\mathbb F_3$, a pro-three quotient
contradicting maximality. The five-term exact sequence proves the
injection. The minimal free pro-three presentation of $P_3$ has
two generators, by the Kummer calculation in A.1; its relation
module has dual $H^2(P_3,\mathbb F_3)=0$. A nontrivial closed
relation subgroup in the Frattini subgroup has a nonzero first
term in the lower exponent-three filtration, so the relation
module cannot vanish. Thus $P_3$ is free on two generators.

At the chosen prime over seven, finite-field Kummer theory gives
the evaluations $(2,0)$ on the classes $(\zeta,\lambda)$:
$\zeta\mapsto2$, $\lambda\mapsto-1$ in $\mathbb F_7$, and their
squares are $\zeta^2$ and one. At the conjugate Frobenius the
values are $(2,2)$. They are independent over $\mathbb F_3$.
Consequently $u,cuc$ generate the Frattini quotient of $P_3$;
the pro-three basis theorem and freeness make them a free basis.
Complex conjugation is an actual order-two element of $G_S$, so
its relations are exactly $c^2=1$, $cuc=v$, and $cvc=u$.

For each finite three-primary quotient, arbitrary cocycle value
on $u$, its prescribed conjugate value on $v$, and value zero on
$c$ respect all these relations. The free pro-three universal
property gives a cocycle; taking the compatible inverse limit
gives (A.10.1). On $T$ the normalized cocycles are $T$ modulo
$(u-1)T^c$. Now $T^c=\mathbb Z_3(1,0)$ and
$(T_\alpha-1)(1,0)=3(0,1)$, so
\[
 H^1(G_S,T)=\mathbb Z_3(1,0)\oplus(\mathbb Z/3)(0,1).
\]
The projection of (A.10.1) is the free generator. The finite summand
is the Kummer class of $P$, and filling the cusps identifies the
free quotient with the Tate module of the divisible part used in
(E24.9). Thus the given lift defines the asserted splitting.
The Frobenius scalar at $(7,\zeta-2)$ is $1+3\zeta$, rather than
its conjugate: its reduction has zero differential at that prime,
which selects the root divisible by the prime. This verifies the
matrix used in (A.10.1).
\end{proof}

\begin{statement}{A.11}{Proved}{Executed finite and certified-tail evaluations}
The program \texttt{code/a\_criterion.py} implements (A.8.4) for
finite-support exact residue matrices and for the two infinite
families A.4 and A.6 with their proved tail certificates.
It gives the following verified results:
\begin{enumerate}
\item the E.28 quotient of $(\kappa_6,\kappa_{109})$ has correction
$-1/12$ at $109$, of three-adic valuation $-1$; all $729$ central
lift pairs modulo $27$ fail the tame relation;
\item the E.31 origin retains $\kappa_3(\zeta)$, its local
Frobenius and inertia values are $2$ and $1$, and all nine
central lift pairs fail; the entire coset $\alpha+3\mathcal K_0$
has the same excluded image;
\item all the off-three residues vanish for every section
unramified outside three;
\item the E.35 norm polynomial is
$X^3+AY^3+A^2Z^3-3AXYZ$, its example $A=2,b=3$ has witness
$(1,1,0)$, and its finite conjugation calculation has span-rank
histogram $(1,32,48)$ and $153$ orbits over all $81$ abelian
assignments and $729$ generator assignments.
\end{enumerate}
The infinite family A.6 is certified by correction valuations at
least $j$, whereas A.4 is certified by primitive corrections at
every selected prime. These two conclusions use their complete
tail proofs, with no inference from a finite sample.
\par\smallskip\noindent\emph{Manuscript comparison: \S51.16,
folios 733v--736v, scans 172--178, formulas (371)--(379).}
\end{statement}
\begin{proof}
For exact finite-support data, the program calculates $-D_q^{-1}r_q$
by rational linear algebra and tests the three-adic valuation of
every coordinate. All unlisted corrections are zero, so this is
exactly (A.8.4). The two family constructors apply the adjugate
bound of A.6 and the minimal-denominator proof of A.4, respectively.
Their output records these bounds explicitly.

The comparison for E.28 is the actual covering homomorphism
$N\to\widehat F_2^{(3)}$ followed by its class-two quotient.
The two coordinate inertia exponents multiply by three, so the
bracket multiplier is nine. At $109$, choose the residue generator
$6^4$ of the three-primary roots of unity; the Kummer value of six
is one. The resulting differential is $109-1=108$ and the residue
is nine, giving $-9/108=-1/12$. The program also enumerates both
central coordinates modulo $27$ and checks the same nonzero defect
for all $27^2$ possibilities.

For E.31 use the actual function
$h=(x+y)/(x+\zeta y)$ on the cubic, with the same tangent at $O$.
Writing $s=z/x$, its leading constant is
$-1/(3\lambda)=\zeta\lambda^{-3}$. At
$\mathfrak q=(3+\zeta)$ over seven, $\zeta=4$, $\lambda=4$ and
$\bar\pi=2-\zeta=5$. At the specified torus point the total
$h$ unit reduces to three, whose square is two, the square of
$\zeta=4$. Thus the Kummer Frobenius exponent is two and the
other coordinate has inertia exponent one. The commutator is
nonzero in the order-$27$ Heisenberg quotient. The morphism on
abelian cohomology to these two $\mathbb F_3$ coordinates is
additive, with the fixed affine origin just computed; adding
$3\mathcal K_0$ leaves it unchanged. This proves the whole coset
assertion, rather than a list of sampled coset elements.

For an unramified section $t_{q,d}=0$ for every $q\ne3$;
every term of (A.8.3) is zero. The program verifies the corresponding
polynomial identity. The existence of lifts to the full $N$ uses
the free pro-three argument, not this finite computation.

Finally the inner action of E.31 is removed by the explicit
arithmetic-group isomorphism
$(n,g)\mapsto(nq_g,g)$, $q_g=(\kappa_3(\zeta)(g),0,0)$.
It sends the abelian pair $(a,b)$ to $(A,b)=(\zeta a,b)$.
The central extension then has cup obstruction
$\kappa_3(A)\cup\kappa_3(b)$. The crossed-product comparison \cite[IV.3.14]{MilneCFT} represents
this class by the cyclic algebra
$L_A\oplus L_AU\oplus L_AU^2$, with
$U\ell=\sigma(\ell)U$ and $U^3=b$. It splits exactly when
$b=N_{L_A/K_0}(w)$: a norm witness gives the representation
$U=w\sigma$ on the three-dimensional $K_0$-space $L_A$;
a splitting conversely gives that form for $U$ on a one-dimensional
$L_A$-module. Its norm multiplication matrix on $1,t,t^2$ is
\[
 \begin{pmatrix}X&AZ&AY\\Y&X&AZ\\Z&Y&X\end{pmatrix}.
\]
The determinant is the displayed polynomial. The finite group law
$\tfrac12(uv'-vu')$ gives the central change $rv-su$ under
conjugation. Thus the complete nonempty fibre is
$K_0^\times/(K_0^{\times3}\langle A,b\rangle)$.
The program checks the determinant identity, the norm witnesses,
and every finite conjugation orbit. This confirms the exact
quotient calculation used in E.35; it does not assert that a lift
in this finite quotient already lifts to $N/\gamma_3N$.
\end{proof}

\begin{statement}{A.12}{Refuted}{A universal finite-precision decision program}
There is no program which, for every completed Kummer input in
A.3 supplied through finite-coefficient queries, terminates and
decides whether its continuous intrinsic obstruction vanishes.
The same failure holds for the full inputs of A.8. Finite-support
inputs and inputs furnished with proved tail certificates have
the exact evaluator in A.11.
There is also a concrete family of inputs described by finite
programs, each returning the actual R4/E36 finite-coefficient
cocycles and each with the uniform known tail bound
$v_3(a_j)\ge j$, on which the vanishing decision is
undecidable. Thus replacing an oracle input by a finite program
for its coefficients does not provide a total evaluator on
this explicitly specified class. The task's phrase ``given
$(\alpha,a)$'' does not itself specify an input encoding;
these two precise input classes state the scope of the
refutation, while A.11 states the implemented domain and
proves all four requested reproductions.
\par\smallskip\noindent\emph{Manuscript comparison: \S51.16,
folios 733v--736v, scans 172--178, formulas (371)--(379).}
\end{statement}
\begin{proof}
Run such a program on the zero input. Its terminating computation
uses finitely many coefficient levels, with maximum $3^n$.
Discard the first $n$ terms of the nonzero obstruction constructed
in A.4. The resulting input is zero modulo $3^n$, since its
remaining coefficients have valuations at least their indices.
Every query in the alleged computation therefore receives exactly
the same answer. Nevertheless its corrections remain primitive at
infinitely many primes, so A.3 proves that its obstruction is
nonzero. The program gives the same answer on two inputs requiring
different answers, a contradiction. The counterexample is a tail
of the actual intrinsic family, rather than an arbitrary formal
sequence outside the section problem. This establishes the stated
precise computational limit and explains the role of the tail
certificates in A.11.

For the finite-program assertion, first make the family A.4
effective. At level $3^r$, the representation $E_w(-1)/3^r$
is determined by the $3^r$-division points of the explicit
generalized Jacobian and the $3^r$-division torsor of the
framed point $P-O$. The latter point and its torus multiple
are the actual framed point of A.9 and (E24.6).
Multiplication by $3^r$ is a finite algebraic morphism on
this semiabelian variety in characteristic zero. Its fiber
equations, the elliptic division polynomials, and the torus
root equations give a finite collection of algebraic numbers
and their exact Galois action. These are effective finite
algebraic operations: roots are represented by polynomials
and isolating data, and their equalities, embeddings, and
permutations are decided by exact algebraic arithmetic.
The exterior-square map and the cochain
$b_{g,P}\wedge e_P$ then give the finite representation
used in A.4, in the fixed existing coordinates.

We can consequently enumerate primes and test the finite
conditions $D_q\equiv0\pmod{3^j}$ and
$u_q\not\equiv0\pmod3$ by reduction in that explicitly
constructed finite Galois extension. Take $q_j$ to be the
first unused larger prime passing these tests. A.4 and
Chebotarev prove termination of each search. Put
$m_j=v_3(\det D_{q_j})$. The matrix $D_{q_j}$ is effective
from the integral Frobenius matrices, and is invertible
over $\mathbb Q_3$ by the weight calculation. Its inverse
has denominators dividing $3^{m_j}$. The minimal exponent
$e_j$ such that $3^{e_j}D_{q_j}^{-1}u_{q_j}$ is integral
is therefore determined by the finite representation
modulo $3^{m_j+1}$; it satisfies $j\le e_j\le m_j$.
This constructs uniformly computable $(q_j,e_j)$ and
all finite cocycles of the actual A.4 input, rather than
a formal sequence of unrelated residue classes.

Now fix an ordinary finite Turing-machine description $T$.
Define its input coefficients by the bounded calculation
\[
 a_{T,q_j}=
 \begin{cases}
  3^{e_j},&T\text{ has not halted within }j\text{ steps},\\
  0,&T\text{ has halted within }j\text{ steps},
 \end{cases}
 \qquad a_{T,q}=0\quad(q\notin\{q_j\}).
 \tag{A.12.1}
\]
Every coefficient query requires only a bounded simulation.
To return the entire input modulo $3^n$, the program
constructs $q_j,e_j$ for $j<n$, performs those finitely
many bounded simulations, and returns the finite Kummer
product with the surviving exponents $3^{e_j}$.
Every omitted coefficient is divisible by $3^n$.
Thus $j\ge n$ is one uniform convergence modulus,
independent of whether $T$ halts.

The input section is the actual equivariant line
$\ell\kappa(a_T)$ of A.3. Its R4 coordinates are
$(0,-\kappa(a_T);-\kappa(a_T),-\kappa(a_T),-\kappa(a_T))$
and its E36 parameter is $-\kappa(a_T)$, with the
same fixed origins as A.8. The finite Kummer product
and its algebraic roots supply the corresponding
finite cocycle; the preceding division-point
construction supplies every required origin
cochain at the requested level. Hence these
programs give legitimate inputs in the existing
coordinate system and include their finite
cochain data.

If $T$ halts, (A.12.1) has finite support.
All of its local corrections are integral by
the definition of $e_j$, so its obstruction
vanishes. If $T$ does not halt, (A.12.1) is
the A.4 phantom input and its obstruction
is nonzero. A total vanishing evaluator for
these finite program descriptions would
therefore decide whether an arbitrary
Turing machine halts. The halting theorem
proves the second asserted impossibility.
This conclusion concerns the precise
uniform family (A.12.1); it does not assign
an unstated encoding to all possible
symbolic descriptions of completed inputs.
\end{proof}

\begin{statement}{A.13}{Proved}{The full finite-$S$ complex, with its integral differential}
With the basis and the Frobenius $u$ in A.7 and A.10, put
$M=\mathcal K$, $M^+=M^{c=1}$ and $d=(G_\alpha-1)|_{M^+}$.
There are compatible explicit quasi-isomorphisms
\[
 R\Gamma(\mathbb Z[1/3],M/3^n)
       \simeq[M^+/3^nM^+\xrightarrow{d}M/3^nM]
\]
in degrees zero and one, and
\[
 R\Gamma_{\rm cts}(G_S,M)\simeq[M^+\xrightarrow d M].
 \tag{A.13.1}
\]
The first lattice has rank $21$ and the second rank $44$.
The twenty-one nonzero Smith factors of $d$ are seven units,
nine factors three, and five factors nine. In particular
\[
\begin{aligned}
 H^0(G_S,M/3^n)&\simeq(\mathbb Z/3)^9
                    \oplus(\mathbb Z/3^{\min(n,2)})^5,\\
 H^1(G_S,M/3^n)&\simeq(\mathbb Z/3^n)^{23}
                    \oplus(\mathbb Z/3)^9
                    \oplus(\mathbb Z/3^{\min(n,2)})^5.
\end{aligned}
\tag{A.13.2}
\]
The full integral matrices are printed by
\texttt{code/a\_k\_lattice.py}; thus (A.13.1) specifies the complex
and not only its Euler characteristic.
\par\smallskip\noindent\emph{Manuscript comparison: \S51.16,
folios 733v--736v, scans 172--178, formulas (371)--(379).}
\end{statement}
\begin{proof}
Every finite-coefficient cocycle factors through $\Gamma$.
On $G_{K_0,S}$ its graph in the semidirect product of the coefficient
three-group and the action image is a homomorphism to a three-group,
so it kills the kernel defining $P_3$. This gives the required
factorization; it is compatible with $c$. Inflation is therefore
an isomorphism in degrees zero and one. Both groups have vanishing
higher three-primary cohomology for these coefficients, by A.1
and the free pro-three kernel of $\Gamma$ with its order-two
quotient. Thus inflation is a quasi-isomorphism.

Normalize a one-cocycle on $\Gamma$ to vanish on $c$, using
averaging over its group of order two. Its value on the free
generator $u$ is arbitrary in $M/3^n$, and uniquely determines
its value on $cuc$ and on the entire group. The remaining
coboundaries come precisely from $M^{c=1}/3^n$ and have value
$(G_\alpha-1)m$ at $u$. This construction is a cochain map from
the displayed two-term complex to continuous cochains: in degree
zero it is inclusion of the $c$-invariants; in degree one it sends
a vector to the unique normalized cocycle. It induces the stated
cohomology isomorphisms, hence is the claimed quasi-isomorphism.
Reduction and inverse limits commute with these maps. The invariant
lattice commutes with reduction because two is a unit.

For the matrix calculation, the columns of $1+G_c$ span $M^+$.
The program selects the twenty-one columns independent modulo three;
Nakayama's lemma makes them a $\mathbb Z_3$-basis. It applies
$G_\alpha-1$ and performs the same certified Smith elimination as
in A.7. The resulting valuations are seven zeros, nine ones,
and five twos. Kernel and cokernel modulo $3^n$ give (A.13.2).
The matrix has trivial integral kernel, consistent with
$H^0(G_S,M)=0$, and its integral cokernel is (A.7.1).
\end{proof}

\subsection{An exact two-prime decomposition at seven}

Put $\mathfrak q=(2-\zeta)$ and
$\overline{\mathfrak q}=(3+\zeta)$ in $\mathbb Z[\zeta]$.
Their residue fields are $\mathbb F_7$, with $\zeta=2$ at
$\mathfrak q$ and $\zeta=4$ at $\overline{\mathfrak q}$.
Write
\[
 S_7=\{\lambda,\mathfrak q,\overline{\mathfrak q},\infty\},
 \quad G=G_{K_0,S_7},\quad P_7=G(3),\quad
 D_{\mathfrak p}=G_{K_{0,\mathfrak p}}(3).
\]
Here $(3)$ denotes the maximal pro-three quotient. Choose an
embedding of the decomposition group at $\mathfrak q$ into $G$,
and choose the embedding at $\overline{\mathfrak q}$ by conjugating
it with the fixed complex conjugation $c$. The star in the following
statement is the free product in the category of pro-three groups.

\begin{statement}{A.14}{Proved}{The two local factors generate the entire global pro-three group}
The natural morphism
\[
 D_{\mathfrak q}*_{3}D_{\overline{\mathfrak q}}
       \longrightarrow P_7
 \tag{A.14.1}
\]
is an isomorphism. Each factor has the explicit pro-three
presentation
\[
 D_{\mathfrak q}=
 \langle\phi,\tau\mid\phi\tau\phi^{-1}=\tau^7\rangle_{3}
       \simeq\mathbb Z_3(1)\rtimes\mathbb Z_3,
 \tag{A.14.2}
\]
where $\phi$ is arithmetic Frobenius and $\tau$ is tame inertia.
The relative pro-three quotient of
$G_{\mathbb Q,\{3,7,\infty\}}$ over $\mathbb Q(\zeta)$ is
therefore
\[
 \Gamma_7\simeq
 (D_{\mathfrak q}*_{3}D_{\overline{\mathfrak q}})
           \rtimes\langle c\rangle,\qquad c^2=1,
 \tag{A.14.3}
\]
with $c$ exchanging the conjugately embedded factors. In particular
the local morphisms in (A.14.1) are injective; this is part of the
conclusion, not an assumption about realization of local extensions.
\par\smallskip\noindent\emph{Manuscript comparison: \S51.16,
folios 733v--736v, scans 172--178, formulas (371)--(379); the
one-prime higher-lifting problem in B.4.}
\end{statement}
\begin{proof}
We first prove exactly the comparison with a maximal pro-three
quotient that will be used. Let $H$ be any profinite group,
$H\twoheadrightarrow H(3)$ its maximal pro-three quotient, and
$J$ its kernel. A nonzero invariant character
$\eta:J\to\mathbb F_3$ would have a closed kernel normal in $H$.
The resulting quotient $H/\ker\eta$ would be an extension of the
pro-three group $H(3)$ by the group $\mathbb F_3$, hence a pro-three
group. The defining maximality would make its map from $H$ kill
$J$, a contradiction. Thus
\[
 H^1(J,\mathbb F_3)^{H(3)}=0,
 \qquad
 H^2(H(3),\mathbb F_3)\hookrightarrow H^2(H,\mathbb F_3).
 \tag{A.14.4}
\]
The second assertion is the inflation map in the five-term
Hochschild--Serre sequence. No assertion about all higher
cohomology of this quotient is needed.

Let $U=\operatorname{Spec}\mathbb Z[\zeta,1/(3\cdot7)]$.
The ring is a localization of a principal ideal domain, so
$\operatorname{Pic}(U)=0$. The Kummer sequence, the fixed
isomorphism $\mathbb F_3\simeq\mu_3$ sending $1$ to $\zeta$,
and Milne II.2.9 \cite{MilneADT} identify
\[
 H^1(G,\mathbb F_3)
    =\langle\zeta,\lambda,\pi,\bar\pi\rangle_{\mathbb F_3},
 \qquad \pi=2-\zeta,\quad\bar\pi=3+\zeta,
 \tag{A.14.5}
\]
and identify $H^2(G,\mathbb F_3)$ with
$\operatorname{Br}(U)[3]$. Milne II.2.1 and Remark II.2.2(a)
give its explicit invariant description
\[
 \operatorname{Br}(U)[3]=
 \{(x_\lambda,x_{\mathfrak q},x_{\bar{\mathfrak q}})
       \in(\tfrac13\mathbb Z/\mathbb Z)^3:
          x_\lambda+x_{\mathfrak q}+x_{\bar{\mathfrak q}}=0\}.
 \tag{A.14.6}
\]
Consequently restriction from the full group $G$ to the two
local groups at seven is an isomorphism in degree two with these
coefficients. Its inverse on invariants is
$(x,y)\mapsto(-x-y,x,y)$. In particular it is injective.
Combined with (A.14.4) and the natural local inflation maps, this
proves the injection
\[
 H^2(P_7,\mathbb F_3)\hookrightarrow
 H^2(D_{\mathfrak q},\mathbb F_3)
       \oplus H^2(D_{\overline{\mathfrak q}},\mathbb F_3).
 \tag{A.14.7}
\]
Indeed a class restricting to zero in both pro-three local groups
inflates to a full global class with zero restrictions, and is
therefore zero by (A.14.6) and (A.14.4).

For degree one, the two local Kummer groups each have dimension
two: one uniformizer class and one residue-unit class. The
unramified Kummer values of $\zeta$ and $\lambda$ are
\[
 \begin{array}{c|cc}
 &\kappa_\zeta&\kappa_\lambda\\\hline
 \phi_{\mathfrak q}&2&0\\
 \phi_{\bar{\mathfrak q}}&2&2
 \end{array}\quad\text{in }\mathbb F_3.
 \tag{A.14.8}
\]
These values follow by taking the square of the residue and
expressing it as a power of the residue of $\zeta$. For example
$\lambda=-1$ at $\mathfrak q$, while $\lambda=4=\zeta$ at
$\overline{\mathfrak q}$. The determinant of (A.14.8) is
$4=1$ modulo three. The valuations of $\pi$ and $\bar\pi$
at the two places form the identity matrix. Thus the restriction
\[
 H^1(P_7,\mathbb F_3)\longrightarrow
 H^1(D_{\mathfrak q},\mathbb F_3)
       \oplus H^1(D_{\overline{\mathfrak q}},\mathbb F_3)
 \tag{A.14.9}
\]
is an isomorphism. With Frobenius lifts fixing the respective
uniformizer roots, the complete matrix, in row order
$(\phi_{\mathfrak q},\tau_{\mathfrak q},
\phi_{\bar{\mathfrak q}},\tau_{\bar{\mathfrak q}})$ and column
order $(\zeta,\lambda,\pi,\bar\pi)$, is
\[
 \begin{pmatrix}
 2&0&0&2\\0&0&1&0\\2&2&1&0\\0&0&0&1
 \end{pmatrix}.
 \tag{A.14.10}
\]
Here both tame generators are oriented using the same global
$\zeta$. Hence $c\tau_{\mathfrak q}c^{-1}
=\tau_{\bar{\mathfrak q}}^{-1}$. Using the conjugately
transported second generator instead negates the last row; it
does not change invertibility. This specifies the orientation
in the factor exchange in (A.14.3).

For completeness we give the group-theoretic detection argument.
Let $B=D_{\mathfrak q}*_{3}D_{\overline{\mathfrak q}}$ and
$f:B\to P_7$ the natural morphism. Its image is compact and
therefore closed. A proper closed subgroup of a pro-three group
is contained in an open subgroup of index three after passing
to a finite quotient. Such a subgroup would yield a nonzero
character of $P_7$ trivial on the image of $f$, contrary to
(A.14.9). Thus $f$ is surjective.

Restriction $H^2(B,\mathbb F_3)\to\bigoplus_{\mathfrak p\mid7}
H^2(D_{\mathfrak p},\mathbb F_3)$ is injective. To see this
directly, a degree-two class is represented by a central
extension of $B$ by $\mathbb F_3$. Zero restriction provides
homomorphic splittings on the two factors. The extension is a
pro-three group, so the universal property of the free pro-three
product combines these into a splitting on $B$. Equation
(A.14.7) in particular shows that
$f^*:H^2(P_7,\mathbb F_3)\to H^2(B,\mathbb F_3)$ is injective.
Put $R=\ker f$. Its five-term sequence and (A.14.9) now give
$H^1(R,\mathbb F_3)^{P_7}=0$.

A nontrivial closed normal subgroup $R$ of a pro-three group $B$
always has a nonzero invariant character to $\mathbb F_3$.
For a proof, choose a finite three-group quotient $B_0$ in
which its image $R_0$ is nontrivial. The nonzero vector space
$R_0/\Phi(R_0)$ is a finite module for the three-group $B_0$.
Its coinvariants are nonzero: the augmentation ideal of
$\mathbb F_3[B_0]$ is nilpotent, so equality of that ideal times
the module with the whole module would force the module to be
zero. A nonzero linear functional on the coinvariants gives
an invariant character of $R$, a contradiction. Thus $R=1$
and (A.14.1) is proved.

At residue characteristic seven every finite extension of
three-power degree is tame. The maximal pro-three extension
is obtained from the unramified three-power tower and compatible
three-power roots of a uniformizer. Its tame inertia is
$\mathbb Z_3(1)$, its unramified quotient is $\mathbb Z_3$,
and arithmetic Frobenius acts on inertia by seventh powers.
Choosing Frobenius to fix the uniformizer roots gives exactly
(A.14.2). Finally the kernel of $G\to P_7$ is characteristic
in $G$ and stable under $c$. The chosen order-two element $c$
therefore splits the relative quotient over $\mathbb Q$ and
exchanges the conjugately chosen local factors, proving
(A.14.3).

The elementary residue and matrix computations, including the
orientation change, are reproduced in
\texttt{code/a\_two\_prime.py} and its output file.
\end{proof}

\begin{statement}{A.15}{Proved}{The exact local equation for every section supported at three and seven}
Let $Q$ be any of the geometric pro-three quotients of $N$
under consideration, with the action defined by the unit
tangential section at $O$. Write $F$ for the action of
$\phi_{\mathfrak q}$ on $Q$. Sections unramified outside
$\{3,7,\infty\}$ have the exact parameterization
\[
 \operatorname{Sec}_{\{3,7\}}(Q)\simeq
 \{(a,b)\in Q^2:aF(b)a^{-1}=b^7\}/Q^c,
 \tag{A.15.1}
\]
where
\[
 h\cdot(a,b)=\bigl(h^{-1}aF(h),\ h^{-1}bh\bigr)
       \quad(h\in Q^c).
 \tag{A.15.2}
\]
The parameterization is natural under all the quotient maps.
For a surjective equivariant morphism $R\twoheadrightarrow Q$
among these geometric pro-three groups, a class represented by
$(a,b)$ lifts exactly when there are lifts $A,B\in R$ satisfying
\[
 A F(B)A^{-1}=B^7.
 \tag{A.15.3}
\]
The same statement holds for $R=N$. Moreover a fixed class-two
section supported at these primes lifts to $N$ exactly when
it lifts to every lower-central quotient $N/\gamma_rN$.
The equivalence includes compatible full lifts; no compatibility
of separately chosen representatives is assumed.
\par\smallskip\noindent\emph{Manuscript comparison: \S51.16,
folios 733v--736v, scans 172--178, formulas (371)--(379); the
one-prime higher-lifting problem in B.4.}
\end{statement}
\begin{proof}
The action on $N$ defined by the tangential basepoint is
unramified away from three. Over $K_0$, its action on
$N/\Phi(N)=V/3$ has three-group image, as proved in A.1.
The kernel of
$\operatorname{Aut}(N)\to\operatorname{Aut}(N/\Phi(N))$
is pro-three. One can see this on the finite lower-exponent-three
quotients: successive kernels embed into homomorphism groups
between finite $\mathbb F_3$-vector spaces, and are three-groups;
their inverse limit is the stated kernel. Hence the action image
over $K_0$ on $N$, and on each $Q$, is pro-three. The graph of
any $Q$-valued cocycle on $G_{K_0,S_7}$ is a homomorphism into
the pro-three semidirect product of $Q$ with this action image.
It factors through $P_7$ by its defining maximality. Thus the
sections in question factor through the exact group $\Gamma_7$
of A.14.

Every cocycle on $\langle c\rangle$ with values in a pro-three
group is a coboundary. For finite three-groups this follows
inductively along $c$-stable central elementary abelian subgroups:
the induction leaves an abelian cocycle, and division by two
solves it. For a pro-three group apply the finite assertion to
$c$-stable finite quotients; the nonempty closed sets of
conjugating elements have the finite-intersection property, so
compactness supplies a common solution. We may therefore
represent any section by a cocycle $\beta$ with $\beta(c)=1$.
The remaining conjugacies are precisely by $Q^c$.

The restriction of a normalized cocycle to the first local
factor is determined by
$a=\beta(\phi_{\mathfrak q})$ and
$b=\beta(\tau_{\mathfrak q})$. Tame inertia acts trivially
on $Q$, because the reference action is unramified at seven.
The single relation in (A.14.2) therefore becomes
$aF(b)a^{-1}=b^7$. Conversely any pair satisfying this relation
gives a continuous splitting over the local pro-three group:
send the two generators to $(a,\phi)$ and $(b,\tau)$ in
$Q\rtimes D_{\mathfrak q}$ and use its pro-three presentation.
On the conjugate factor define
\[
 \beta(cdc^{-1})=c\bigl(\beta(d)\bigr)
       \quad(d\in D_{\mathfrak q}).
\]
The free-product property extends these local splittings
uniquely over $P_7$, and the displayed identity extends them
with $\beta(c)=1$ over $\Gamma_7$. Conjugation of the two
generator values by $h\in Q^c$ gives exactly (A.15.2).
This proves (A.15.1) and its naturality.

The fixed-point map $R^c\to Q^c$ is surjective. Indeed, for a
lift $r$ of a fixed element of $Q$, the discrepancy
$r^{-1}c(r)$ is a cocycle in the pro-three kernel of $R\to Q$.
The preceding coprime argument removes it by an element of that
kernel. Thus changes of normalized representatives can be
lifted, and (A.15.3) gives the stated criterion for the
conjugacy class, not only for an individually chosen pair.

Finally fix the pair $(a,b)$ in the class-two quotient. In
the compact space of its lifts to $N^2$, the condition
\[
 A F(B)A^{-1}B^{-7}\in\gamma_rN
\]
defines a closed subset. Lifting at level $r$ makes this set
nonempty, by lifting the two elements from $N/\gamma_rN$
to $N$; the representative issue has just been resolved.
These sets are nested as $r$ increases. Compactness provides
a pair belonging to all of them. The intersection
$\bigcap_r\gamma_rN$ is trivial: any nontrivial element
survives in a finite three-group quotient, and every such
quotient is nilpotent. The resulting pair satisfies the full
equation in $N$, proving the last assertion.
\end{proof}

Equations (A.15.1)--(A.15.3) remove the additional global
embedding problem at the specified prime seven. They do not
evaluate the full local equation for every ramified mixed
inertia vector. In particular A.14--A.15 alone do not settle
the all-nilpotent assertion in B.4 or the general ramified
metabelian classification in A.6.

\begin{statement}{A.16}{Proved}{A positive-density set with the same exact one-prime reduction}
Define the following explicit set of rational primes:
\[
 \Sigma=\{q:q\equiv4\text{ or }7\pmod9,
               \quad3^{(q-1)/3}\not\equiv1\pmod q\}.
 \tag{A.16.1}
\]
It has natural density $2/9$. For every $q\in\Sigma$, write
$(q)=\mathfrak q\overline{\mathfrak q}$ in $\mathbb Z[\zeta]$.
The natural local morphisms give an isomorphism
\[
 G_{K_0,\{\lambda,\mathfrak q,\overline{\mathfrak q},\infty\}}(3)
   \simeq D_{\mathfrak q}*_{3}D_{\overline{\mathfrak q}},
 \quad D_{\mathfrak q}=
       \langle\phi,\tau\mid\phi\tau\phi^{-1}=\tau^q\rangle_3.
 \tag{A.16.2}
\]
Its relative quotient over $\mathbb Q$ is the semidirect product
by complex conjugation exchanging the factors. The exact section
parameterization and lifting assertions of A.15 hold with $7$
replaced throughout by $q$. In particular the defining equation
is $aF_q(b)a^{-1}=b^q$, with the same action of $Q^c$ on pairs.
The primes in $\Sigma$ up to $1000$ are
\[
\begin{gathered}
7,13,31,43,79,97,139,157,211,223,229,241,\\
277,283,313,331,337,349,373,409,421,457,463,\\
571,601,607,673,691,709,733,751,769,823,859,877,907.
\end{gathered}
 \tag{A.16.3}
\]
\par\smallskip\noindent\emph{Manuscript comparison: \S51.16,
folios 733v--736v, scans 172--178, formulas (371)--(379); the
one-prime higher-lifting problem in B.4.}
\end{statement}
\begin{proof}
Let $q\equiv1\pmod3$ be a rational prime different from three,
choose one residue root of $\zeta$ at $\mathfrak q$, and write
\[
 a=\kappa_\zeta(\phi_{\mathfrak q})
       =\frac{q-1}{3}\bmod3,
 \qquad b=\kappa_\lambda(\phi_{\mathfrak q})\in\mathbb F_3.
\]
Conjugation sends $\zeta$ to $\zeta^{-1}$ and
$\lambda$ to $-\zeta^2\lambda$. Its action on the Kummer
coefficients $\mu_3$ is multiplication by minus one. Thus the
conjugately chosen Frobenius at $\overline{\mathfrak q}$ has
Kummer vector $(a,a-b)$, with both characters measured using
the same global $\zeta$. In particular the determinant of the
two unramified-unit rows is
\[
 \det\begin{pmatrix}a&b\\a&a-b\end{pmatrix}
       =a(a+b)\quad\text{in }\mathbb F_3.
 \tag{A.16.4}
\]
The identity $\lambda^2=-3\zeta$ gives
$\kappa_3=2b-a=-a-b$. Hence (A.16.4) is nonzero precisely
when $a\ne0$ and the residue of $3$ is not a cube. These are
exactly the two elementary conditions in (A.16.1).

Since $\mathbb Z[\zeta]$ is principal, choose a generator $\pi$
of $\mathfrak q$ and its conjugate for
$\overline{\mathfrak q}$. The four S-unit Kummer classes are
$\zeta,\lambda,\pi,\bar\pi$. Their inertia submatrix is the
identity, while (A.16.4) is their unramified-unit submatrix.
For $q\in\Sigma$ this proves the degree-one isomorphism used
in A.14. The Brauer invariant proof is unchanged, since exactly
three finite places are inverted. All group-theoretic maps and
kernel arguments in A.14 consequently apply to these actual
local morphisms and prove (A.16.2). The local tame presentation
has exponent $q$; $q\equiv1\pmod3$ ensures that its Frobenius
action on three-inertia factors through $\mathbb Z_3$.
The construction of normalized cocycles and the compactness
argument in A.15 now give all the stated section assertions.

For the density assertion put
\[
 L=K_0(\zeta_9,\sqrt[3]{\lambda}).
\]
The Kummer classes of $\zeta$ and $\lambda$ are independent:
valuation at $\lambda$ first forces the latter exponent to be
zero modulo three, and $\zeta$ is not a cube in $K_0$, as its
unit group is $\mu_6$. Hence $[L:K_0]=9$ and its Galois group
is $\mathbb F_3^2$, with the two displayed Kummer coordinates.
The relation $\bar\lambda=-\zeta^2\lambda$ makes $L$ stable
under complex conjugation: a conjugate cube root is
$-\zeta_9^2\sqrt[3]{\lambda}$. Thus $L/\mathbb Q$ is
Galois of degree eighteen. The chosen complex conjugation
acts on the normal subgroup by
\[
 (a,b)\longmapsto(a,a-b).
\]
The admissible vectors are exactly
$(1,0),(1,1),(2,0),(2,2)$; this four-element subset is stable
under conjugation in the full group. A rational prime outside
three belongs to $\Sigma$ exactly when its Frobenius class
belongs to this subset. Chebotarev
\cite[V.3.23]{MilneCFT} therefore gives natural density
$4/18=2/9$. The exhaustive elementary check through $1000$,
including both residue roots, the conjugation identity, the
$\kappa_3$ identity and the determinant test, is in
\texttt{code/a\_two\_prime.py} and its output.
\end{proof}

\subsection{A strict restriction on ramification in the higher quotients}

\begin{statement}{A.17}{Proved}{Tame abelian inertia in a metabelian section lies on one cusp line}
Let $k$ be a finite extension of $\mathbb Q_q$ with $q\ne3$,
and let
\[
 s:G_k\longrightarrow(N/N'')\rtimes G_k
\]
be a continuous section. For a generator of tame three-inertia,
write $t$ for its geometric value and $w$ for its image in
$V=N^{\rm ab}$. Then
\[
 w=0\quad\text{or}\quad
 w\in\mathbb Z_3 e_d\setminus\{0\}
       \text{ for a unique cusp }d\in D(k).
 \tag{A.17.1}
\]
In particular the assertion T1 of Lead 2, refuted for class
two by A.5, holds for the tame abelian inertia of every
metabelian section through the actual quotient
$N/N''\twoheadrightarrow N/\gamma_3N$.
\par\smallskip\noindent\emph{Manuscript comparison: \S51.16,
folios 733v--736v, scans 172--178, formulas (371)--(379).}
\end{statement}
\begin{proof}
Put $Q=N/N''$ and let $Q_k$ denote the size of the residue
field of $k$. The good marked model and its integral tangent
make the reference action of inertia on $N$ trivial.
Consequently the restriction of the cocycle of $s$ to inertia
is a homomorphism; wild inertia and the tame factors prime
to three have trivial image in $Q$. For an arithmetic
Frobenius lift $F=s(\phi_k)$ the exact relation in the
arithmetic quotient is
\[
 FtF^{-1}=t^{Q_k}.
 \tag{A.17.2}
\]
It follows on abelianization that $w$ is in $W$: its image
in $T_3E$ is a $Q_k$-eigenvector, and that Tate module has
Frobenius eigenvalues of complex absolute value
$\sqrt{Q_k}$ and is torsion-free.

For $n\ge1$ define the actual subgroups of the full geometric
group
\[
 M_n=\ker(N\to V/3^n)=N'N^{3^n},\qquad
 H_n=\{g\in N:g\bmod M_n\in\langle w\bmod3^n\rangle\}.
 \tag{A.17.3}
\]
Here products and powers denote closed subgroups. The group
$H_n$ is open and normal in $N$, contains $N'$, and is stable
under the action induced by the local section. Indeed inertia
acts trivially on $V$, while Frobenius sends the generating
vector $w$ to $Q_kw$, and $Q_k$ is a three-adic unit.
Crucially, the inclusion $N'\subset H_n$ gives
\[
 N''=[N',N']\subset[H_n,H_n],\qquad
 H_n^{\rm ab}\xrightarrow{\sim}(H_n/N'')^{\rm ab}.
 \tag{A.17.4}
\]
The displayed isomorphism is induced by the quotient map;
its kernel is explicitly
$N''[H_n,H_n]/[H_n,H_n]=0$.

The subgroup $(H_n/N'')s(G_k)$ of $Q\rtimes G_k$ is open.
Pull it back through the actual morphism
$N\rtimes G_k\to Q\rtimes G_k$. The resulting arithmetic
subgroup defines a finite etale cover $Y_n\to U_k$, with
geometric group $H_n$. This cover has the smooth proper tame
compactification needed for the weight argument in B.7.
We spell out why that argument applies to the present
metabelian section, without assuming a full section lift.
The cocycle values on arithmetic inertia are powers of $t$
and therefore lie in $H_n/N''$. Thus the pulled-back
subgroup contains the original arithmetic inertia.
Original inertia centralizes $N$, so its action on every
coset in $N/H_n$ is trivial. Prime-to-$q$ specialization
therefore descends the cover to the good marked model.
At each boundary point its compactification is the tame
normalization $z^e=x$ after an unramified extension, with
$e$ prime to $q$; it is smooth there and etale away from
the boundary. This is the specialization comparison used
in B.7, with the same hypotheses.

Choose any lift $\widetilde t\in H_n$ of $t$, and any
lift $\widetilde F$ of the Frobenius value in this arithmetic
cover subgroup. Relation (A.17.2) says
\[
 \widetilde F\widetilde t\widetilde F^{-1}
                      \widetilde t^{-Q_k}\in N''.
\]
By (A.17.4), the class of $\widetilde t$ in $H_n^{\rm ab}$
is therefore an exact $Q_k$-eigenvector. Let $X_n$ denote
the smooth proper compactification of $Y_n$. Its boundary
sequence is the integral exact sequence
\[
 0\longrightarrow
 \mathbb Z_3(1)[X_n\setminus Y_n]/\mathbb Z_3(1)\delta_n
 \longrightarrow H_n^{\rm ab}
 \longrightarrow T_3J_{X_n}\longrightarrow0.
 \tag{A.17.5}
\]
Frobenius on the last term has complex absolute values
$\sqrt{Q_k}$, by the smooth proper finite-field weight
theorem used in B.7. It has no $Q_k$-eigenvector even
integrally, because the lattice is torsion-free. Thus the
class of $\widetilde t$ is in the actual boundary submodule
in (A.17.5).

The finite cyclic three-group $H_n/M_n$ is generated by
the image of $\widetilde t$. The boundary submodule
therefore maps onto this group. At least one boundary
inertia subgroup maps onto it: otherwise all its images
would be contained in the unique maximal subgroup of this
cyclic group. Such a boundary generator is conjugate in
$N$ to $c_d^e$ for some cusp $d\in D$ and some positive
integer $e$. Its image in $V/3^n$ is $e e_d$.
The assertion that this element generates $H_n/M_n$
means that $w\bmod3^n$ is a scalar multiple of
$e_d\bmod3^n$. If the cyclic group is trivial, the same
conclusion holds with scalar zero.

There are only nine cusps. One cusp occurs for an unbounded
subsequence of $n$, and its line $\mathbb Z_3e_d$ is closed
in $V$. Equivalently, $e_d$ is primitive and the quotient
of $V$ by that line is a free $\mathbb Z_3$-module, so
congruence membership for unbounded $n$ gives exact
membership. This proves $w\in\mathbb Z_3e_d$.
Distinct cusp lines have zero intersection: in
$\mathbb Z_3[D]/\mathbb Z_3\delta$ an equality
$a e_d=b e_{d'}$ with $d\ne d'$ has zero coefficients
at the other seven cusps, forcing the diagonal coefficient
and then $a,b$ to be zero. Thus the cusp is unique when
$w\ne0$. Frobenius sends $e_d$ to $Q_ke_{\phi_k d}$.
Together with $F(w)=Q_kw$, this forces $\phi_k d=d$.
The cusp scheme is unramified at $q$, so this is precisely
$d\in D(k)$, proving (A.17.1).
\end{proof}

\begin{statement}{A.18}{Proved}{Class-two mixed examples have no metabelian or marked-depth lift}
Let $I_{\rm mark}$ be the closed normal subgroup generated
by the inertia at $Q_c$, so that the marked depth quotient
in the task is $N/\gamma_3(I_{\rm mark})$.
The class-two sections of A.5 and A.6 have no lift to either
\[
 N/N''\quad\text{or}\quad N/\gamma_3(I_{\rm mark}).
 \tag{A.18.1}
\]
This holds for every class-two lift of their specified
$N^{\rm cc}$ sections, and continues to hold after every
finite number-field extension. Thus their section germs
also have no lift to these two quotients.

More generally, a local section of the marked depth quotient
over a finite extension $k/\mathbb Q_q$, $q\ne3$, has
tame abelian inertia $w$ whose image in
$V/\mathbb Z_3e_{Q_c}$ is zero or lies on a single cusp
line belonging to $D\setminus\{Q_c\}$ in that quotient.
The images of both section maps in (A.18.1) to class two
are therefore proper; the entire unramified family of
A.20 still belongs to both images.
\par\smallskip\noindent\emph{Manuscript comparison: \S51.16,
folios 733v--736v, scans 172--178, formulas (371)--(379).}
\end{statement}
\begin{proof}
The two comparison maps to class two are genuine quotient
maps. Indeed $N''\subset\gamma_4N\subset\gamma_3N$,
and $\gamma_3(I_{\rm mark})\subset\gamma_3N$ because
$I_{\rm mark}\subset N$.
At each prime used in A.5 or A.6, the specified abelian
tame inertia is a nonzero scalar multiple of
\[
 w_0=e_O+e_P.
\]
It is on no cusp line. In the basis of $W$ given by the
eight cusps other than $O$, its coefficient at $P$ is zero
and its coefficients at the other seven cusps are all
minus one. This vector is neither a multiple of a basis
vector nor of $e_O$, whose eight coefficients are all
minus one. Multiplication by a nonzero three-adic scalar
does not change this conclusion. A.17 consequently rules
out every metabelian lift.

For the marked quotient use the explicit filling morphisms
\[
 N/\gamma_3(I_{\rm mark})\twoheadrightarrow N/I_{\rm mark}
 \xrightarrow{\sim}
 \pi_1(E_{\overline{\mathbb Q}}setminus
                    (D\setminus\{Q_c\}))^{(3)}.
 \tag{A.18.2}
\]
The second map is induced by filling the puncture $Q_c$.
It is an isomorphism because a finite three-group etale
cover of $U$ extends etale across that puncture exactly
when its inertia there is trivial. Normal closure kills
precisely that inertia condition, and the finite covers
on the two sides are thereby identified. The same map
on abelianizations is
\[
 V\longrightarrow V/\mathbb Z_3e_{Q_c},
\]
as follows from the boundary exact sequence, with the
single summand at $Q_c$ removed. These morphisms preserve
the arithmetic action and the tangent at $O$.

The resulting open elliptic curve has eight cusps and good
marked reduction at $q\ne3$. The proof of B.7 applies
verbatim to its full geometric pro-three group: it uses
only the finite-cover boundary sequence, smooth proper
weights, and compactness of the lifted cusp set. Hence
every nonzero tame inertia of the projected full section
lies in one cusp inertia subgroup. Its abelian image is
zero or on that cusp line, which proves the more general
assertion of the statement.

In the filled curve the images of the seven cusps other
than $O,Q_c$ form a basis of its boundary lattice. The
image of $w_0$ has coefficient zero at $P$ and minus one
at the other six basis cusps. It is not zero and is on no
cusp line by the same coordinate comparison. Thus (A.18.2)
rules out the marked-depth lifts of A.5 and A.6 as well.

Under a finite extension of local fields the tame inertia
generator maps to a nonzero integer multiple of the old
one, with that integer equal to the ramification index
after compatible choices. The new vector is consequently
a nonzero scalar multiple of the same mixed vector;
changes of decomposition group merely conjugate and
permute its two distinct cusps. The two line exclusions
remain valid over every such extension, including those
of degree divisible by three. A hypothetical lift over
a finite number field would give a forbidden local lift
at a place above one of the specified ramified primes.
This proves the finite-extension and germ assertions.
Finally the full unramified lifts of A.20 project through
both quotient maps, while A.5 supplies a class outside
both images. The images are therefore nonempty and proper.
The integral line exclusions before and after filling are
independently certified by unit two-by-two minors in
\texttt{code/a\_higher\_checks.py} and its output.
\end{proof}

\begin{statement}{A.19}{Proved}{The complete local metabelian image at every good prime}
Let $k/\mathbb Q_q$ be finite with $q\ne3$, put
$Q=N/N''$, and retain the unit tangential origin at $O$.
Every continuous section of $Q\rtimes G_k$ is either
unramified or contained in the decomposition group of a
unique lifted $k$-rational cusp. It is unramified exactly
when its abelianization is unramified.

More precisely, with $F$ denoting the reference arithmetic
Frobenius action and $Q_k$ the residue-field cardinality,
the local section set is the disjoint union
\[
 \operatorname{Sec}_k(Q)=
 Q/\{a\sim h^{-1}aF(h):h\in Q\}
 \ \sqcup\!
 \coprod_{d\in D(k)}
 \{u\in k^{\times\wedge_3}:v_k(u)\ne0\}.
 \tag{A.19.1}
\]
The first term consists of unramified sections, with $a$
their geometric Frobenius value. The term at $d$ uses its
actual unit tangent section $s_{d,1}$ and the map
$g\mapsto c_d^{\kappa(u)(g)}s_{d,1}(g)$, followed by
the quotient to $Q$. Its parameter group is
\[
 k^{\times\wedge_3}
      \simeq\mathbb Z_3\oplus
                 \mu_{3^{v_3(Q_k-1)}}(k).
\]
The morphism
$\operatorname{Sec}_k(N)\to\operatorname{Sec}_k(Q)$
is surjective and is bijective on the ramified parts.
The ramified image of
$\operatorname{Sec}_k(Q)\to
\operatorname{Sec}_k(N/\gamma_3N)$
is exactly the full cuspidal packet of B.8, with the same
parameters, including the central inertia and Frobenius
coordinates. In particular no extra local metabelian
section remains over a class-two section with a noncuspidal
central coordinate.
\par\smallskip\noindent\emph{Manuscript comparison: \S51.16,
folios 733v--736v, scans 172--178, formulas (371)--(379).}
\end{statement}
\begin{proof}
We give the inverse-limit and integral arguments that
strengthen the line condition of A.17 to an actual cusp
subgroup. Put $M=Q'=N'/N''$. Let $t\in Q$ be the tame
inertia value, $w\in V$ its abelianization, and
$\widetilde t\in N$ a lift. Use the groups $M_n,H_n$
of (A.17.3), and set
\[
 H_\infty=\bigcap_nH_n
    =\{g\in N:g^{\rm ab}\in\mathbb Z_3w\}.
 \tag{A.19.2}
\]
All these groups contain $N'$. The same good-reduction
cover argument as in A.17 makes the image of
$\widetilde t$ in every $H_n^{\rm ab}$ lie in its integral
boundary submodule $B_n$.

First we justify the precise inverse limits. For any
descending sequence of closed subgroups of a profinite
group, their images in a fixed finite quotient eventually
stabilize to the image of their intersection. This follows
by compactness: a compatible nonempty sequence of fibers
over the fixed finite quotient has nonempty intersection.
The images of their closed commutator subgroups therefore
stabilize to the commutator of that intersection's image.
Applying this to the $H_n$ and testing in every finite
quotient of $N$ proves
\[
 \bigcap_n H_n'=H_\infty',\qquad
 H_\infty^{\rm ab}\xrightarrow{\sim}
                    \varprojlim_n H_n^{\rm ab}.
 \tag{A.19.3}
\]
For surjectivity of the second map, representatives of a
compatible family define nested compact cosets of $H_n'$;
their intersection gives a representative in $H_\infty$.
Thus the image of $\widetilde t$ lies in
$\varprojlim B_n\subset H_\infty^{\rm ab}$.

Let $C_n(e)$ be the geometric cusps of the cover $Y_n$
above a base cusp $e$. Since $H_n$ contains $N'$, these
finite sets and their inertia exponents are explicitly
\[
 C_n(e)=V/(\mathbb Z_3w+\mathbb Z_3e_e+3^nV),
 \qquad
 j_n(e)=\min\{j\ge1:j e_e\in
                         \mathbb Z_3w+3^nV\}.
 \tag{A.19.4}
\]
The minimum is a positive power of three. A basis vector
at a cusp in $C_n(e)$ represents a conjugate of
$c_e^{j_n(e)}$. For the map from level $n+1$ to level $n$,
the boundary-module transition is the pushforward on cusp
sets multiplied by $j_{n+1}(e)/j_n(e)$. This follows
directly from the map on tame inertia generators, so the
transition is specified integrally, including ramification.
There is an exact sequence of compact modules
\[
 0\to\mathbb Z_3(1)\delta_n\to
    \bigoplus_{e\in D}\mathbb Z_3(1)[C_n(e)]
       \to B_n\to0.
 \tag{A.19.5}
\]
The diagonal transition multiplies by $[H_n:H_{n+1}]$:
over each cusp, the number of its preimages times their
ramification index equals the cover degree.

For every inverse sequence of compact abelian groups the
first derived limit is zero. Indeed, in the standard
difference equation $x_n-f_n(x_{n+1})=y_n$, any finite
set of equations can be solved backwards. Their solution
sets are closed in the compact product of the groups and
have the finite-intersection property, which solves all
equations simultaneously. This applies to the diagonal
groups in (A.19.5), without asserting a Mittag--Leffler
property for their underlying discrete groups.

Suppose first that $w=0$. Then $H_\infty=N'$,
$j_n(e)=3^n$ for every cusp, and every cusp-module
transition in (A.19.5) has a factor three. Its inverse
limit is zero: each coefficient at a fixed level is
divisible by arbitrarily high powers of three. The
diagonal transitions have the same vanishing inverse
limit, and the preceding derived-limit argument gives
$\varprojlim B_n=0$. By (A.19.3) the image of
$\widetilde t$ in $N'/N''=M$ is zero. Hence
\[
                       w=0\quad\Longrightarrow\quad t=1.
 \tag{A.19.6}
\]

Now suppose $w\ne0$. By A.17 there is a unique
$d\in D(k)$ and $a\in\mathbb Z_3\setminus\{0\}$ with
$w=a e_d$. Write $a=3^s u$, $u\in\mathbb Z_3^\times$.
For $n>s$ the exponents in (A.19.4) are
\[
 j_n(d)=3^s,\qquad j_n(e)=3^n\quad(e\ne d).
 \tag{A.19.7}
\]
For the second equality, distinct cusp vectors form a
primitive independent pair in $V$: this is immediate in
$\mathbb Z_3[D]/\mathbb Z_3\delta$ by choosing a third
cusp to eliminate. Thus no smaller power of $e_e$ can
belong to the line generated by $a e_d$ modulo $3^n$.
The transition degree is $3^9$ for $n>s$. Consequently
the diagonal inverse limit is zero, all summands at
$e\ne d$ have zero inverse limit, and the remaining
summand gives the exact isomorphism
\[
 \varprojlim_n B_n\simeq
 \mathbb Z_3(1)[[A_d]],\qquad A_d=V/\mathbb Z_3e_d.
 \tag{A.19.8}
\]
Here the right side means the completed permutation module
$\varprojlim_n\mathbb Z_3(1)[A_d/3^nA_d]$, with
pushforward transition maps. Its basis measure at a coset
represented by $g\in N$ maps to the class of
$g c_d^{3^s}g^{-1}$ in $H_\infty^{\rm ab}$.
The isomorphism follows from (A.19.5), including its
diagonal term and the vanishing of that term's first
derived limit; no closure or rationalization replaces
the integral boundary module.

The actual Frobenius action on this module is $Q_k$
times an affine permutation of $A_d$. To specify its
origin, let $v_F\in V$ be the geometric Frobenius value
of the given section and use the cusp-path cocycle of
A.8, so that reference Frobenius sends
$c_d$ to $b_{F,d}^{-1}c_d^{Q_k}b_{F,d}$.
The affine map on the boundary cosets is
\[
 \Psi(x)=F_Vx+\overline{v_F-b_{F,d}}\quad(x\in A_d).
 \tag{A.19.9}
\]
The bar means the quotient by $\mathbb Z_3e_d$.
This proves the sign and the tangential origin of the
translation: the section contributes left translation
by $v_F$, and the reference path contributes
$-b_{F,d}$. Changing the based path conjugates this
affine description by a translation and leaves the
underlying boundary action unchanged.

We use the following elementary integral-measure fact.
For a finite free $\mathbb Z_3$-module $A$, let
$\Psi(x)=Bx+v$ be an affine automorphism with
$B^m-1$ invertible over $\mathbb Q_3$ for every positive
integer $m$. Then
\[
 \mathbb Z_3[[A]]^{\Psi}=
 \begin{cases}
 \mathbb Z_3\delta_{x_0},&x_0=(1-B)^{-1}v\in A,\\
 0,&x_0\notin A.
 \end{cases}
 \tag{A.19.10}
\]
Here is a proof with the necessary uniformity. The sole
possible periodic point in $A\otimes\mathbb Q_3$ is
$x_0$, since translation by $x_0$ turns $\Psi$ into
$B$ and $B^m-1$ is invertible. A positive power
$\Theta=\Psi^e$ has pro-three closure in the affine
automorphism group: take $e$ to kill its finite action
on $A/3A$. Thus every orbit of $\Theta$ on a finite
quotient $A/3^nA$ has three-power length.
Fix a residue cell in $A$ not containing $x_0$, and
let the corresponding orbit at level $n$ have length
$3^r$. For each $h$, there is a level $m\ge n$ at
which no point of the cell is fixed modulo $3^m$ by
$\Theta^{3^{r+h}}$. Otherwise the nested nonempty
compact solution sets would give an actual periodic
point in that cell. Every orbit at level $m$ meeting
this cell therefore has length divisible by
$3^{r+h+1}$. An invariant coefficient is constant on
such an orbit, and its pushforward to the chosen
level-$n$ cell is multiplied by the ratio of orbit
lengths, divisible by $3^{h+1}$. Hence that level-$n$
coefficient is divisible by every power of three and
is zero. All cells away from $x_0$ have zero coefficient.
If $x_0\notin A$ every coefficient is zero. If
$x_0\in A$, compatibility leaves exactly one constant
mass on its residue cells, giving (A.19.10).

Apply (A.19.10) to (A.19.9). The eigenvalues of $F_V$
on $A_d\otimes\mathbb Q_3$ have complex absolute
values $Q_k$ or $\sqrt{Q_k}$; the rational isogeny
comparison is (A.1.4), followed by the quotient of its
boundary part by $e_d$. Hence no positive power has
eigenvalue one, as required. The class of
$\widetilde t$ is a $Q_k$-eigenvector in
$H_\infty^{\rm ab}$, so its measure in (A.19.8),
after removing the Tate twist, is invariant under
$\Psi$. Its total mass is $u$, because every boundary
basis measure has abelian image $3^s e_d$ whereas
$\widetilde t$ has image $a e_d$. Thus this mass is
nonzero. Formula (A.19.10) proves that $x_0\in A_d$
and that the measure is exactly $u\delta_{x_0}$.
Choose $g\in N$ representing $x_0$. Then
\[
 \widetilde t\equiv g c_d^a g^{-1}\pmod {H_\infty'}.
 \tag{A.19.11}
\]

This congruence gives actual conjugacy in the metabelian
group. The image $J$ of $H_\infty$ in $Q$ fits into
$1\to M\to J\to\mathbb Z_3a e_d\to1$ and splits
by the subgroup generated by $t$, or by $c_d^a$.
Since $M$ is abelian, its commutator subgroup is
\[
 J'=(T-1)M,
 \tag{A.19.12}
\]
where $T$ is conjugation by any element with abelian
image $a e_d$. The image is already closed, as the
continuous image of the compact group $M$. Write
$h=g c_d^a g^{-1}$ in $Q$. Equations (A.19.11)--
(A.19.12) give $t h^{-1}=(T-1)m$ for some $m\in M$,
in additive notation for $M$. The right side equals
$m^{-1}h m h^{-1}$. Hence $t=m^{-1}h m$ in $Q$:
the tame inertia lies in one actual cusp inertia
subgroup of the metabelian group.

We recall and justify the centralizer calculation
needed to conclude for the whole section. For every
nonzero integer $b$,
\[
                    C_Q(c_d^b)=\overline{\langle c_d\rangle}.
 \tag{A.19.13}
\]
This is Yamaguchi, Proposition 1.1.6
\cite{AYamaguchi2023}, with the full class of finite
three-groups and $m=2$. Each $c_d$ is a free generator
in some basis of $N$: in the surface presentation one
eliminates a different cusp, which is available because
there are nine cusps. For a direct proof in this case,
commutation in class two first makes the abelian image
of any centralizer element a multiple of $e_d$, since
$\bigwedge^2V$ is torsion-free. Remove that power of
$c_d$ to obtain an element of $M$. The completed
abelian-cover chain complex of a graph with ten free
generators embeds
\[
 M\hookrightarrow\mathbb Z_3[[V]]^{10},
 \qquad \partial(f_i)=\sum_i(T_i-1)f_i.
\]
One obtains this injection by taking the inverse limit
of the degree-one chain sequences of its finite
abelian covers; (A.19.3) with $H_\infty=N'$ identifies
the limiting first homology with $M$. This is also
Proposition 1.1.2 of the cited paper. The ring
$\mathbb Z_3[[V]]\simeq\mathbb Z_3[[X_1,\ldots,X_{10}]]$
is an integral domain, and $T_d^b-1$ is nonzero.
Its multiplication is therefore injective on $M$,
which proves (A.19.13).

After a geometric conjugation take $t=c_d^a$.
Its closed generated subgroup equals that generated
by $c_d^{3^s}$. Compare the actual Frobenius value
with the Frobenius value of the unit tangent section
at $d$. Both conjugate $t$ to $t^{Q_k}$, so their
geometric ratio centralizes a nonzero integral power
of $c_d$. By (A.19.13) that ratio lies in its inertia
group. Frobenius and inertia generate the local
arithmetic image topologically; hence the whole
section is contained in the cusp decomposition group.
The cusp lift is unique: distinct base cusps have
distinct abelian lines, and for two lifts of the same
cusp any conjugator identifying a nonzero inertia
element belongs to (A.19.13).

Sections in this decomposition group are exactly
$H^1(G_k,\mathbb Z_3(1))=k^{\times\wedge_3}$ with
the specified tangential origin. A nonzero valuation
is exactly nontrivial tame inertia. Their conjugacies
in $Q$ are exactly their conjugacies inside the cusp
decomposition group, by the same uniqueness and
centralizer argument. These packets therefore give
the second, disjoint term in (A.19.1).
For unramified sections the cocycle factors through
$G_k/I_k=\widehat{\mathbb Z}$. A geometric value
$a\in Q$ at arithmetic Frobenius determines it and
always defines a continuous section: the closed
procyclic subgroup generated by $(a,1)$ in
$Q\rtimes\widehat{\mathbb Z}$ maps isomorphically
onto $\widehat{\mathbb Z}$. Indeed the map from
$\widehat{\mathbb Z}$ sending $1$ to this element
composes with projection to the identity. Conjugation
is precisely the relation in the first term of
(A.19.1).

Finally choose any lift of an unramified Frobenius
value to $N$ and repeat this procyclic construction.
Every cuspidal packet already comes from the actual
full cusp decomposition group. Thus all local
metabelian sections lift to $N$. B.7 and the injective
packet parameterization give bijectivity on the
ramified parts. Projecting the same actual packets to
class two gives exactly B.8; its injective parameters
also show that both the central inertia and the
Frobenius coordinate are retained. This proves all
claims.
\end{proof}

\subsection{Complete global images and fibers for one auxiliary prime}

Put $A=N/\gamma_3N$ and, for a geometric quotient $B$ mapping
to $A$, put $J_B=\ker(B\to A)$. For $q\ne3$ let
$\mathscr L_q\subset\operatorname{Sec}_{\mathbb Q_q}(A)$
be the union of the unramified local section classes and
the ramified cuspidal packets in B.8. Thus this subset uses
the actual class-two inertia and Frobenius values and their
unit-tangent origins, including all central coordinates.

\begin{statement}{A.24}{Proved}{The complete one-prime metabelian image and its fibers for $q\in\Sigma$}
For each $q\in\Sigma$ of A.16, put $S_q=\{3,q,\infty\}$.
Then
\[
 \operatorname{Sec}_{S_q}(N)\longrightarrow
 \operatorname{Sec}_{S_q}(N/N'')
       \quad\text{is surjective}.
 \tag{A.24.1}
\]
For either $B=N$ or $B=N/N''$, the exact image in class
two is
\[
 \operatorname{im}\bigl(\operatorname{Sec}_{S_q}(B)
                  \to\operatorname{Sec}_{S_q}(A)\bigr)
   =\operatorname{res}_q^{-1}(\mathscr L_q).
 \tag{A.24.2}
\]
All fibers are as follows. A target outside the displayed
subset has empty fiber. A target unramified at $q$ has
precisely the unramified fiber of (A.20.2); every lift in
the present fiber is unramified outside three. For a target
whose restriction at $q$ is ramified and belongs to
$\mathscr L_q$, choosing one normalized lift gives a
bijection of its fiber with the set of right cosets
\[
                         J_B/J_B^c.
 \tag{A.24.3}
\]
Here $J_B^c$ is the fixed subgroup of the actual complex
conjugation, and it is not asserted to be normal.
The same description for the ramified fibers of
(A.24.1) is $N''/(N'')^c$.
\par\smallskip\noindent\emph{Manuscript comparison: \S51.16,
folios 733v--736v, scans 172--178, formulas (371)--(379);
Tasks A.6 and B.4 for this explicit infinite set of primes.}
\end{statement}
\begin{proof}
Use the exact relative free product of A.16 and normalize
all cocycles to have value one on $c$. A normalized target
section is determined by its pair $(a,b)$ on the chosen
local factor at $q$. A local full lift of its conjugacy
class can be made a lift of this exact pair: lift a
conjugator through the surjection of geometric groups and
conjugate the local lift. A.15 then extends that pair,
together with its $c$-conjugate pair and value one at $c$,
to a global section over $G_{\mathbb Q,S_q}$.
A.19 supplies such a local full lift for every local
metabelian section. This proves (A.24.1). B.7--B.8 and
A.19 give exactly the local subset $\mathscr L_q$, so
the same construction proves both directions of
(A.24.2).

If the target is unramified at $q$, its abelian inertia
is zero. The inertia of every full or metabelian lift is
then trivial, by B.7 or (A.19.6), respectively. Such a
lift factors through $G_{\mathbb Q,\{3,\infty\}}$:
the kernel of the map from $G_{\mathbb Q,S_q}$ to this
group is the closed normal subgroup generated by the
inertia groups above $q$, and the reference action is
also trivial on these groups. Thus its entire fiber is
exactly the one already determined in A.20.

Now fix a ramified target in the image and fix a full or
metabelian global lift $\widetilde\beta$ with
$\widetilde\beta(c)=1$ and exact projection $\beta$.
Such an exact representative exists because the
fixed-point map $B^c\to A^c$ is surjective, as proved in
A.15. The ramified local packet injects into class two
by B.8 and A.19. Consequently all local $B$-section
lifts of the target pair are conjugate in $B$.

The target local section has trivial centralizer in $A$.
Indeed a centralizer element has abelian image fixed by
the twisted Frobenius, whose action on $V$ has no
eigenvalue one. It therefore lies in $\bigwedge^2V$;
the Frobenius eigenvalues there also have absolute value
strictly greater than one, so that central element is
zero. The local section centralizer in $B$ is likewise
trivial: any such element centralizes its nonzero cusp
inertia and so lies in that cusp group, by B.4 or
(A.19.13); Frobenius multiplies this group by $q$, with
no nonzero fixed element. Hence every local pair lifting
the exact target pair is uniquely of the form
\[
       \widetilde\beta|_{D_q}^{\,j}\quad(j\in J_B),
\]
where the superscript means geometric conjugation by
$j$, with cocycle convention
$\beta^j(g)=j^{-1}\beta(g)g(j)$.
The conjugator belongs to $J_B$, since its projection
centralizes the target local section.

Extend each such local pair to a global normalized
cocycle by A.15. Every element of the global fiber is
obtained this way. A conjugation between two resulting
global normalized cocycles belongs to $B^c$ and its
projection centralizes $\beta$; the latter centralizer
is already trivial after restriction at $q$. Thus the
conjugator lies in $J_B^c$. Composition of conjugations
is $(\widetilde\beta^{,j})^h
=\widetilde\beta^{,jh}$, so the identifications are
exactly right multiplication by $J_B^c$. This proves
(A.24.3). For the map to the metabelian quotient the
same argument uses its trivial ramified local section
centralizer and kernel $N''$, giving the last assertion.
\end{proof}

\begin{statement}{A.25}{Proved}{The relative arithmetic presentation at every inert prime $q\equiv2,5\pmod9$}
Put
\[
 \Xi=\{q\text{ prime}:q\equiv2\text{ or }5\pmod9\}.
 \tag{A.25.1}
\]
This set has natural density $1/3$ and contains $2$.
For $q\in\Xi$ let $\mathfrak q=(q)$ be the unique
prime of $K_0$ above $q$, and write
$P_q=G_{K_0,\{\lambda,\mathfrak q,\infty\}}(3)$.
After conjugating the local embedding compatibly with
the fixed global complex conjugation, there is an
isomorphism
\[
 P_q\simeq D_{q^2}*_{3}\mathbb Z_3^{\langle z\rangle},
 \qquad
 D_{q^2}=\langle\phi,\tau\mid
                  \phi\tau\phi^{-1}=\tau^{q^2}\rangle_3,
 \tag{A.25.2}
\]
and its relative arithmetic quotient is
\[
 \Gamma_q\simeq
 (D_{q^2}*_{3}\mathbb Z_3^{\langle z\rangle})
                   \rtimes\langle c\rangle,
 \qquad
 c\phi c^{-1}=\phi,\quad
 c\tau c^{-1}=\tau^{-1},\quad czc^{-1}=z^{-1}.
 \tag{A.25.3}
\]
Here $c^2=1$. Thus $\sigma=zc$ is another involution,
with no further relation beyond (A.25.3). The source
category in these presentations consists of profinite
groups whose kernel over $\langle c\rangle$ is pro-three;
the kernel free product is precisely the pro-three free
product, as displayed.

For an inert prime $q\equiv8\pmod9$, the two unit
Frobenius values used in this construction are both zero
modulo three. In particular $q=17$ is the smallest prime
not covered by $\Sigma\cup\Xi$; (A.25.2) is not asserted
there. The union $\Sigma\cup\Xi$ has natural density
$5/9$.
\par\smallskip\noindent\emph{Manuscript comparison: \S51.16,
folios 733v--736v, scans 172--178, formulas (371)--(379).}
\end{statement}
\begin{proof}
The three S-unit Kummer characters over $K_0$ are
$\kappa_\zeta,\kappa_\lambda,\kappa_q$.
At the inert prime the residue field is $\mathbb F_{q^2}$.
Write $a=(q^2-1)/3\bmod3$. The two unramified values
are
\[
 (\kappa_\zeta(\phi),\kappa_\lambda(\phi))=(a,2a).
 \tag{A.25.4}
\]
Indeed $\lambda^{q-1}=-\zeta^2$ in this residue field,
because Frobenius sends $\zeta$ to $\zeta^2$ and
$\bar\lambda=-\zeta^2\lambda$. Consequently
$\lambda^{(q^2-1)/3}=\zeta^{2(q+1)/3}$.
The sign is one for odd $q$, since $(q+1)/3$ is even,
and is one in characteristic two as well. Since
$q-1=1$ modulo three, this exponent is $2a$.
The value $a$ is nonzero exactly for
$q\equiv2,5\pmod9$. The character $\kappa_q$ gives
the independent tame-inertia valuation.

Thus restriction in degree one from $P_q$ to the local
group $D_{q^2}$ has rank two, with a one-dimensional
kernel. In degree two the argument of A.14 applies
with just two finite places in $S$: the Brauer
invariants at $\lambda$ and $\mathfrak q$ sum to zero,
so restriction to $\mathfrak q$ is injective. Inflation
from the maximal pro-three quotient is injective by
(A.14.4). Choose one element complementing the local
image in the global Frattini quotient. The resulting
map $D_{q^2}*_{3}\mathbb Z_3\to P_q$ is an
isomorphism by the five-term kernel proof of A.14.
In particular the local morphism is injective before
any assertion about an equivariant complement is made.

The relative local group over $\mathbb Q_q$ is
$D_{q^2}\rtimes C_2$. It is explicitly obtained from
the tower of unramified extensions with degrees
$2\cdot3^n$ and compatible three-power roots of the
uniformizer $q$. Its order-two unramified element
acts on tame three-inertia by inversion and commutes
with the pro-three unramified generator $\phi$;
the latter acts on inertia by $q^2$. Its morphism
to the global relative quotient is injective, since
it is injective on the pro-three kernel just proved
and on the quotient $C_2$. Its order-two complement
is conjugate to the fixed global $c$ by an element
of $P_q$, using the coprime complement argument of
A.15. Conjugate the local embedding accordingly.
This yields the first two action formulas in
(A.25.3) with the actual chosen $c$.

It remains to produce the equivariant free generator.
On the three global Kummer characters, conjugation is
\[
 c(\kappa_\zeta)=\kappa_\zeta,
 \quad c(\kappa_\lambda)=\kappa_\zeta-\kappa_\lambda,
 \quad c(\kappa_q)=-\kappa_q.
 \tag{A.25.5}
\]
This uses conjugation on the fixed coefficient
$\mu_3$ as well as on the S-units. The Frattini
quotient therefore has signs $(1,2)$, whereas the
local subgroup has signs $(1,1)$ by its explicit
$\phi,\tau$ action. Over $\mathbb F_3$, division
by two gives an invariant complement, and that
one-dimensional complement has minus sign.
Choose an element $x\in P_q$ lifting a nonzero
vector in it and put
\[
                         z=x\,c(x)^{-1}.
\]
Then $c(z)=z^{-1}$ exactly, while its Frattini
class is twice that chosen vector. It and the
local Frattini image span the whole quotient.
The map with this new generator in place of the
initial complement is again an isomorphism by
the same degree-one and degree-two argument.
This proves both (A.25.2) and the full equivariant
presentation (A.25.3). In terms of $\sigma=zc$,
the last relation is exactly $\sigma^2=1$;
substitution $z=\sigma c$ gives the reverse
morphism of presentations, proving that no
additional relation was omitted.

Chebotarev in $\mathbb Q(\zeta_9)/\mathbb Q$
assigns density $1/6$ to each invertible residue
class modulo nine, so (A.25.1) has density $1/3$.
It consists of inert primes, whereas $\Sigma$
consists of split primes. Their densities add
to $1/3+2/9=5/9$. Formula (A.25.4) gives zero
for both units at $q\equiv8\pmod9$; thus the
local degree-one restriction at $17$ has rank
one, not two. This identifies the precise failure
of the present free-product construction there.
The exact residue-field computations, involution
matrices, and all primes through $1000$ are
reproduced in \texttt{code/a\_inert\_prime.py}
and its output. The list has $58$ primes in
$\Xi$ in that range.
\end{proof}

\begin{statement}{A.26}{Proved}{The complete one-prime metabelian image and fibers for $q\in\Xi$}
For every $q\in\Xi$ and $S_q=\{3,q,\infty\}$,
the conclusions (A.24.1)--(A.24.2) hold. Thus every
metabelian section supported at these primes lifts
to $N$, and for $B=N$ or $N/N''$ the exact image
in class two is
\[
 \operatorname{im}\bigl(\operatorname{Sec}_{S_q}(B)
                  \to\operatorname{Sec}_{S_q}(A)\bigr)
       =\operatorname{res}_q^{-1}(\mathscr L_q).
 \tag{A.26.1}
\]
For targets unramified at $q$, the fiber is exactly
the unramified fiber of A.20. For a ramified target
in this image, choose a normalized lift and let
$X_0$ be its cocycle value at the involution
$\sigma=zc$ of A.25. On $J_B$ put
\[
                     \Theta=\operatorname{Ad}(X_0)\circ\sigma.
\]
This is an involution. The fiber is in bijection
with either of the following explicitly identified
sets:
\[
 \{j\in J_B:j\Theta(j)=1\}
       \simeq J_B^{\Theta}\backslash J_B,
 \qquad h\longmapsto h^{-1}\Theta(h).
 \tag{A.26.2}
\]
The quotient in the right member is a set of left
cosets, without an assertion of normality. The same
ramified-fiber formula for the surjection from full
to metabelian sections uses $J_B=N''$ and the
involution supplied by the chosen full lift.
\par\smallskip\noindent\emph{Manuscript comparison: \S51.16,
folios 733v--736v, scans 172--178, formulas (371)--(379);
Tasks A.6 and B.4, including the prime two.}
\end{statement}
\begin{proof}
Normalize the cocycle at $c$ to one. Presentation
(A.25.3), with $z$ replaced by $\sigma c$, gives
the following exact data for a section with values
in a geometric quotient $B$:
\[
 \begin{gathered}
 a\in B^c,\qquad c(b)=b^{-1},\qquad
 aF_{q^2}(b)a^{-1}=b^{q^2},\\
 X\in B,\qquad X\sigma(X)=1.
 \end{gathered}
 \tag{A.26.3}
\]
The first row is its section on the relative local
group at $q$, with values at $\phi,\tau$. The
second is its value at the additional involution.
The remaining conjugations are by $B^c$, acting
on the first row as in (A.15.2) and on the
second row by $X\mapsto h^{-1}X\sigma(h)$.
The pro-three universal property in (A.25.2)
proves that these data define continuous sections,
and the given presentation proves that they
exhaust them.

For every equivariant pro-three surjection
$R\twoheadrightarrow B$, the second-row data
lift. Indeed $H^1(\langle\sigma\rangle,B)=1$
by the coprime argument of A.15, so write
$X=h^{-1}\sigma(h)$, lift $h$ to $\widetilde h$
and take $\widetilde h^{-1}\sigma(\widetilde h)$.
This is a lift satisfying the exact involution
relation. A lift of the local section in the
first row can be made to have value one on $c$
without changing its projection: its value at
$c$ is a cocycle in the pro-three kernel, and
the same coprime argument conjugates it to one.
The two rows can therefore be lifted independently
and combined using (A.26.3).

A.19 supplies the first-row lift to $N$ for every
metabelian section. For a class-two section B.7--
B.8 and A.19 say that this first row lifts exactly
when its local class belongs to $\mathscr L_q$.
This proves the asserted surjectivity and
(A.26.1). If the target is unramified at $q$,
all its full and metabelian lifts are unramified
there by their zero-inertia results. The same
inertia-kernel argument as in A.24 therefore
identifies the fiber with A.20.

Now fix a ramified target in the image and a
normalized lift. Its local lifting class is
unique by the injectivity of the ramified cusp
packet in class two. As in A.24, its local
geometric centralizer is trivial both in the
target and in $B$. Thus local lifts of the
exact target with value one at $c$ form one
free conjugacy orbit under $J_B^c$.
Use that conjugacy to fix the entire local
first row to the chosen lift. No residual
geometric conjugation remains, because the
fixed local section has trivial centralizer.
The remaining global fiber is consequently
the literal set of second-row lifts of the
target value at $\sigma$.

Write a second-row lift uniquely as $jX_0$,
$j\in J_B$. Its involution equation is
\[
 jX_0\sigma(j)\sigma(X_0)=1
       \quad\Longleftrightarrow\quad j\Theta(j)=1.
\]
The identity $X_0\sigma(X_0)=1$ proves that
$\Theta^2=1$. Coprime $H^1$ for this action
on $J_B$ gives every solution as
$h^{-1}\Theta(h)$. Two choices of $h$ give
the same solution precisely when they differ
by left multiplication by $J_B^{\Theta}$.
This proves (A.26.2), including its coset
orientation. The same argument with target
$N/N''$ and kernel $N''$ gives the last
assertion.
\end{proof}

\begin{statement}{A.27}{Proved}{The normalized integral Frobenius at an inert prime and the explicit Heisenberg test at seventeen}
Retain the free pro-three basis $u,v=cuc$ of $P_3$ from
A.10, and write $[u,v]=uvu^{-1}v^{-1}$. For an inert
rational prime $q\ne3$, conjugate its local order-two
complement in $P_3\rtimes\langle c\rangle$ to $\langle c\rangle$.
Let $\varphi_q\in P_3$ be the resulting arithmetic Frobenius
of the prime of $K_0$ of norm $q^2$. Its image in the
integral class-two quotient is exactly
\[
 \varphi_q\bmod\gamma_3P_3
       =u^{a_q}v^{a_q}[u,v]^{-a_q^2/2},
 \qquad a_q=\frac{\log(-q)}{\log7}\in\mathbb Z_3.
 \tag{A.27.1}
\]
The image in (A.27.1) is independent of the conjugator
used to identify the local order-two complement with
$\langle c\rangle$.

Put $\theta=\zeta_9$, $\beta^3=\lambda$, and
\[
 u_i=1-\zeta^i\theta\quad(0\le i\le2),\qquad
 \omega=u_1u_2^2,
 \qquad
 L=K_0(\theta,\beta,\omega^{1/3}).
 \tag{A.27.2}
\]
Then $L/K_0$ is the full exponent-three class-two
quotient of $P_3$: its Galois group has order $27$ and
presentation
\[
 \langle\sigma,\tau,\rho\mid
  \sigma^3=\tau^3=\rho^3=1,\quad
  \rho\text{ central},\quad[\sigma,\tau]=\rho\rangle.
 \tag{A.27.3}
\]
The field is Galois over $\mathbb Q$, and complex
conjugation acts by
\[
 c\sigma c=\tau\sigma,
 \qquad c\tau c=\tau^{-1},
 \qquad c\rho c=\rho^{-1}.
 \tag{A.27.4}
\]
For every prime $q\equiv8\pmod9$, the Frobenius of
the prime of $K_0$ above $q$ is trivial in (A.27.3).
In particular this holds at $17$, although
$v_3(a_{17})=1$ in the integral formula (A.27.1).
For example,
\[
 a_{17}\equiv78\pmod{81},\qquad
 -a_{17}^2/2\equiv36\pmod{81}.
 \tag{A.27.5}
\]
\par\smallskip\noindent\emph{Manuscript comparison: \S51.16,
folios 733v--736v, scans 172--178, formulas (371)--(379);
the first inert prime outside A.25 in Tasks A.6 and B.4.}
\end{statement}
\begin{proof}
The image of the local unramified group at $q$ in the
relative quotient over $K_0$ is procyclic with quotient
$C_2$ and pro-three kernel. Its order-two subgroup is
the prime-to-three component of this procyclic group.
All complements to $P_3$ in $P_3\rtimes C_2$ are
$P_3$-conjugate, by the coprime argument of A.15.
After the specified conjugation, $\varphi_q$ commutes
with $c$.

There is an exact integral logarithm description
of the free class-two pro-three group on $u,v$:
its underlying lattice has basis $U,V_0,[U,V_0]$,
and multiplication is
$X*Y=X+Y+\frac12[X,Y]$. Here
$U=\log u$ and $V_0=\log v$. The fraction $1/2$
belongs to $\mathbb Z_3$; multiplication and the
inverse maps are integral polynomial maps, so this
identifies the pro-three group and the lattice with
their compact topologies. The involution interchanges
$U,V_0$ and negates $[U,V_0]$. Thus its fixed
subgroup is exactly
$\{\exp(a(U+V_0)):a\in\mathbb Z_3\}$.
The cyclotomic character sends both $u,v$ to $7$
and sends $\varphi_q$ to $q^2$. Since the logarithm
is injective on $1+3\mathbb Z_3$, it follows that
$2a\log7=2\log(-q)$, giving $a=a_q$.
The product formula in this lattice yields
(A.27.1). Two allowed conjugators differ by an
element of $P_3^c$; its image belongs to the displayed
procyclic fixed subgroup and therefore commutes with
the image of $\varphi_q$. This proves the asserted
independence, with its exact ambiguity identified.

For the field in (A.27.2), direct multiplication gives
\[
 u_0u_1u_2=\lambda,
 \qquad
 \frac{\sigma(\omega)}{\omega}=
       \frac{u_0^3}{\lambda},
 \qquad c(\omega)=-\omega,
 \tag{A.27.6}
\]
where initially $\sigma(\theta)=\zeta\theta$.
The classes of $\zeta$ and $\lambda$ are independent
in $K_0^\times/K_0^{\times3}$: the valuation of
$\lambda$ at the prime $(\lambda)$ is one, whereas
$\zeta$ is a unit and is not a cube in $K_0$.
Consequently $K_0(\theta,\beta)/K_0$ has degree nine
and $\lambda$ is not a cube in $K_0(\theta)$.
The kernel of
\[
 K_0(\theta)^\times/K_0(\theta)^{\times3}
       \longrightarrow
 K_0(\theta,\beta)^\times/
                       K_0(\theta,\beta)^{\times3}
\]
is the line generated by $\lambda$. To see this
directly, a cube root in the cyclic cubic extension
has an eigenvector under its Galois group and hence
is a scalar multiple of $1,\beta$, or $\beta^2$.
If the class of $\omega$ belonged to that line,
the class of $\sigma(\omega)/\omega$ would be
trivial. This contradicts (A.27.6). Thus adjoining
$\gamma=\omega^{1/3}$ has degree three.

Choose the cube roots $\beta,\gamma$ so that actual
complex conjugation satisfies
$c(\beta)=-\theta^2\beta$ and $c(\gamma)=-\gamma$.
Each choice is possible by multiplying an initial
cube root by a power of $\zeta$; the ratios before
and after that multiplication differ by the
corresponding power of $\zeta$. The following
explicit automorphisms are then well-defined:
\[
\begin{array}{c|ccc}
 &\theta&\beta&\gamma\\\hline
 \sigma&\zeta\theta&\beta&u_0\gamma/\beta\\
 \tau&\theta&\zeta\beta&\gamma\\
 \rho&\theta&\beta&\zeta\gamma\\
 c&\theta^{-1}&-\theta^2\beta&-\gamma
\end{array}
 \tag{A.27.7}
\]
Their powers and products, using (A.27.6), prove
(A.27.3)--(A.27.4). All the radicands are units
away from three, and the degrees are powers of
three. The Kummer polynomials and their derivatives
therefore show that the extension is unramified
at every finite prime away from three.
It is consequently a quotient of $P_3$.
The free exponent-three class-two group on two
generators has the normal forms
$\sigma^a\tau^b\rho^z$, with $a,b,z\in\mathbb F_3$,
and order $27$. The displayed quotient has this
order, so the quotient map identifies its kernel
exactly with the closed subgroup
$P_3^3\gamma_3P_3$.

For a rational prime $q\equiv-1\pmod9$, the
rational Frobenius in $\operatorname{Gal}(L/\mathbb Q)$
acts on $\theta$ as $\theta\mapsto\theta^{-1}$.
It consequently has the form $g c$ with
$g=\tau^b\rho^z$. Formula (A.27.4) gives
$(gc)^2=g c(g)=1$. Its square is the Frobenius
over $K_0$, proving its triviality in (A.27.3).
Equivalently, at every residue choice of $\theta$
in $\mathbb F_{q^2}$, (A.27.6) gives
$\omega^q=-\omega$, whence
\[
 \omega^{(q^2-1)/3}
       =(-1)^{(q+1)/3}=1.
\]
The exponent on the right is even because the
odd prime $q$ is $-1$ modulo nine.

The program \texttt{code/a\_frobenius17.py}
checks (A.27.6) in the exact ring
$\mathbb Z[\theta]=\mathbb Z[T]/(T^6+T^3+1)$,
all $729$ products under (A.27.4), and all three
choices of $\theta$ at each of the $29$ primes
$q\equiv8\pmod9$ through $1000$. At seventeen
one choice gives $\theta=9+11\zeta$ and
$\omega=3+6\zeta$, with cube roots
$2+\zeta$, $16+\zeta$, and $16+15\zeta$.
The same program computes $a_{17}$ by unique
successive lifting of $7^a=-17$, without a
floating-point logarithm. It gives (A.27.5)
and the output through modulus $3^{12}$.
Finally $v_3(\log(-17))=v_3(-18)=2$ and
$v_3(\log7)=1$, proving $v_3(a_{17})=1$.
\end{proof}

\begin{statement}{A.28}{Proved}{The exact equivariant one-relator presentation at every inert prime}
Let $q\ne3$ be any inert rational prime, and
let $P_q=G_{K_0,\{3,q\}}(3)$ and
$\Gamma_q=P_q\rtimes\langle c\rangle$.
After the local-complement identification used
in A.27, choose a tame generator $t$ and an
arithmetic Frobenius $\phi$ at the prime over
$q$, with
\[
 c(t)=t^{-1},\qquad c(\phi)=\phi,
 \qquad \phi t\phi^{-1}=t^{q^2}.
 \tag{A.28.1}
\]
Choose a lift $x\in P_q$ of $u\in P_3$ and
put $y=cxc$. The lift can and will be chosen
with $\kappa_q(x)=0$ modulo three. Let
$F=F_3(x,y,t)$ be the free pro-three group
with involution $c(x)=y$, $c(y)=x$,
$c(t)=t^{-1}$. Its natural equivariant
map to $P_q$ is surjective. There is a
$c$-fixed pro-three word $W\in F$ mapping
to the actual local Frobenius $\phi$,
and for every such choice the natural map
gives the exact presentations
\[
 \begin{split}
 P_q&\simeq
 F/\overline{\langle\!\langle
          WtW^{-1}t^{-q^2}\rangle\!\rangle},\\
 \Gamma_q&\simeq
 \bigl(F/\overline{\langle\!\langle
          WtW^{-1}t^{-q^2}\rangle\!\rangle}\bigr)
                         \rtimes\langle c\rangle.
 \end{split}
 \tag{A.28.2}
\]
These are continuous pro-three presentations,
with the indicated finite operator group in
the second line. In particular
\[
 P_q^{\mathrm{ab}}\simeq
 \mathbb Z_3^2\oplus\mathbb Z_3/(q^2-1),
 \qquad
 \dim_{\mathbb F_3}H^2(P_q,\mathbb F_3)=1.
 \tag{A.28.3}
\]
At $q=17$, $W$ lies in $\Phi(F)$, and the
relator in (A.28.2) lies in the third
Zassenhaus subgroup of $F$. Its image on
killing $t$ is the actual Frobenius of A.27;
in particular its integral class-two image
is given by (A.27.1) and (A.27.5).

For any of the geometric quotients $B$ under
consideration, (A.28.2) gives the following
exact section parameterization. Normalize
the value at $c$ to one; assign the values
$a,c(a),b$ at $x,y,t$. If
$\Omega_W(a,b)$ denotes the value at $W$
of this continuous free-group cocycle,
the remaining equations are precisely
\[
 c(b)=b^{-1},\qquad
 \Omega_W(a,b)F_{q^2}(b)\Omega_W(a,b)^{-1}
                                      =b^{q^2}.
 \tag{A.28.4}
\]
Classes are the orbits under $B^c$ acting by
\[
 (a,b)\longmapsto
       (h^{-1}a\,x(h),\ h^{-1}bh).
 \tag{A.28.5}
\]
\par\smallskip\noindent\emph{Manuscript comparison: \S51.16,
folios 733v--736v, scans 172--178, formulas (371)--(379);
Tasks A.6 and B.4 beyond the prime families A.16 and A.25.}
\end{statement}
\begin{proof}
The Kummer group for $K_0$ with ramification
allowed at three and $q$ has the basis
$\zeta,\lambda,q$. The first two classes
are pulled back from $P_3$, and the last
has value one on $t$. Adjusting an initial
lift of $u$ by a power of $t$ gives the
specified value $\kappa_q(x)=0$; because
$c(\kappa_q)=-\kappa_q$, the value at
$y$ is also zero. Thus $x,y,t$ form a
basis of the Frattini quotient, and the
continuous map $F\to P_q$ is onto.
The involution is equivariant by its
definition on these three generators.

The local generators (A.28.1) can be
chosen in the relative tame local group:
the unramified $\mathbb Z_3$ generator
and its order-two component commute,
and the latter acts on three-primary
inertia as the Teichm\"uller scalar $-1$.
Coprime complement conjugacy identifies
this order-two element with the given
global $c$. The local pro-three group
maps injectively into $P_q$. For an
explicit realization, the extension
\[
 K_0(\mu_{3^\infty},q^{1/3^\infty})/K_0
 \tag{A.28.6}
\]
is pro-three and unramified outside three
and $q$. At $q$ its cyclotomic part
contains every unramified extension of
three-power degree: $q^2\in1+3\mathbb Z_3$
has infinite order, so the local
cyclotomic degrees are unbounded. The
uniformizer roots supply every tame
three-power ramification degree. The
standard tame description then identifies
its local compositum with the maximal
pro-three extension of $K_{0,q}$.
Consequently the local decomposition
map into $P_q$ has trivial kernel.

Since $\phi$ is $c$-fixed and $F\to P_q$
is an equivariant pro-three surjection,
the coprime $H^1$ argument of A.15 proves
$F^c\to P_q^c$ is onto. This supplies
$W\in F^c$ as stated. Put
$r=WtW^{-1}t^{-q^2}$ and $G_W=F/(r)$,
where $(r)$ denotes the closed normal
subgroup it generates. As
\[
 c(r)=Wt^{-1}W^{-1}t^{q^2}
        =t^{-q^2}r^{-1}t^{q^2},
\]
this subgroup is stable under $c$.
There is a surjective equivariant map
$G_W\to P_q$, inducing an isomorphism
on $H^1(-,\mathbb F_3)$. Here
$r\in\Phi(F)$ since $q^2\equiv1\pmod3$.

The degree-two comparison is the exact
one proved in A.14: full-to-maximal-pro-three
inflation is injective, and the full
global group has
\[
 H^2(G_{K_0,\{3,q\}},\mathbb F_3)
 \simeq
 \{(z_3,z_q)\in(\tfrac13\mathbb Z/\mathbb Z)^2:
                         z_3+z_q=0\}.
 \tag{A.28.7}
\]
The Kummer, Picard, and Brauer comparisons
are precisely Milne II.2.9, II.2.1, and
II.2.2(a) \cite{MilneADT}, applied to the
localization of the same principal ideal
domain. Thus restriction
$H^2(P_q,\mathbb F_3)\to
H^2(D_{q^2},\mathbb F_3)$ is injective.
The assignments $\phi\mapsto W$,
$t\mapsto t$ define a continuous map
$D_{q^2}\to G_W$, whose composite to
$P_q$ is the actual local map. It follows
that inflation
$H^2(P_q,\mathbb F_3)\to
H^2(G_W,\mathbb F_3)$ is injective.
The five-term kernel argument proved
in A.14 now shows that $G_W\to P_q$
is an isomorphism. This proves (A.28.2),
including the stated operator action.

Abelianizing its single relator gives
exactly $(1-q^2)t$ and no relation on
$x,y$, proving the first formula of
(A.28.3). This abelianization has
nonzero finite torsion; $P_q$ is therefore
not free pro-three. The minimal relation
criterion used in A.10 gives
$H^2(P_q,\mathbb F_3)\ne0$, while its
injection into $H^2(D_{q^2},\mathbb F_3)$
bounds the dimension by one. The latter
local group has dimension one, as follows
from the Hochschild--Serre sequence for
$\mathbb Z_3(1)\rtimes\mathbb Z_3$ with
trivial $\mathbb F_3$ coefficients.
This proves the second formula.

At seventeen the Frobenius values on
$\zeta,\lambda$ are zero by (A.25.4).
The chosen local Frobenius fixes the
uniformizer roots, so its value on
$q$ is also zero modulo three. Conjugating
the local embedding does not change
these degree-one values. Therefore
$W\in\Phi(F)$. If $D_n(F)$ denotes the
Zassenhaus filtration, then
$\Phi(F)=D_2(F)$, so $[W,t]\in D_3(F)$;
also $t^{q^2-1}=t^{288}\in D_9(F)$.
The relator thus belongs to $D_3(F)$.
The map killing $t$ identifies the
quotient with the unramified group
$P_3$ and sends $W$ to the normalized
Frobenius already computed in A.27.

Finally a normalized cocycle on
$F\rtimes C_2$ is exactly the arbitrary
value $a$ on $x$, its conjugate on $y$,
and the value $b$ on $t$ satisfying
$c(b)=b^{-1}$. The continuous value at
the pro-three word $W$ is defined by
the free pro-three universal property;
this defines $\Omega_W$ without replacing
$W$ by a truncated word. Arithmetic
inertia at $q\ne3$ acts trivially on
the original geometric group, while
$W$ acts as its actual Frobenius
$F_{q^2}$. Evaluating $r$ gives exactly
the second equation (A.28.4), and
conjugating a normalized cocycle gives
(A.28.5). These computations prove both
directions of the parameterization.
The word $W$ in this assertion retains
its possible dependence on $t$;
(A.27.1) computes its stated quotient,
not that full dependence. Solving
(A.28.4) in the full geometric group
for every prescribed metabelian target
at seventeen remains the specific
unresolved step recorded in the ledger.
\end{proof}

\begin{statement}{A.29}{Proved}{Local cuspidality at every finite derived level and its global one-prime consequences}
Write $N^{(0)}=N$ and
$N^{(m+1)}=\overline{[N^{(m)},N^{(m)}]}$,
and put $Q_m=N/N^{(m)}$ for $m\ge2$.
For every finite extension $k/\mathbb Q_q$,
$q\ne3$, every continuous section with
geometric quotient $Q_m$ is either unramified
or lies in exactly one actual cusp packet.
Its inertia is trivial exactly when its
abelian inertia is zero. Every such local
section lifts to $N$, and the map from full
to $Q_m$ sections is bijective on the
ramified classes. Thus the exact local
description in A.19 holds with $Q_m$ in
place of $N/N''$, including the actual
central Frobenius coordinates of its image
in class two.

For every $q\in\Sigma\cup\Xi$ and every
$m\ge2$, the global map
\[
 \operatorname{Sec}_{\{3,q,\infty\}}(N)
       \longrightarrow
 \operatorname{Sec}_{\{3,q,\infty\}}(Q_m)
 \tag{A.29.1}
\]
is onto. For ramified targets its fiber
is the right-coset set
$N^{(m)}/(N^{(m)})^c$ when $q\in\Sigma$,
and the left-coset set
$(N^{(m)})^\Theta\backslash N^{(m)}$
when $q\in\Xi$, with the chosen-lift
involution $\Theta$ of A.26. Unramified
fibers are the exact unramified fibers
of A.20 for the map $N\to Q_m$.
\par\smallskip\noindent\emph{Manuscript comparison: \S51.16,
folios 733v--736v, scans 172--178, formulas (371)--(379);
the successive derived approximations in Tasks A.6 and B.4.}
\end{statement}
\begin{proof}
We first identify the precise peripheral
intersection statement used below.
Each boundary generator $c_d$ extends
to a free pro-three basis of $N$:
in the standard surface presentation,
eliminate a different boundary generator.
Yamaguchi's Proposition 1.1.6
\cite{AYamaguchi2023} gives the centralizer
for every nonzero integer power. For
$a=3^s u$ with $u\in\mathbb Z_3^\times$,
continuous exponentiation gives
$c_d^{3^s}=(c_d^a)^{u^{-1}}$ and
$c_d^a=(c_d^{3^s})^u$. The two elements
generate the same closed procyclic
subgroup and hence have the same
centralizer. We therefore obtain
\[
 C_{Q_j}(c_d^a)=\langle c_d\rangle
       \quad(j\ge2,\ a\in\mathbb Z_3\setminus\{0\}).
 \tag{A.29.2}
\]
If $g c_d^a g^{-1}=c_e^b$ with $a,b\ne0$,
abelianization gives $a e_d=b e_e$.
The distinct boundary lines are disjoint
away from zero, and each $e_d$ is primitive.
It follows that $d=e$ and $a=b$.
Formula (A.29.2) then forces
$g\in\langle c_d\rangle$. This proves
the required intersection and normalizer
claims through explicit maps to $V$ and
the stated centralizer theorem.

We induct on $m$, starting with A.19.
Let a local section in $Q_m$ have tame
value $t$, and denote its image in
$Q_{m-1}$ by $\bar t$. By induction
$\bar t$ is either one or conjugate
to $c_d^a$, $a\ne0$. In the latter case
choose such a conjugator and lift it
through $Q_m\to Q_{m-1}$; this changes
the representative by an allowed
geometric conjugation and makes
$\bar t=c_d^a$ exactly.

Choose a decreasing cofinal sequence
of characteristic open normal subgroups
$M_n$ of $Q_{m-1}$, and let $H_n\subset N$
be the inverse image of
$\langle\bar t\rangle M_n$. Then
\[
 N^{(m-1)}\subset H_n,
 \qquad N^{(m)}\subset H_n',
 \qquad
 H_\infty:=\bigcap_nH_n
      =\begin{cases}
         N^{(m-1)},&\bar t=1,\\
         N^{(m-1)}\langle c_d^a\rangle,&\bar t=c_d^a.
        \end{cases}
 \tag{A.29.3}
\]
These $H_n$ need not be normal in $N$;
the finite covers and their boundary
points below are described by cosets
and double cosets, respectively.
The characteristic choice and the
section's Frobenius relation make
$H_n/N^{(m)}$ stable under the section.
Pulling back
$(H_n/N^{(m)})s(G_k)$ from
$Q_m\rtimes G_k$ gives an actual
arithmetic cover with geometric group
$H_n$. It contains original arithmetic
inertia because $t\in H_n/N^{(m)}$.
Thus this cover has the good tame
compactification used in A.17--A.19.
The inclusion $N^{(m)}\subset H_n'$
identifies its actual geometric
abelianization with that of
$H_n/N^{(m)}$. In particular the
chosen lift of $t$ has the exact
Frobenius eigenvalue $\#\mathbb F_k$
in $H_n^{\mathrm{ab}}$, rather than
only an eigenvalue in a further quotient.
The compactification weight argument
places it in the integral boundary
submodule at every $n$.

As proved in A.19 by stabilization in
every finite quotient and compact
lifting of cosets,
\[
 H_\infty^{\mathrm{ab}}
       \simeq\varprojlim_n H_n^{\mathrm{ab}}.
 \tag{A.29.4}
\]
We compute the inverse limit of the
boundary modules with their actual
transition maps. The boundary points
above $e$ are
$H_n\backslash N/I_e$; the loop at
$H_ngI_e$ is the appropriate conjugate
of $c_e^{j_n(g,e)}$, where
\[
 j_n(g,e)=[I_e:I_e\cap g^{-1}H_ng].
 \tag{A.29.5}
\]
Transition maps sum over inverse images
and multiply each term by the ratio
of these ramification indices. The
diagonal relation has transition map
equal to the covering degree.
The index $[N:H_n]$ is unbounded,
because the group $Q_{m-1}$ cannot
have an open procyclic subgroup:
its quotient $V$ has rank ten.
The diagonal tower consequently has
zero limit, and its first derived
limit is zero by the compact equation
argument in A.19. The inverse limit
of boundary modules is therefore the
inverse limit of the weighted free
boundary modules in (A.29.5).

If $\bar t=1$, every intersection
$I_e\cap g^{-1}H_\infty g$ is trivial,
since each cusp group injects into
$Q_{m-1}$ by its primitive abelian image.
Every ramification index tends to
infinity. This growth is uniform on
each clopen set of the compact double
coset space: otherwise nested nonempty
closed sets with bounded indices would
give a point with nontrivial limiting
intersection. The weighted transition
maps are thus divisible by arbitrary
powers of three at each fixed finite
stage. Their inverse limit is zero.
By (A.29.4), $t$ is zero in
$H_\infty^{\mathrm{ab}}
=N^{(m-1)}/N^{(m)}$, proving $t=1$.

Now let $\bar t=c_d^a$ and write
$a=3^s\varepsilon$ with
$\varepsilon\in\mathbb Z_3^\times$.
A bounded index in (A.29.5) is
equivalent to a nontrivial intersection
of $g I_e g^{-1}$ with
$\langle c_d^a\rangle$ in $Q_{m-1}$.
By (A.29.2), this forces $e=d$ and
the image of $g$ to belong to $I_d$.
There is therefore exactly one
surviving boundary double coset
in $H_\infty\backslash N/I_d$.
Its ramification index stabilizes at
$3^s$. Its $3^s$ sheets under
$I_d/\langle c_d^a\rangle$ constitute
one boundary point, not $3^s$
independent boundary generators.
On every clopen set away from this
point the same compactness argument
gives uniform unbounded ramification.
The inverse boundary module is
consequently the single copy of
$\mathbb Z_3(1)$ generated by the
class of $c_d^{3^s}$. Projecting
$H_\infty$ to $\langle c_d^a\rangle$
shows that the class of $t$ is
$\varepsilon[c_d^{3^s}]=[c_d^a]$
in (A.29.4).

Put $A_m=N^{(m-1)}/N^{(m)}$ and
let $T$ denote conjugation by $c_d^a$.
The image of $H_\infty$ in $Q_m$ is
$A_m\rtimes\langle c_d^a\rangle$.
Its closed commutator subgroup is
exactly $(T-1)A_m$: the image is
closed because $A_m$ is compact.
The equality of abelianizations just
proved gives
$t(c_d^a)^{-1}\in(T-1)A_m$, which
is equivalent to
$t=h^{-1}c_d^a h$ for some $h\in A_m$.
This proves actual cusp conjugacy.
After making this conjugation,
compare the Frobenius value with
that of the actual tangential cusp
section. Their ratio centralizes
$c_d^a$, hence lies in $I_d$ by
(A.29.2). The complete local section
therefore lies in the actual cusp
packet. Its parameters and uniqueness
are exactly the Kummer parameters of
A.19 and B.8, and its full lift is
the corresponding cusp section.
In the unramified case an arbitrary
Frobenius value lifts to $N$, as
proved in A.19. This completes
the induction and the local statement.

For the global consequences, the
arithmetic presentations A.15--A.16
and A.25 are independent of the
geometric quotient. Replace the
local metabelian lifting input in
A.24 and A.26 by the local result
just proved. The local ramified
lift is still unique and has trivial
full local geometric centralizer;
the independent involution datum
in A.26 still lifts by coprime
$H^1$. The identical explicit maps
then prove (A.29.1) and its stated
fibers with kernel $N^{(m)}$.
For an unramified target, the local
zero-inertia result forces all full
lifts to be unramified at $q$, so
the comparison with A.20 is exact.
\end{proof}

\begin{statement}{A.30}{Proved}{The mixed Heisenberg test at seventeen and the exact first relation degree}
Use the presentation at $q=17$ in A.28.
For its Frobenius word and defining relator,
the stronger exact conclusions are
\[
 W\in D_3(F)\setminus D_4(F),\qquad
 r=WtW^{-1}t^{-289}
           \in D_4(F)\setminus D_5(F).
 \tag{A.30.1}
\]
Here $D_n(F)$ is the Zassenhaus filtration,
equivalently the filtration by degree in
the mod-three Magnus expansion of the
free pro-three group $F=F_3(x,y,t)$.
If $x,y,t$ have Magnus variables $X,Y,T$,
the degree-four part of $r-1$ has
coefficient $2$ on both $X^3T$ and
$Y^3T$. These coefficients are
independent of the other degree-three
coefficients of $W-1$.

The arithmetic input detecting the
only possible degree-two component
of $W$ is the following explicit
Heisenberg extension. Let
\[
 \alpha^3=3,\qquad \beta^3=17,
 \qquad u_0=2+\alpha^2,
 \qquad u_i=\sigma^i(u_0),
 \quad \sigma(\alpha)=\zeta\alpha,
\]
and set
\[
 \omega=u_1u_2^2,
 \qquad L_{3,17}=K_0(\alpha,\beta,\omega^{1/3}).
 \tag{A.30.2}
\]
This field is Galois over $K_0$, has
the group (A.27.3), and is unramified
outside three and seventeen. At the
prime with $\alpha\equiv7\pmod{17}$,
the Frobenius fixing the uniformizer
roots has trivial image in the center
of this Heisenberg quotient.
\par\smallskip\noindent\emph{Manuscript comparison: \S51.16,
folios 733v--736v, scans 172--178, formulas (371)--(379);
the next arithmetic relation at the first remaining inert prime
in Tasks A.6 and B.4.}
\end{statement}
\begin{proof}
The norm calculation is exact:
\[
 u_0u_1u_2=2^3+\alpha^6=17,
 \qquad
 \frac{\sigma(\omega)}{\omega}
          =\frac{u_0^3}{17}.
 \tag{A.30.3}
\]
The classes of $3$ and $17$ in
$K_0^\times/K_0^{\times3}$ are
independent, by their valuations at
the primes above three and seventeen.
The degree and automorphism argument
of A.27 applies with these two
radicands and the norm witness
$u_0$. In particular the last
cubic extension in (A.30.2) is
nontrivial: otherwise the class
of $\omega$ over $K_0(\alpha)$
would lie on the line generated
by $17$, contradicting (A.30.3).
For $\gamma^3=\omega$, the explicit
automorphisms are
\[
\begin{array}{c|ccc}
 &\alpha&\beta&\gamma\\\hline
 \sigma&\zeta\alpha&\beta&u_0\gamma/\beta\\
 \tau&\alpha&\zeta\beta&\gamma\\
 \rho&\alpha&\beta&\zeta\gamma
\end{array}
 \tag{A.30.4}
\]
and $[\sigma,\tau]=\rho$.
Choose positive real cube roots
$\alpha,\beta$ in the fixed complex
embedding. Complex conjugation
interchanges $u_1,u_2$ and satisfies
\[
 c(\omega)\omega=(17/u_0)^3,
 \qquad c(\gamma)=17/(u_0\gamma).
\]
The second equality follows from
$|\gamma|^2=17/u_0>0$, for any
choice of the cube root $\gamma$.
Thus the field is Galois over
$\mathbb Q$, with
\[
 c\sigma c=\sigma^{-1},\qquad
 c\tau c=\tau^{-1},\qquad
 c\rho c=\rho.
 \tag{A.30.5}
\]
The radicands are units away from
three and seventeen, so the
Kummer polynomials are unramified
at every other finite prime.

At the chosen prime over
seventeen, the residue field of
$K_0$ is $\mathbb F_{17^2}$ and
$\alpha=7$ is a cube root of
$3$. In the basis $1,\zeta$,
direct finite-field multiplication
gives
\[
 (u_0,u_1,u_2)=(0,\,4+2\zeta,\,2+15\zeta),
 \qquad \omega=7+10\zeta,
 \qquad \omega^{96}=1.
 \tag{A.30.6}
\]
Here $\omega$ is a unit, so its
cubic residue is exactly the
Frobenius value on $\gamma$.
The specified local Frobenius
fixes $\alpha$ and the uniformizer
root $\beta$, and (A.30.6) shows
that it fixes $\gamma$ as well.
Its Heisenberg image is therefore
the identity. In particular this
value is unchanged by conjugating
the local embedding to align its
order-two complement with $c$.

We now compare this arithmetic
calculation with the actual free
presentation. By A.28, $W\in D_2(F)$
and $c(W)=W$. The quotient
$D_2(F)/D_3(F)$ has basis
\[
 [x,y],\quad[x,t],\quad[y,t]
 \quad\text{over }\mathbb F_3.
\]
Indeed commutators generate this
quotient, cubes already have
Magnus degree three, and these
three degree-two commutator
polynomials are linearly independent.
The involution negates the first
basis vector and sends the last
two to each other's negatives.
Its fixed subspace is exactly
the line generated by
$[x,t]-[y,t]$.

Under the quotient (A.30.2), the
Frattini images of $x,y,t$ on
the two Kummer coordinates
$(3,17)$ are
\[
 x\longmapsto(1,0),\qquad
 y\longmapsto(2,0),\qquad
 t\longmapsto(0,1).
 \tag{A.30.7}
\]
For the first coordinate this
uses $\kappa_3=-\kappa_\zeta-
\kappa_\lambda$ and the two
Frobenius rows of A.10. The
second uses the lift choice
in A.28. Consequently the
fixed wedge generator maps
to $(1-2)\rho=-\rho\ne0$.
The trivial Frobenius value
in (A.30.6) forces the class
of $W$ on that line to vanish.
This proves $W\in D_3(F)$.

In $F^{\mathrm{ab}}$, its
coordinates on $x,y$ are
$a_{17},a_{17}$ by A.27,
and its coordinate on $t$
is zero because it is fixed
by the involution which
inverts $t$. We have
$a_{17}\equiv6\pmod9$.
Setting $Y=T=0$ in its
Magnus expansion therefore
gives $(1+X)^{a_{17}}$.
The coefficient of $X^3$ is
\[
 \binom{a_{17}}3
       \equiv a_{17}/3\equiv2\pmod3.
\]
The same holds for $Y^3$.
Thus $W\notin D_4(F)$.
If $W_3$ is its homogeneous
degree-three term, the first
possible term of $[W,t]-1$
is $W_3T-TW_3$. The coefficient
of $X^3T$ in this expression
is exactly $2$, and likewise
for $Y^3T$; no other monomial
of $W_3$ can contribute to
either coefficient. The factor
$t^{1-289}=t^{-288}$ has first
possible Magnus degree nine,
since $288=9\cdot32$.
It follows that $r$ has exact
degree four, proving (A.30.1).

The exact ring identities,
all three residue choices,
and the nonzero Magnus
coefficients with arbitrary
remaining cubic terms are
reproduced by
\texttt{code/a\_frobenius17\_mixed.py}
and its output. This calculation
determines the first relation
degree; the remaining coefficients
of $W$ and the full lifting
equation (A.28.4) continue to
be retained in the unresolved
global problem.
\end{proof}

\begin{statement}{A.31}{Proved}{The nonzero cubic Frobenius term at seventeen and its failure to be a cube}
Let $L_0=K_0(\alpha)$ with $\alpha^3=3$, and
let $H=G_{L_0,\{3\}}(3)\subset P_3$.
Put
\[
 \pi=\frac{\lambda}{\alpha},\qquad
 \eta=\alpha^2-2,\qquad
 \delta=1+\zeta^2\alpha+\zeta\pi.
 \tag{A.31.1}
\]
The four Kummer classes of
$\zeta,\pi,\eta,\delta$ form an
$\mathbb F_3$-basis of $H^1(H,\mathbb F_3)$.
At every prime of $L_0$ above seventeen,
the Frobenius row on this basis is
\[
                         (0,0,0,1).
 \tag{A.31.2}
\]

Choose $m\in P_3^c$ whose abelian
coordinates in the basis $u,v$ are
$(1,1)$, and put $n=vu^{-1}$.
Then $m,n$ are a free pro-three
basis, $c(m)=m$, $c(n)=n^{-1}$,
and the normalized Frobenius
of A.27 satisfies the exact congruence
\[
 \varphi_{17}\equiv
       m^6[n,[n,m]]\pmod{D_4(P_3)}.
 \tag{A.31.3}
\]
Moreover $\varphi_{17}$ is not a cube
of any element of $P_3$.
\par\smallskip\noindent\emph{Manuscript comparison: \S51.16,
folios 733v--736v, scans 172--178, formulas (371)--(379);
the explicit nonprimitive Frobenius in Tasks A.6 and B.4.}
\end{statement}
\begin{proof}
The subgroup $H$ is exactly
$\ker(\kappa_3:P_3\to\mathbb F_3)$.
Indeed every finite Galois three-extension
of $L_0$ unramified outside three has
a Galois closure over $K_0$ with
three-group Galois group: take the
compositum of its three conjugates,
whose kernel over $L_0$ embeds in
a product of three-groups. This
proves the stated comparison with
the open subgroup of $P_3$.
The subgroup is free of rank four
by the pro-three Schreier formula.
Thus its first cohomology has
dimension four.

The element $\pi$ is a unit away
from three. For $\sigma(\alpha)=\zeta\alpha$,
direct multiplication gives
\[
 N_{L_0/K_0}(\eta)=1,
 \qquad
 \delta^3=\zeta\eta^2\sigma(\eta),
 \qquad
 N_{L_0/K_0}(\delta)=\zeta.
 \tag{A.31.4}
\]
Consequently $\eta$ and $\delta$
are units at every finite prime.
The following four Frobenius
evaluations prove independence
of their Kummer classes:
\[
\begin{array}{c|c|c|rrrr}
 q&\bar\zeta&\bar\alpha&
 \kappa_\zeta&\kappa_\pi&\kappa_\eta&\kappa_\delta\\\hline
 2&\zeta&1&1&2&0&2\\
 2&\zeta&\zeta&1&1&2&2\\
 2&\zeta&\zeta^2&1&0&1&0\\
 61&13&4&2&0&1&0
\end{array}
 \tag{A.31.5}
\]
The first three residue fields are
$\mathbb F_4$ with
$\zeta^2+\zeta+1=0$; the last
is $\mathbb F_{61}$. In each row,
raise the unit residue to
$(\#\mathbb F-1)/3$ and express
the answer as a power of the
displayed residue of $\zeta$.
The integer determinant of the
four-by-four matrix in (A.31.5)
is $2$, hence is nonzero modulo
three. The dimension calculation
then proves the basis assertion.

At seventeen choose the residue
field $\mathbb F_{17^2}$ and first
take $\alpha=7$. The four residues
are
\[
 \zeta,\qquad5+12\zeta,\qquad
 13,\qquad16+3\zeta.
\]
Their $96$th powers give
$(1,1,1,\zeta)$. The other two
choices $\alpha=7\zeta,7\zeta^2$
give the same four Kummer values.
This proves (A.31.2), including
the effect of any conjugation
permuting the three primes above
seventeen.

Coprime $H^1$ for the kernel
of $P_3\to P_3^{\mathrm{ab}}$
makes $P_3^c\to(P_3^{\mathrm{ab}})^c$
onto. This gives $m$ with the
specified abelian vector. Together
with $n=vu^{-1}$ its Frattini
vectors have determinant $2$,
so $m,n$ are a free basis.
The actions of $c$ are exactly
those stated, and
\[
 \kappa_3(m)=0,\quad
 \kappa_3(n)=1,\quad
 \kappa_\zeta(m)=1
       \quad\text{in }\mathbb F_3,
 \tag{A.31.6}
\]
by the two rows of A.10.
Thus conjugation by $n$ acts
on $L_0$ as $\sigma$.
One completely specified choice
of $m$ uses only continuous
pro-three square roots. Put
$d_0=uc(uv)^{-1/2}$ and
$h_0=(d_0c)^{1/2}\in P_3$, and
take $m=h_0^{-1}(uv)h_0$.
Indeed $(uc)^2=uv$, so $d_0$
is an involution commuting
with $uv$. The identity
$c(d_0c)c=(d_0c)^{-1}$ gives
$h_0ch_0^{-1}=d_0$ and hence
$c(m)=m$. Its abelian vector
is still $(1,1)$. Square roots
here are the unique continuous
powers $g^{1/2}$ in the closed
procyclic group generated by
$g$; this defines every word
and comparison in the formula.

We give the finite-group
comparison which permits use
of (A.31.2) in degree three.
A Schreier basis of $H$ is
$n^3,m,nmn^{-1},n^2mn^{-2}$.
Modulo $\Phi(H)$ these form
a basis of an elementary
abelian group of order $3^4$.
Conjugation by $n$ fixes
$n^3$ and cyclically permutes
the other three basis vectors.
Hence $P_3/\Phi(H)$ has
nilpotence class at most three,
exponent dividing nine, and
commutator subgroup of exponent
three. The defining product
formula
\[
 D_4(P_3)=P_3^9\,\gamma_2(P_3)^3\,
                         \gamma_4(P_3)
\]
therefore proves
$D_4(P_3)\subset\Phi(H)$.
Every Kummer character in
(A.31.5), in particular
$\ell=\kappa_\delta$, annihilates
this subgroup.

The vector space
$D_3(P_3)/D_4(P_3)$ has basis
\[
 m^3,\quad n^3,\quad
 [m,[m,n]],\quad[n,[n,m]].
 \tag{A.31.7}
\]
Collection gives these
generators; their Magnus
polynomials are respectively
$M^3$, $N^3$,
$M^2N+MNM+NM^2$, and
$N^2M+NMN+MN^2$, which are
independent. On this basis
$c$ has signs $+,-,-,+$.
The Frobenius lies in $D_3$
by A.27 and is fixed by $c$,
so its class has form
$A m^3+B[n,[n,m]]$.
Its abelian coordinates in
the $m,n$ basis are
$(a_{17},0)$, whence
$A=a_{17}/3=2$ modulo three.

For $h\in H$, Kummer theory
identifies
$\ell(nhn^{-1})$ with
$\kappa_{\sigma^{-1}\delta}(h)$.
Since $-2=1$ in $\mathbb F_3$,
we obtain
\[
 \begin{split}
 \ell([n,[n,m]])
   &=((\sigma^{-1}-1)^2\kappa_\delta)(m)\\
   &=\kappa_{N_{L_0/K_0}\delta}(m)
     =\kappa_\zeta(m)=1.
 \end{split}
 \tag{A.31.8}
\]
Also $\ell(m^3)=0$. Formula
(A.31.2) gives
$\ell(\varphi_{17})=1$,
so $B=1$. This proves the
full degree-three congruence
(A.31.3), rather than only
its cyclotomic component.

Finally, a cube root of
$\varphi_{17}$ in $P_3$
would have abelian vector
$(a_{17}/3,a_{17}/3)$ in
the $u,v$ basis. Its
$\kappa_3$ value is zero,
so it would belong to $H$.
The character $\kappa_\delta$
annihilates every cube in
$H$, contradicting (A.31.2).
This proves the last assertion.

The exact unit identities
and root are verified by
\texttt{a\_frobenius17\_unit\_root.py}
in \texttt{code/};
the residue matrix and all
three Frobenius rows at
seventeen are verified by
\texttt{a\_frobenius17\_sunits.py}
in the same directory.
Their outputs are included.
An initial attempted list
$\zeta,\pi,\eta,\sigma(\eta)$
had rank three, not four;
the exact cubic root
$\delta$ in (A.31.4)
supplies the missing unit.
No basis assertion relies
on that rejected list.
\end{proof}

\begin{statement}{A.32}{Proved}{The exact marked global image and ramified fibers on $\Sigma\cup\Xi$}
Put
$R=N/\gamma_3(I_{\mathrm{mark}})$,
$A=N/\gamma_3N$, and
$J_R=\ker(R\to A)$.
For every $q\in\Sigma\cup\Xi$,
the three global class-two images
coincide and are exactly
\[
 \begin{split}
 \operatorname{im}\bigl(\operatorname{Sec}_{S_q}(N)
                                   \to\operatorname{Sec}_{S_q}(A)\bigr)
 &=\operatorname{im}\bigl(\operatorname{Sec}_{S_q}(N/N'')
                                   \to\operatorname{Sec}_{S_q}(A)\bigr)\\
 &=\operatorname{im}\bigl(\operatorname{Sec}_{S_q}(R)
                                   \to\operatorname{Sec}_{S_q}(A)\bigr)
   =\operatorname{res}_q^{-1}(\mathscr L_q),
 \end{split}
 \tag{A.32.1}
\]
where $S_q=\{3,q,\infty\}$ and
$\mathscr L_q$ is the actual local
unramified and cusp-packet set of A.24.

Fix a ramified class-two target in
this image and let $d\in D(\mathbb Q_q)$
be its unique local cusp. Define its
local marked lifting-class set by
\[
 \mathcal L_{q,d}=
 \begin{cases}
  \{*\},&d\ne Q_c,\\
  \mathcal C_{\mathbb Q_q,Q_c},&d=Q_c,
 \end{cases}
 \tag{A.32.2}
\]
using exactly the reference and
quotient of B.12--B.13.
Restriction gives a surjection from
the global marked fiber onto
$\mathcal L_{q,d}$.
For each element of this local
class set, choose a normalized
global marked lift above it.
The inverse image of that element
is precisely
\[
 \begin{cases}
 J_R/J_R^c,&q\in\Sigma,\\
 J_R^\Theta\backslash J_R,&q\in\Xi,
 \end{cases}
 \tag{A.32.3}
\]
with right and left cosets,
respectively. In the second case
$\Theta=\operatorname{Ad}(X_0)\circ\sigma$
is the actual involution specified
by the chosen lift in A.26.
Thus (A.32.2)--(A.32.3) describe
the entire fiber without identifying
different local marked classes or
asserting a canonical product
decomposition. Targets unramified
at $q$ have exactly the unramified
fiber of A.20 for $R\to A$.

The exact image of
$\operatorname{Sec}_{S_q}(N)\to
\operatorname{Sec}_{S_q}(R)$ consists
of the classes unramified at $q$,
the unmarked ramified cusp classes,
and the marked ramified classes
whose local parameter in
$\mathcal C_{\mathbb Q_q,Q_c}$ is zero.
Its nonempty ramified fibers are
given by (A.32.3) with
$J_R$ replaced by
$\gamma_3(I_{\mathrm{mark}})$ and
the corresponding chosen full lift.
In particular there are global
marked sections supported at
$\{3,q\}$ which have no full lift,
although their class-two image
has a full lift. These examples
still have no full lift after
any finite number-field extension.
\par\smallskip\noindent\emph{Manuscript comparison: \S51.16,
folios 733v--736v, scans 172--178, formulas (371)--(379);
the marked higher image and fibers in Task A.6.}
\end{statement}
\begin{proof}
The local marked classification
B.12--B.14 proves that its exact
class-two image is $\mathscr L_q$.
The local fiber over an unmarked
ramified cusp is a singleton;
the local fiber over the marked
cusp is the full group
$\mathcal C_{\mathbb Q_q,Q_c}$.
Each class in the latter group
is an actual local section,
with the Frobenius representative
$c_{Q_c}^r zF_0$ and inertia
$c_{Q_c}^a$ of (B.12.3).
Its class-two image retains
$a,r$ and kills $z$ through
the actual morphism
$\gamma_2I_{\mathrm{mark}}/
 \gamma_3I_{\mathrm{mark}}\to A$.

For $q\in\Sigma$, a normalized
global section is its arbitrary
local pair modulo $R^c$, by
A.15--A.16. Any prescribed local
marked lifting class therefore
globalizes. To make its projection
the exact local pair of the
given global class-two target,
lift the required local geometric
conjugator through $R\to A$.
There is no further arithmetic
relation. For $q\in\Xi$, first
normalize the local $c$-value
to one using coprime $H^1$
in the pro-three kernel.
The independent involution
value in the second row of
(A.26.3) lifts the prescribed
class-two value by the same
coprime argument. It combines
with the selected local row
to give the required global
marked section. This proves
surjectivity onto (A.32.2)
and the marked equality in
(A.32.1). The other two
equalities are A.24 and A.26.

We check the centralizers used
to pass from local to global
conjugation. The class-two local
section has trivial geometric
centralizer: its image in $V$
would be Frobenius-fixed, hence
zero by the weight calculation,
and the remaining element in
$\bigwedge^2V$ is zero for the
same reason. At an unmarked
ramified cusp in $R$, its
projection to the filled full
group has trivial centralizer
by B.7, and the kernel has
no inertia-fixed element in
either $M$ or $Z$ by the
free label actions in B.14.
Hence its centralizer in
$R$ is zero. At the marked
cusp the centralizer of the
inertia element is
$\langle c_{Q_c}\rangle Z$
by (B.12.6). A full-section
centralizer in this group
has zero cusp coordinate
because multiplication by
$q-1$ is injective on
$\mathbb Z_3$; its central
coordinate is also zero,
because $q^2\bigwedge^2\Phi-1$
is injective on $Z$.
This last injectivity follows
on every finite free permutation
quotient from its eigenvalues,
and then in the inverse limit.
It is unaffected by the
central Frobenius parameter
$z$. Thus every ramified
local marked section used
here has trivial geometric
centralizer.

Fix a local marked lifting
class. Exact local lifts of
the target pair form one
free conjugacy orbit under
$J_R$. Indeed their conjugator
projects to a centralizer of
the target, which has just
been proved trivial. For
$\Sigma$, the remaining
global conjugators are
$J_R^c$. With the cocycle
convention of A.24, their
action on the local orbit
parameter is right multiplication,
giving $J_R/J_R^c$.
For $\Xi$, normalized local
rows form a free orbit under
$J_R^c$. Fix one such row
by this conjugation; no
geometric conjugator remains.
The remaining data are exactly
the lifts of the independent
involution value. As in
(A.26.2), writing them as
$jX_0$ gives $j\Theta(j)=1$
and the bijection
$J_R^\Theta\backslash J_R
\to\{j:j\Theta(j)=1\}$.
This proves (A.32.3) with
its coset orientations.

The same presentation arguments
apply to $N\to R$. B.7 and
B.12--B.14 show that a local
marked class has a full lift
exactly for the parameters
listed in the statement.
The lifting class is unique
when ramified. Thus the local
condition is necessary, the
arithmetic presentations make
it sufficient, and the same
centralizer and coset argument
gives the asserted fibers.
If a class-two or marked
target is unramified at $q$,
the zero-abelian-inertia results
of B.7 and B.12 force every
higher lift here to be
unramified there as well.
The comparison with A.20
is consequently exact.

For existence of the last
examples, start with any
nonzero marked tame exponent
and choose a nonzero element
of $\mathcal C_{\mathbb Q_q,Q_c}$;
B.13 gives such elements,
including an explicit central
commutator detected in a
finite label quotient. The
preceding construction gives
the global marked section,
while its class-two image
belongs to (A.32.1).
The group $\mathcal C_{k,Q_c}$
is the unramified first
cohomology of the actual
central coefficient module
$Z$ in B.12. On this module
the cusp inertia acts trivially.
Its tame first-cohomology
component has Frobenius
$Q\bigwedge^2\Phi$, with
no invariant vector, so
it is the whole
$H^1_{\mathrm{cts}}(G_k,Z)$.
For any finite local extension,
restriction followed by
corestriction on this continuous
cohomology is multiplication
by its degree; this is the
usual finite coset-sum chain
map and is valid for compact
coefficients. B.13 proves
that the source group is
torsion-free. Its restriction
is therefore injective.
A nonzero central parameter
remains nonzero at every
place of every finite
number-field extension above
$q$, and the local full
image still has central
parameter zero by B.12.
This proves persistence of
the failure of full lifting.
\end{proof}

\begin{statement}{A.33}{Proved}{A class-three lift for every finite-support locally cuspidal class-two section}
Let $S'$ be any finite set of rational primes containing three,
with the real place included. Put $A=N/\gamma_3N$ and
$M_3=\gamma_3N/\gamma_4N$. Every section of $A$ over
$G_{\mathbb Q,S'}$ whose restriction at each
$q\in S'\setminus\{3\}$ belongs to the actual local image
$\mathscr L_q$ of B.8 has a section lift to $N/\gamma_4N$
over the same arithmetic group. Its lifting fiber is a torsor under
\[
 H^1(G_{\mathbb Q,S'},M_3)
 =H^1(G_{\mathbb Q,\{3,\infty\}},M_3)
 =\mathbb Z_3^{166}\oplus(\mathbb Z/3)^7
   \oplus(\mathbb Z/9)^{56}
   \oplus(\mathbb Z/27)^{25}
   \oplus(\mathbb Z/81)^3.
 \tag{A.33.1}
\]
These assertions concern the actual common class-three quotient
of the full and metabelian groups.

There is also an exact obstruction to choosing this global lift
with prescribed actual cuspidal local class-three restrictions.
Let $J\subset S'\setminus\{3\}$ be the primes where the
class-two target is ramified. Choose any global class-three lift
$s_3$. At each $q\in J$ the unique actual cuspidal class-three
lifting class determines a difference
$\Delta_q(s_3)\in M_3/(F_q-1)M_3$. The element
\[
 \left[(\Delta_q(s_3))_{q\in J}\right]
 \in\operatorname{coker}\left(
 H^1(G_{\mathbb Q,\{3,\infty\}},M_3)
 \xrightarrow{\operatorname{res}}
 \bigoplus_{q\in J}M_3/(F_q-1)M_3\right)
 \tag{A.33.2}
\]
is independent of $s_3$. It vanishes exactly when a global
class-three lift can have all the indicated cuspidal restrictions.
Every full or metabelian global lift of the class-two target
forces (A.33.2) to vanish.
\end{statement}
\begin{proof}
The action on the central kernel $M_3$ is the actual action
induced from $V$; conjugation by the class-two section does not
change it. Its finite quotients have three-group image over
$K_0$. The composition-series argument in A.1 therefore gives
\[
 H^i(G_{\mathbb Q,\{3,\infty\}},M_3/3^n)=0
 \quad(i\ge2).
\]
For $X=\operatorname{Spec}\mathbb Z[1/3]$, apply the actual
localization triangle (A.1.2) with $M_3/3^n$ and remove the
finite set $S'\setminus\{3\}$. Its degree-two term injects
into the sum of the residue groups
$H^1(\mathbb F_q,M_3(-1)/3^n)$. These are precisely the
local groups $H^2(G_{\mathbb Q_q},M_3/3^n)$, by the tame
inertia calculation in A.1. Thus the actual global lifting
obstruction, restricted to each local arithmetic group, is its
corresponding local lifting obstruction. An actual local full
lift exists by membership in $\mathscr L_q$, so all these
restrictions are zero.

The finite-$S'$ degree-one groups with finite coefficients are
finite. Their first derived limit consequently vanishes.
Taking inverse limits therefore gives an injection
\[
 H^2_{\rm cts}(G_{\mathbb Q,S'},M_3)
 \longrightarrow\bigoplus_{q\in S'\setminus\{3\}}
       H^2_{\rm cts}(G_{\mathbb Q_q},M_3),
 \tag{A.33.3}
\]
and the continuous obstruction is zero. This proves existence
of a continuous global class-three lift, without a passage
from unrelated finite-level choices.

Every rational Frobenius eigenvalue on $M_3(-1)$ has complex
absolute value strictly greater than one. Indeed $M_3$ is
the degree-three multiplication kernel of A.22, and every
eigenvalue is a product of three eigenvalues on $V$; their
absolute values are at least $q^{3/2}$ before the twist.
Hence $H^0(\mathbb F_q,M_3(-1))=0$. Degree-one localization
and the same finite-$S'$ inverse-limit argument identify
the two first-cohomology groups in (A.33.1). The Smith
calculation of A.22 proves the displayed decomposition.
The target section has trivial geometric centralizer: its
image in $V$ is Frobenius-fixed and therefore zero, and
its remaining image in $\bigwedge^2V$ is zero for the
same reason. The ordinary central-extension lifting action
is consequently a torsor on section classes, with no further
stabilizer quotient.

Locally the same weight calculation kills the tame term in
$H^1(G_{\mathbb Q_q},M_3)$, so this group is exactly
$M_3/(F_q-1)M_3$. A ramified class-two cusp packet fixes
its actual full cuspidal lifting class, by B.7--B.8, and
therefore fixes its class-three image. The difference from
$s_3$ is consequently an element of this first-cohomology
group. Replacing $s_3$ by a global central twist adds the
restriction of that twist to every local difference. This
proves the independence and the asserted equivalence in
(A.33.2). A.19 forces every ramified local metabelian
restriction to be the actual cusp packet, and B.7 gives
the same conclusion for the full group. Their class-three
images therefore make every local difference zero, proving
the final necessity assertion.
\end{proof}

\begin{statement}{A.34}{Proved}{The restriction map at seventeen and a class-three lift that stops at class four}
Use the actual arithmetic generators $m,n,c$ of A.31.
On the degree-three lattice $M_3$, the restriction map
\[
 H^1(G_{\mathbb Q,\{3,\infty\}},M_3)
 \longrightarrow H^1(G_{\mathbb Q_{17}},M_3)
                         =M_3/(F_{17}-1)M_3
 \tag{A.34.1}
\]
has image of dimension $15$ after reduction in the local
target modulo three. The target modulo three has dimension
$164$, and the cokernel of (A.34.1) modulo three therefore
has dimension $149$.

Fix free geometric generators $x_0,\ldots,x_9$ whose
abelianizations are the integral basis of A.7, with the
first eight entries the non-$O$ cusp generators. In $M_3$
put $e=[[x_8,x_0],x_9]$. The exact normalized cocycle
\[
 \beta(c)=0,\qquad \beta(m)=(1+c)e,
                  \qquad\beta(n)=0
 \tag{A.34.2}
\]
is unramified outside three. Start with the actual full
$O$-cuspidal section whose tangent parameter is $17$,
project it to $N/\gamma_4N$, and twist by $\beta$.
The resulting global section $s_3$ has the same class-two
image as that full cuspidal section, but $s_3$ has no lift
to $N/\gamma_5N$. In particular it has no full or
metabelian lift, and no lift to the marked quotient with
marked cusp $Q_c$.

Thus a class-two target with a full lift can have a
class-three lifting class with a nonzero next obstruction.
The theorem does not assert non-lifting of that class-two
target: its original full cuspidal lift is specified above.
\end{statement}
\begin{proof}
All the following matrices act on the reduction $\overline M_3$.
Write $U,V_0,C$ for the actual matrices of $u,v,c$ in A.22,
and put
\[
 M=UV_0,\qquad L=V_0U^{-1},\qquad
 A_0=M-1,\qquad B_0=L-1,\qquad
 \Pi_+=(1+C)/2.
 \tag{A.34.3}
\]
The action over $K_0$ is abelian by the explicit isogeny
and CM comparisons in A.7. Thus these really are the
actions of $m,n$, independently of the chosen $c$-fixed
lift $m$. The exact matrices give
\[
 A_0^3=B_0^3=A_0^2B_0^2=0,
 \quad A_0B_0=B_0A_0,
 \quad CMC=M,
 \quad CLC=L^{-1}.
 \tag{A.34.4}
\]
In particular the fourth power of the augmentation ideal
of this coefficient action is zero modulo three.
Normalized cocycles are prescribed freely by
\[
       a\in\ker(C-1),\qquad b\in\ker(LC+1)
 \tag{A.34.5}
\]
as their values at $m,n$. These subspaces have dimensions
$164,166$. They are reductions of the corresponding
integral eigensublattices, since the two operators $C$
and $LC$ are involutions and two is invertible.
The exact free presentation with $c(m)=m$ and
$c(n)=n^{-1}$ proves that every integral such pair is
an actual normalized cocycle.

For any cocycle, evaluation on $D_5(P_3)$ is zero modulo
three. Evaluation on $\gamma_5(P_3)$ contains four
augmentation operators and is zero by (A.34.4).
The coefficient action kills all commutators, so evaluation
on cubes of commutators is three times a value and is
zero. Evaluation on ninth powers is also zero: the action
of every element has order dividing three, and its
nine-term norm is three times its three-term norm.
These are the factors generating $D_5(P_3)$.
Consequently (A.31.3) leaves only the three possible
degree-four Hall terms
\[
 K_1=[m,[m,[m,n]]],\quad
 K_2=[n,[m,[m,n]]],\quad
 K_3=[n,[n,[m,n]]].
\]
Their cocycle values, and the value of the specified
degree-three term, are respectively
\[
\begin{array}{c|l}
 m^6[n,[n,m]]&(2A_0^2+B_0^2)a-A_0B_0b\\
 K_1&-A_0^2B_0a+A_0^3b\\
 K_2&-A_0B_0^2a+A_0^2B_0b\\
 K_3&-B_0^3a+A_0B_0^2b.
\end{array}
 \tag{A.34.6}
\]
These formulas follow by repeated use of
$f(gh)=f(g)+g f(h)$; commuting coefficient actions give
$f([g,h])=(1-h)f(g)+(g-1)f(h)$.

The action of the normalized rational Frobenius at
seventeen is $C$ modulo three. Indeed its square is
$\varphi_{17}$, and (A.27.1) gives the latter action as
$M^{a_{17}}$ with $a_{17}\in3\mathbb Z_3$, whereas
$M^3=1$ modulo three. The rational Frobenius is exactly
$c\varphi_{17}^{1/2}$ in its decomposition group.
The local target modulo three is therefore
$\overline M_3/(C-1)\overline M_3\simeq\overline M_3^C$.
Its Frobenius value is half the $\Pi_+$ projection of
the value at $\varphi_{17}$.

The program \texttt{a\_frobenius17\_eval\_m3.py} in the
\texttt{code} directory applies the four operators in
(A.34.6) to the two eigenspaces (A.34.5), using the full
integral basis and matrices in
\texttt{a\_degree3\_matrices.json}.
All operations here are exact in $\mathbb F_3$.
After projection by $\Pi_+$, the three degree-four
operators have ranks $0,1,0$. The specified
degree-three operator has rank $15$. Adding the middle
operator with coefficient $0$, $1$, or $2$ gives the
same column space, of rank $15$ in every case.
The saved image basis and the elimination routine
certify these comparisons, so unknown higher coefficients
of the actual Frobenius cannot change the image.
Coboundaries restrict to zero in the local quotient;
(A.34.5) supplies all integral normalized cocycles.
This proves the rank and cokernel assertions for the
actual map (A.34.1).

For (A.34.2), the same calculation has zero contribution
from every unknown degree-four term. In the multidegree
basis of A.22, $e$ has zero-based index $87$, namely
$(\{0,8,9\},8)$. Let $z$ be the reduction modulo three
of its local Frobenius value. It is $C$-fixed. The program
records all its $36$ nonzero coordinates and forms the
actual Lie bracket $[z,e_O]$ in the degree-four free Lie
algebra by its injection into the free associative algebra.
Here
$e_O=-\sum_{j=0}^7e_j$ in the fixed abelian basis.
For reproducibility, a single nonzero coefficient is
\[
 [z,e_O]_{X_0X_0X_0X_2}=1\quad\text{in }\mathbb F_3.
 \tag{A.34.7}
\]
The complete polynomial has $1098$ nonzero words. It is
saved in \texttt{a\_frobenius17\_\allowbreak eval\_m3.json}.
The code converts each multidegree difference to its
integral Hall bracket $[[x_i,x_j],x_k]$, expands the
brackets, and checks directly with the actual matrix on
$V$ that
\[
                    C[z,e_O]=-[z,e_O].
 \tag{A.34.8}
\]
Thus this is a nonzero element of the actual degree-four
free Lie lattice modulo three, with its stated sign.

Locally choose the uniformizer Kummer roots for $17$
so that the reference cusp section has inertia $c_O$
and Frobenius $F_0$. The twist (A.34.2) changes no
inertia, since it is unramified there, and changes
Frobenius to $zF_0$ modulo $\gamma_4N$.
In a putative lift to $N/\gamma_5N$, its inertia must
be $c_Ok$ for some
$k\in L_4:=\gamma_4N/\gamma_5N$. An arbitrary choice
of the lifted Frobenius differs only by a central
$L_4$ value, which has no effect on its conjugation
of inertia. The tame relation therefore gives the
necessary equation
\[
                 (F_{17}-17)k=-17[z,e_O]
                         \quad\text{in }L_4.
 \tag{A.34.9}
\]
Modulo three its left side is $(C+1)k$, in the plus
eigenspace. By (A.34.7)--(A.34.8) its right side is a
nonzero minus eigenvector. This is impossible because
two is a unit modulo three. Conjugating a prospective
lifting class would not avoid the equation: lift its
class-three conjugator through the surjection
$N/\gamma_5N\to N/\gamma_4N$ to make its projection
the prescribed representative, and then apply the
same calculation.

This proves the local obstruction and therefore the
failure of a global class-four lift. A.19 would make
any local metabelian lift of this ramified class an
actual cusp section and hence give a full local lift,
contradicting (A.34.9). B.14 gives the identical
consequence for a marked lift because $O$ is an
unmarked cusp. These exact quotient comparisons prove
the remaining non-lifting assertions. The original
$O$-cusp section remains a full lift of the shared
class-two image, as asserted.
\end{proof}

\begin{statement}{A.35}{Proved}{A fifteen-dimensional toric family with zero first mod-three gluing class}
Let $B\in(\bigwedge^2W)^{c=-1}$. This eigensublattice has
rank fifteen. The exact prescriptions
\[
 \gamma_B(c)=0,\qquad \gamma_B(m)=0,
                  \qquad\gamma_B(n)=B
 \tag{A.35.1}
\]
define an unramified-outside-three class-two central cocycle.
Twist the class-two image of the global $O$-cuspidal section
with tangent parameter $17$ by $\gamma_B$.
Its restriction at seventeen is the same actual class-two
cusp section, including its Frobenius coordinate. This target
has a global class-three lift with unchanged local inertia.
For the exact gluing class (A.33.2), its image in the
cokernel of (A.34.1) modulo three is zero.
No assertion of integral vanishing of (A.33.2) is made here.
\end{statement}
\begin{proof}
Over $K_0$ the toric lattice is scalar cyclotomic. Thus
$m$ acts on $\bigwedge^2W$ by $49^2$, whereas $n$ acts
as one. Its complex-conjugation signs are $(13,15)$,
from the signs $(3,5)$ on $W$. The free presentation of
A.31 proves (A.35.1). At the normalized local group,
$\varphi_{17}$ is fixed by $c$. The cocycle value at this
element lies both in $\mathbb Z_3B$ and in the plus
eigenspace of $c$, so is zero. Its value on tame inertia
is zero because the cocycle is unramified. Its value on
rational Frobenius $c\varphi_{17}^{1/2}$ is also zero.
This proves the exact local class-two assertion.

Here is an actual global class-three lift. In A.28 first
choose $x$ with $\kappa_{17}(x)=0$ exactly: right
multiplication by a unique power of the tame generator
achieves this, and the correction is in the Frattini
subgroup when the initial value is zero modulo three.
The conjugate $y$, and the resulting generators $m,n$,
then also have zero Kummer value. Keep the tame section
value $c_O$. Choose $h\in\gamma_2N/\gamma_4N$ which is
a product of toric commutators with degree-two image
$-B/2$. Give $m$ section value one and $n$ section value
$h^{-1}(nc)(h)$ relative to the reference cusp section.
The latter value obeys the exact involution relation
because $(nc)^2=1$. Its degree-two image is $B$.
The fixed word $W$ in A.28 has zero degree-two difference
from the reference, by the same $c$-fixed argument used
for $\varphi_{17}$. Its remaining difference in class
three is central, and therefore does not change its
conjugation of $c_O$. The actual defining relation
$WtW^{-1}=t^{17^2}$ is satisfied. The universal property
of that proved presentation gives the claimed global
class-three lift.

We calculate its Frobenius difference modulo three,
retaining the actual path origins. For every cusp $d$,
let $z_d$ be its framed Abel--Jacobi point and $h_d=3z_d$
the torus point of A.9. Choose compatible division
points $z_{d,j}$ and toric division points $h_{d,j}$.
The compatible torsion points
$3z_{d,j}-h_{d,j}$ form $v_d\in V$, and their actual
proper image in $T/3$ is $d-O$. Thus the path cocycle
has the exact identity
\[
 3b_{g,d}=\kappa(h_d)(g)+g v_d-v_{gd}.
 \tag{A.35.2}
\]
The toric component of $v_d$ is constrained by the
actual torsion point $z_d-h_{d,1}$; it has not been
replaced by an arbitrary lift. Altering the chosen
$z$-division points changes $v_d$ by $3V$ and changes
$b$ by the corresponding actual coboundary. We may
therefore choose its proper component to be the
specified integral cusp label $t_d=(a_d,b_d)$ of A.7.
For $g$ over $K_0$ the proper component of (A.35.2) is
$(T_g-1)t_d/3$.

Put
\[
 H_m=2I_2,\qquad
 H_n=(T_n-I_2)/3
 =\frac17\begin{pmatrix}-4&-1\\1&-5\end{pmatrix}.
 \tag{A.35.3}
\]
These are the proper components of the path operators
for $m,n$, since $T_m=7I_2$. The upper-right blocks in
the actual integral action on $V$ are
\[
                  \Phi_m=14\Psi,\qquad
                  \Phi_n=-\Psi H_n.
 \tag{A.35.4}
\]
Modulo three all the diagonal $K_0$ actions on $V$
are one. On a toric commutator $[e_d,e_e]$, its
change in degree three is the derivation with cusp
values $-[H_g t_d,e_d]$, together with toric
terms depending on the full framed origins.
The subsequent augmentation kills every toric term;
on a term with one proper $T$ factor it replaces
that factor by $\Phi_g$.

More explicitly, use the abelian coefficient group
$\gamma_2N/\gamma_4N$, with the action twisted by the
reference cusp section. In the submodule generated
by toric commutators and their images, three successive
augmentation operators are zero modulo three. A first
augmentation produces degree-three terms with at most
one $T$ factor; the second produces only $W$ terms;
the third is zero. The $c$-normalization used above
changes the chosen $n$-value only by degree-three
terms with at most one $T$ factor, so these changes
also vanish after two augmentations. The tame
augmentation sends a toric commutator to its bracket
with $e_O$, a pure $W$ term, and its product in either
order with every other augmentation is zero here.

For noncommuting coefficient operators $A_0=m-1$ and
$B_0=n-1$, the leading crossed-derivation value at
$[n,[n,m]]$, when the $m$-value is zero and the
$n$-value is $B$, is
\[
                   (A_0B_0-2B_0A_0)B.
 \tag{A.35.5}
\]
This follows by differentiating brackets:
$D[X,Y]=X D(Y)-Y D(X)$, including the action
$[X,Y]$ on the source value. All the higher group
corrections contain three augmentations and are
zero. The term $m^6$ contributes zero. A.30 proves
$W\in D_3(F_3(m,n,t))$ for the actual three-generator
word, so there are no degree-two tame-letter terms.
Its degree-three terms involving the tame letter have
two augmentations and vanish by the identities just
proved; a restricted tame cube has zero tame source
value. The unknown terms in $D_4$ have three augmentations
and are also zero.
Thus (A.31.3) computes the whole Frobenius difference
in this calculation.

For each cusp, the proper-path contribution to
(A.35.5) is, up to the common cusp-bracket sign,
\[
 \Phi_mH_n-2\Phi_nH_m
     =14\Psi H_n+4\Psi H_n
     =18\Psi H_n\equiv0\pmod3.
 \tag{A.35.6}
\]
The framed toric terms were killed before this
calculation, so no unknown origin term remains.
The exact program \texttt{a\_frobenius17\_toric\_twists.py}
in the \texttt{code} directory expands all twenty-eight
toric commutators in the free associative algebra,
using (A.35.3)--(A.35.4), and verifies (A.35.6).
Twenty-seven have nonzero individual summands, but
their stated difference is always zero.
Hence the chosen class-three lift has zero local
Frobenius difference in $M_3/3$. Inertia was unchanged
by construction. Its class in the local quotient
and therefore the gluing cokernel modulo three is
zero, which proves the statement.
\end{proof}

\begin{statement}{A.38}{Proved}{An actual Frobenius basis with completely specified finite class-two action}
Let $F_2=\langle x,y\rangle$ be the geometric pro-three group of
$\mathbb P^1-\{0,1,\infty\}$, based at the unit tangent at zero,
and let
\[
 N=\ker(F_2\longrightarrow(\mathbb Z/3)^2),\qquad
 K=N^3\gamma_3N.
 \tag{A.38.1}
\]
Here $N^3$ is the closed subgroup generated by cubes. There is
an actual Frobenius $u$ at the prime $(7,\zeta-2)$ of $K_0$ such
that $u,v=cuc$ are a free basis of $P_3$, their integral action
on $V$ is exactly A.7, and their action on $N/K$ is the following
fully specified finite action. Complex conjugation is induced by
$x\mapsto x^{-1},y\mapsto y^{-1}$. The restriction of Frobenius
$u$ to $N/K$ equals the restriction of the following
automorphism of $F_2/K$:
\[
 x\longmapsto x^7,\qquad y\longmapsto f_0^{-1}y^7f_0.
 \tag{A.38.2}
\]
In the ordered Schreier basis
\[
\begin{split}
 &(yxy^{-1}x^{-1},\ y^2xy^{-2}x^{-1},\ y^3,\
 xyxy^{-1}x^{-2},\ xy^2xy^{-2}x^{-2},\ xy^3x^{-1},\\
 &\hspace{25mm}x^3,\ x^2yxy^{-1},\ x^2y^2xy^{-2},\
 x^2y^3x^{-2}),
\end{split}
 \tag{A.38.3}
\]
write $N/K$ as $\mathbb F_3^{10}\oplus\bigwedge^2\mathbb F_3^{10}$
with product $(a,z)(b,w)=(a+b,z+w+\tfrac12a\wedge b)$.
The word-class representative $f_0$ defining this finite
automorphism has coordinates
\[
 a_f=(0,0,1,1,2,0,2,0,1,2),
 \tag{A.38.4}
\]
and, in lexicographic order $i<j$, the forty-five central
coordinates are
\[
\begin{split}
 z_f={}&(2,1,2,1,1,2,0,1,2,\ 1,1,0,0,0,2,2,0,1,\\
       &1,1,1,1,1,0,1,1,0,0,1,2,0,0,1,0,1,1,\\
       &0,0,1,2,0,2,2,0,0).
\end{split}
 \tag{A.38.5}
\]
Thus (A.38.2)--(A.38.5) specify every coordinate of the images
of the ten geometric generators. They specify a representative
for the induced finite action; the actual Ihara word can differ
by the kernel of its map to that action. The arithmetic basis
is obtained from any initial
Frobenius basis by at most four explicitly described conjugations
below. The underlying Galois action and tangential origin are
retained; all cocycles and all arithmetic Frobenius words are
transported by the resulting actual source automorphism.
\end{statement}
\begin{proof}
We first specify the geometric comparison. The map from the
Fermat curve is
\[
 t=-z^3/x^3,
 \qquad (-z/x)^3=t,\qquad(-y/x)^3=1-t.
 \tag{A.38.6}
\]
It is the connected degree-nine cover classified by
$F_2\to(\mathbb Z/3)^2$. At $O$, the specified parameter is
$s=z/x$, so $t=-s^3$. The change from the tangent $-1$ to the
unit tangent is the based pullback
$t^{1/3^n}\mapsto-s^{3/3^n}$, using the compatible rational
odd roots of $-1$. It has zero Kummer cocycle. This retains the
given tangential arithmetic splitting. The Schreier transversal
$x^ay^b$, $0\le a,b<3$, gives exactly (A.38.3). It also gives
the actual conjugation actions of $x$ and $y$ on $N/K$; using
these actions and the stated class-two product is a finite
word-collection algorithm for (A.38.2).

For the chosen tangential paths the actual Ihara action is
$x\mapsto x^7$, $y\mapsto f^{-1}y^7f$, with $f\in F_2'$.
Its symmetry, hexagon and
pentagon are the actual Galois relations
\cite[\S\S I.3, I.6]{ASchneps2005}; we use the hexagon order
\[
 f(x,y)x^3 f(z,x)z^3 f(y,z)y^3=1,
 \qquad z=(xy)^{-1}.
 \tag{A.38.7}
\]
The comparison with A.7 is the integral unimodular matrix
saved by \texttt{a\_ihara\_lattice.py}. It is obtained from the
actual order-nine deck translations, and intertwines both
complex conjugation and the integral Frobenius homology matrix.
Consequently the finite equations below constrain the actual
arithmetic action in the specified geometric basis.

We give the finite elimination with its source-lifting
conditions. The exact collector for (A.38.3) works in
$F_2/K$, of order $3^{57}$. Its intersection with the lower
central and Zassenhaus filtrations is computed by taking
literal group commutators, cubes, and their invariant spans.
Through degree seven the kernel of the comparison from the
free source has one degree-six generator, $[x,y]^3$, and four
degree-seven generators. The literal words for all five are
included in the certificate and are individually checked to
map to one in $F_2/K$. No source direction below degree six is
lost. The exact commutator calculation is separately verified
by the full class-two multiplication law, retaining its
quadratic conjugation terms.

The homology equation, $f\in F_2'$, symmetry, and (A.38.7)
leave three possible homology rows and an affine central
system over each row. Requiring a lift of these same equations
over $\mathbb Z/9$ retains all next-digit variables. We then
impose the full group pentagon through source degree seven,
including the five actual source-kernel words just specified.
For precision, the pentagon is evaluated in the fourth-strand
free group of $P B_4=F_3\rtimes P B_3$. The Artin maps are
\[
 r_1:(X,Y,Z)\mapsto(Y,Y^{-1}XY,Z),\qquad
 r_2:(X,Y,Z)\mapsto(X,Z,Z^{-1}YZ).
 \tag{A.38.8}
\]
Put $A_{12}=r_1^2$, $A_{23}=r_2^2$,
$A_{13}=r_2r_1^2r_2^{-1}$, and
$A_{14}=\operatorname{Inn}(X)$,
$A_{24}=\operatorname{Inn}(Y)$,
$A_{34}=\operatorname{Inn}(Z)$.
These are the actual pure Artin maps
\cite[(2.1)--(2.3)]{DE_CohenFalkRandell2012}; direct free-word
reduction verifies every relation and
$A_{12}A_{13}A_{23}=\operatorname{Inn}((XYZ)^{-1})$.
The five factors have pairs
\[
\begin{split}
 &(A_{12},A_{23}A_{24}),\quad (A_{13}A_{23},A_{34}),\quad
 (A_{23},A_{34}),\\
 &(A_{12}A_{13},A_{24}A_{34}),\quad (A_{12},A_{23}).
\end{split}
 \tag{A.38.9}
\]
If they are denoted $F_0,\ldots,F_4$, the residual is
$F_0F_1F_4^{-1}F_3^{-1}F_2^{-1}$. Forgetting the fourth
strand sends this residual to one, so it is an element of the
actual free fiber and its displayed action is inner.

All evaluations use the denominator-free Magnus substitution
$X_i\mapsto1+X_i$ in the associative algebra over $\mathbb F_3$,
truncated after degree eight. If $\Delta_A=\Phi_A-1$, then
the source word $1+\sum_wc_ww$ acts as
$1+\sum_wc_w\Delta_{w_1}\cdots\Delta_{w_k}$.
Each $\Delta_A$ raises augmentation degree. Terms of source
degree at least eight therefore do not affect the target
generators through degree eight. Conversely, the first
nonzero degree-$d$ term $h_d$ of the fiber residual gives
$[h_d,X_i]$ in degree $d+1$ of its inner action. A positive-degree
tensor commuting with all three generators is zero. Hence
the computation detects precisely all source residuals through
degree seven. This uses the actual free fiber and the displayed
maps throughout.

The degree-two source coefficient must also be retained
integrally. Write $f=[x,y]^{f_2}$ modulo $\gamma_3F_2$.
The integral Heisenberg collector applied to the hexagon with
$m=(\chi-1)/2$ gives
\[
 3f_2=\frac{m(m+1)}2.
 \tag{A.38.10}
\]
This identity is in the torsion-free group $\mathbb Z_3$;
at $\chi=7$ it gives $f_2=2$ before any reduction. Its reduction
modulo nine is well-defined on the source used above, since
the Heisenberg group over $\mathbb Z/9$ kills $D_8(F_2)$.
The five source-kernel words have degree-two coefficients
$(3,0,0,0,0)$, so this final scalar row is imposed on the
combined source word, including those kernel corrections.

Here is the complete finite elimination result. ``Invisible''
means the kernel of the explicitly computed map to the source
metabelian quotient, followed by the intrinsic Johnson map.
\[
\begin{array}{c|rrr}
 \text{equations imposed}&\operatorname{rank}J&
       \operatorname{rank}(J,H)&\operatorname{rank}J|_{\ker H}\\ \hline
 \text{mod-nine symmetry and hexagon}&3&5&2\\
 \text{also full degree-seven pentagon}&3&4&1\\
 \text{also integral }f_2\bmod9&2&3&0
\end{array}
 \tag{A.38.11}
\]
The source metabelian map in this table is the literal Fox map:
write $f=[x,y]^h$ in $F_2'/F_2''$, with
$X=x-1$, $Y=y-1$. In the finite comparison its coefficient
ring is
\[
 \mathbb F_3[X,Y]/(X^6,X^3Y^3,Y^6,X^5Y^2,X^2Y^5).
 \tag{A.38.12}
\]
It is computed by $D_y(f)=Xh$. Thus every map in the last
column of (A.38.11) is an actual quotient map of the same
based free group.

The row reductions give a more precise conclusion. Let $J_0$
be the $45\times10$ central matrix obtained from
(A.38.4)--(A.38.5). There are two specified matrices $L_1,L_2$
such that every actual Frobenius in the initial basis has
\[
 J=J_0+aL_1+bL_2,\qquad
       a=h_{03}-2,\quad b=h_{05}\quad\text{in }\mathbb F_3.
 \tag{A.38.13}
\]
Here $h_{ij}$ is the coefficient of $X^iY^j$ in (A.38.12).
The matrices $L_1,L_2$, of respective ranks seven and one,
are saved with all 450 entries each. The nine matrices in
(A.38.13) are distinct. Formula (A.38.13) is at this point a
necessary arithmetic constraint; the realization follows below.

For reproducibility, the complete pentagon matrix has 29,520
coefficient rows. Its six candidate columns and five source-kernel
columns have rank seven, and the five kernel columns have rank
five. Each of the three branches has coefficient and augmented
rank seven on its 24 retained parameters. The additional scalar
row (A.38.10) has rank one in each branch. The supplied particular
solutions and full nullspaces verify every coefficient row.
The original and degree-pruned pentagon evaluations give identical
matrices; the independent integral Heisenberg audit verifies
(A.38.10) and every filtered solution against all 29,520 rows.
The files in the \texttt{code} directory with prefixes
\texttt{a\_ihara\_pentagon}, \texttt{a\_ihara\_source\_scalar},
\texttt{a\_ihara\_final\_parameters},
\texttt{de\_fermat\_pentagon}, and
\texttt{de\_fermat\_scalar\_audit} retain the exact words,
matrices, row reductions, and executed outputs. These finite
identities establish (A.38.11)--(A.38.13), including the affine
origin.

We now realize the stated value by an arithmetic morphism.
For a current actual basis $u,v=cuc$, two operations are
\[
\begin{array}{c|c|c}
 g&u'=gug^{-1},\quad v'=cu'c&(a',b')\\ \hline
 v&vuv^{-1},\quad c(vuv^{-1})c&(a-1,b+a-1)\\
 {[u,v]}&[u,v]u[u,v]^{-1},\quad c([u,v]u[u,v]^{-1})c&(a,b+1).
\end{array}
 \tag{A.38.14}
\]
These are direct identities of the full $55\times55$ matrices
$\left(\begin{smallmatrix}U&0\\J&\bigwedge^2U\end{smallmatrix}\right)$
and the fixed complex-conjugation matrix. The program
\texttt{a\_ihara\_frobenius\_conjugacy.py} verifies all nine
inputs of both operations. Applying the first operation $a$
times, using the representatives $a\in\{0,1,2\}$, changes the
pair to $(0,b+a(a-1)/2)$. Applying the second operation the
negative of that last second coordinate times gives $(0,0)$.
At most four operations are needed. Every step conjugates an
actual Frobenius, so the final element is an actual Frobenius
at the same prime of $K_0$.

The assignments $u\mapsto u'$, $v\mapsto cu'c$ preserve the
two independent Frattini classes and commute with $c$. They
therefore define a surjective endomorphism of the free pro-three
group $P_3$. A finitely generated profinite group is Hopfian:
its finitely many open normal subgroups of each fixed index are
permuted by inverse image under a surjection, forcing its kernel
into every open normal subgroup. Thus the endomorphism is an
actual $c$-equivariant automorphism $\psi$.

The full integral homology action is also retained. The
splitting after pullback by $[3]$ in A.7 is unique, hence
$G_{K_0}$-equivariant. Its isogeny
$T_D\times E\to J_{E,D}$ induces the rational equivariant
identification with $W_{\mathbb Q_3}\oplus T_{\mathbb Q_3}$.
The embedding of the chosen integral graph coordinates into
this diagonal representation is
\[
 \iota(w,t)=(w+\Psi t/3,t).
 \tag{A.38.15}
\]
The actual isogeny on Tate modules is $j(w,t)=(w,3t)$ in
the rational diagonal coordinates, or
$\iota^{-1}j(w,t)=(w-\Psi t,3t)$ in graph coordinates.
Thus the equivariant comparison retains its multiplication
by three on the proper factor.
On this rational direct sum every $g\in G_{K_0}$ acts by
$(\chi(g)I_8,T_g)$, where $T_g$ is multiplication in the
commutative algebra $\mathbb Q_3[\zeta]$. These actions commute.
The integral lattice embeds in this rational representation,
so their integral actions commute as well. Arithmetic
conjugation consequently preserves the exact A.7 matrix.

Finally we give the transport of every parameter used later.
For the original action $\rho$, the new action is $\rho\circ\psi$.
The exact semidirect-product isomorphism is
$ (N/K)\rtimes_{\rho\circ\psi}\Gamma\to
 (N/K)\rtimes_\rho\Gamma$, $(h,g)\mapsto(h,\psi(g))$,
where $\psi(c)=c$; its inverse uses $\psi^{-1}$.
A cocycle is pulled back by $z\mapsto z\circ\psi$. In particular,
\[
 z(gug^{-1})=
 z(g)\,\rho(g)(z(u))\,
       \rho(gug^{-1})(z(g)^{-1}),
 \tag{A.38.16}
\]
with the additive version
$z(g)+\rho(g)z(u)-\rho(gug^{-1})z(g)$ for an abelian module.
This formula retains the reference section and its origin.
In the free presentation of A.28, make the same substitutions
on $x,y$ and fix $t$. They again define a $c$-equivariant IA
automorphism, and the actual new Frobenius word is
\[
 W_{q,\mathrm{new}}=\psi^{-1}(W_{q,\mathrm{old}}).
 \tag{A.38.17}
\]
Indeed the old quotient map composed with $\psi$ sends this
word to the same local Frobenius and sends the transformed
relator to $W_{q,\mathrm{old}}tW_{q,\mathrm{old}}^{-1}t^{-q^2}$.
The definitions of $m,n$ by the pro-three square-root operations
commute with this actual automorphism. An IA automorphism of
a free pro-three group acts identically on its associated
restricted Lie algebra. Since A.30 proves $W_{17}\in D_3$,
the initial terms in A.31 and A.37 survive unchanged modulo
$D_4$.

For completeness this also preserves the ordinary weight-three
coefficients modulo three used to identify those terms.
The degree-one and degree-two Hall coefficients of $W_{17}$
are divisible by three. Under a basis-conjugating automorphism,
the changes to a degree-one power with exponent $3a$ have
weight-two and weight-three coefficients divisible by $3a$
or $\binom{3a}{2}$; both are divisible by three. A weight-two
power with exponent $3b$ changes by weight-three multiples
of $3b$. The automorphism is the identity on the third lower
central graded quotient. Hence the stated ordinary coefficients
are unchanged as well. All higher terms, and the integral
lower coordinates when needed, are transported by (A.38.17).
This proves the realization and every specified comparison.
\end{proof}

\begin{statement}{A.39}{Proved}{Global exclusions at seventeen inside the exact local cusp packet}
Use the actual arithmetic basis of A.38, with $m,n$ and the
normalized local Frobenius $\varphi=\varphi_{17}$ transported
as there. Put $M_2=\bigwedge^2V$ and define the integral
normalized cochain module and its actual evaluation map by
\[
\begin{split}
 D&=M_2^{c=1}\oplus M_2^{nc=-1},\\
 E_\varphi:D&\longrightarrow M_2^{c=1},\quad
       (a,b)\longmapsto\gamma_{a,b}(\varphi),\\
 R&=\ker E_\varphi.
\end{split}
 \tag{A.39.1}
\]
Here $\gamma_{a,b}$ is the unique continuous cocycle on
$P_3\rtimes\langle c\rangle$ with values $(a,b,0)$ at
$(m,n,c)$. The lattice $D$ has rank $45$, and $R$ is a
saturated lattice of rank $24$.

For $r\in R$, twist the class-two image of the global
$O$-cuspidal section with tangent parameter $17$ by
$\gamma_r$. This defines an injectively parameterized family
of global class-two section classes, unramified outside
$\{3,17,\infty\}$. Every member has the same exact local
$O$-cusp section at seventeen, including its Frobenius
coordinate, and every member has a global class-three lift.

There is an $\mathbb F_3$-linear obstruction on $R/3R$ of
rank at least three which vanishes on every member lifting to
any of
\[
 N,\qquad N/N'',\qquad N/\gamma_3(I_{Q_c}).
 \tag{A.39.2}
\]
Consequently, from any integral basis $r_1,\ldots,r_{24}$ of
the kernel in (A.39.1), at least three of the resulting
twenty-four class-two section classes have none of the lifts
in (A.39.2). This is an unconditional finite-family existence
certificate. It does not select an unevaluated individual
basis vector as a counterexample. No assertion of full local
liftability at three is included.
An individually specified member of $R$ with nonzero
obstruction, including all integral generator coordinates
and a nine-term detecting functional, is given in B.20.
\end{statement}
\begin{proof}
We first prove the integral assertions, before using any
finite computation. The relations $c(m)=m$, $c(n)=n^{-1}$
give precisely the two conditions in (A.39.1). Since $m,n$
are a free pro-three basis, they are sufficient. The signs of
$c$ on $V$ are $(4,6)$, so its signs on $M_2$ are $(21,24)$.
The involution $nc$ is conjugate to $c$ by a pro-three square
root of $n$, and therefore has the same signs. Its integral
eigenspaces are direct summands, since two is a unit in
$\mathbb Z_3$. Thus $D$ is free of rank $45$.

Since $\varphi$ is $c$-fixed, its cocycle value belongs to
$M_2^{c=1}$. For $h$ in that invariant lattice, the global
coboundary $g\mapsto(g-1)h$ belongs to $D$ and evaluates to
$(F_\varphi-1)h$. The rational eigenvalues of $F_\varphi$ on
$M_2$ are
\[
                  17^4,\qquad -17^3,\qquad17^2.
 \tag{A.39.3}
\]
Indeed on the rational equivariant splitting of A.7 the toric
factor has eigenvalue $17^2$ and the proper factor has scalar
eigenvalue $-17$. None of (A.39.3) is one. The map
$F_\varphi-1$ is consequently invertible over $\mathbb Q_3$
on the $21$-dimensional $c$-fixed target. Evaluation has
rational rank $21$. Its kernel $R$ has rank $24$, and
$D/R\simeq\operatorname{im}E_\varphi$ embeds in a free
$\mathbb Z_3$-module. Thus $R$ is saturated and
\[
                 R/3R\hookrightarrow D/3D
 \tag{A.39.4}
\]
has dimension $24$.

For $r\in R$ the unramified-outside-three cocycle has zero
value at $c$ and $\varphi$ and zero tame inertia value. Its
restriction to the local group at seventeen is therefore zero.
This proves the exact local class-two assertion. The family
is injective on global section classes as well. A geometric
conjugator between two such class-two sections has invariant
image in $V$, which is zero by Frobenius weights. It therefore
lies in $M_2$. Normalization at $c$ makes it $c$-fixed. If the
difference of two members of $R$ were its coboundary, evaluation
at $\varphi$ would give $(F_\varphi-1)h=0$, hence $h=0$ by
(A.39.3).

Here is a definite integral class-three lift of each member.
Write
\[
 L=\gamma_2N/\gamma_4N,
 \qquad 0\longrightarrow M_3\longrightarrow L
                    \longrightarrow M_2\longrightarrow0.
 \tag{A.39.5}
\]
This is an exact sequence of free $\mathbb Z_3$-modules.
Lift the two values of $r$ equivariantly for the involutions
$c$ and $nc$, using their averaging idempotents. In the free
presentation of A.28 take these as the $L$-valued cochain
values at $m,n$, keep the reference tame section value $c_O$,
and keep the value at $c$ equal to one. The resulting crossed
cochain $\beta$ has $\beta(t)=0$. Its value at the actual
word $W$ projects to $E_\varphi(r)=0$, so $\beta(W)\in M_3$.
That value is central in $N/\gamma_4N$ and commutes with the
unchanged tame section. The relation
$WtW^{-1}=t^{289}$ is therefore satisfied exactly. Equivalently,
its crossed value is $(1-T^{289})\beta(W)=0$, where $T$
is the reference tame action on $L$ and is the identity on
$M_3$. The proved one-relator presentation supplies the global
class-three lift.

We describe the finite comparison computing its remaining
local obstruction. It uses the actual action of A.38 and
does not require degree-three images of the geometric
generators. Let $X_i$ be their logarithms in the rational
class-three nilpotent Lie algebra. The logarithm maps $L$
onto the integral lattice with degree-two basis $[X_i,X_j]$
and degree-three Hall basis $[[X_i,X_j],X_k]$, with the usual
Jacobi identifications. To check integrality, the sole
correction needed is
\[
 \log[x_i,x_j]=[X_i,X_j]
       +\tfrac12[X_i+X_j,[X_i,X_j]],
 \tag{A.39.6}
\]
and two is invertible in $\mathbb Z_3$. Thus this is an
actual lattice comparison. If a class-two automorphism has
homology matrix $A$ and central columns $J_i$, its action on
$\overline L=L/3L$ is
\[
 \widetilde G=
 \begin{pmatrix}\bigwedge^2A&0\\
       \Theta(A,J)&\operatorname{Lie}_3(A)\end{pmatrix},
 \quad
 \Theta(A,J)_{ij}=[J_i,A X_j]+[A X_i,J_j].
 \tag{A.39.7}
\]
Unknown degree-three terms in the images of $X_i$ enter these
brackets only in degree four and hence disappear in (A.39.5).
Formula (A.39.7) is therefore determined by A.38.

For the $O$ reference section, $e_O$ is the seventh Schreier
generator of A.38 and
\[
                     T(z,w)=(z,w+[e_O,z]).
 \tag{A.39.8}
\]
The source generators $m,n$ have zero $17$-Kummer value by
A.35 and its transported construction. Their actions on $L$
are consequently the original arithmetic actions. The
program \texttt{a\_ihara\_gamma2\_action.py} constructs
(A.39.7), verifies composition and every involution and tame
relation, and compares the entire degree-three block with the
independent integral lattice of A.13. For the augmentation
ideal generated by $m-1,n-1,t-1$ on $\overline L$, the dimensions
of successive images are
\[
                  375,\ 249,\ 132,\ 47,\ 4,\ 0.
 \tag{A.39.9}
\]
In particular its fifth power is zero. A crossed cochain
therefore kills $D_6$ of the arithmetic free group: the
derivation of the sixth augmentation power lies in the fifth
power applied to the module.

Every source term that can affect the answer is retained.
The group $D_3/D_6$ is abelian of exponent three. On it the
map
\[
                  \operatorname{av}(w)=(w\,c(w))^2
 \tag{A.39.10}
\]
is exactly the $c$-fixed averaging idempotent. A.31 and A.37
give the known cubic part
$m^6[n,[n,m]][t,[t,m]]^2$. The two unknown cubic terms have
Lyndon words $012,021$ in the order $(m,n,t)$. All eighteen
degree-four and all forty-eight degree-five Lyndon terms
are then retained, with (A.39.10) applied to each. These
exhaust the source freedom: the only restricted-power terms
in these degrees are the degree-three generator cubes, and
their coefficients are the already fixed abelian Frobenius
coordinates. The source comparisons of A.38 ensure that
these are the actual transported coefficients of $W$.

The local mod-three $M_3$ target is its $c$-fixed space of
dimension $164$. Divide by the actual rank-fifteen global
restriction image in A.34; denote the resulting linear
projection by $Q$, with $149$-dimensional target. The factor
two from norm Frobenius to rational Frobenius is a unit and
does not affect any of the following ranks. The canonical
integral lifts in (A.39.5) give an ambient linear map on
$D/3D$. Evaluating every term just listed and applying $Q$
shows that only the following three unknown terms survive:
\[
 [m,[n,t]],\qquad [[m,t],n],\qquad [m,[[m,n],n]].
 \tag{A.39.11}
\]
Every degree-five contribution vanishes. We retain the three
actual coefficients as arbitrary members of $\mathbb F_3$;
the rank bound below holds for all twenty-seven triples.

All local conjugators and global class-three corrections are
also retained. The image of a local conjugator in $V$ is
invariant, hence zero by weights. Its degree-two part belongs
to $M_2^{c=1}$ after normalization at $c$. Choose any of its
$21$-dimensional equivariant lifts $h$ to $\overline L$.
Allow the enlarged finite comparison in which
\[
 \beta(t)=(T-1)h,\qquad
 Q\bigl(\beta(W)-(F_\varphi-1)h\bigr)
 \tag{A.39.12}
\]
is evaluated, while the values at $m,n$ are the chosen
cochain lifts. This includes every reduction of a local
comparison coming from a putative lift. Degree-three changes
of $h$ are already killed by the local coboundary quotient.
Allowing every degree-two $h$ only enlarges the set of finite
corrections, so a nonzero obstruction after this allowance
still excludes an integral lift.

Write the map (A.39.12), for its actual coefficient triple
$\epsilon$, as $G_\epsilon r+H_\epsilon h$. The complete
finite calculation is
\[
\begin{array}{c|rrrr}
 \epsilon\in\mathbb F_3^3&
 \operatorname{rank}G_\epsilon&
 \operatorname{rank}H_\epsilon&
 \operatorname{rank}(G_\epsilon,H_\epsilon)&
 \operatorname{rank}(G_\epsilon\bmod\operatorname{im}H_\epsilon)\\ \hline
 \text{each of the twenty-seven triples}&30&6&30&24.
\end{array}
 \tag{A.39.13}
\]
Here $G_\epsilon$ has $45$ columns and $H_\epsilon$ has
$21$ columns. The images of all $H_\epsilon$ are the same
six-dimensional subspace, so the last column can be computed
in one fixed $143$-dimensional target. For a reproducible
specification of all their
entries, take the ascending pivot columns of $1+c$ and
$1-nc$ on $M_2/3$ for the two source bases, lift them by the
same operators on $L/3$, and evaluate the words above using
\[
\begin{split}
 \beta(ab)&=\beta(a)+a\beta(b),\\
 \beta([a,b])&=(1-aba^{-1})\beta(a)
                         +(a-[a,b])\beta(b).
\end{split}
 \tag{A.39.14}
\]
Use the literal Jacobi reductions in the basis of A.13 and
the projection $Q$ just defined. These formulas and A.38
specify every entry without an arithmetic oracle. The program
\texttt{a\_frobenius17\_mixed\_gluing.py}, run with
\texttt{--local-conjugator}, saves the full matrices, all
twenty-seven coefficient cases, and their exact elimination
ranks. Its certificate also saves each of the sixty-eight
candidate word maps before discarding the zero ones. In
particular (A.39.13) retains the two mixed inertia terms and
the quartic term, rather than assigning their values.
The independent program \texttt{bc\_gluing\_audit.py} instead
uses truncated noncommutative Magnus polynomials and the
derivation formula
$\beta((g_1-1)\cdots(g_d-1))=
(g_1-1)\cdots(g_{d-1}-1)\beta(g_d)$.
It reproduces every entry of all sixty-eight word maps,
the common six-dimensional image, and all twenty-seven
rank computations.

Any two global class-three lifts of the fixed class-two
target differ by an $M_3$ cocycle. Such a cocycle is unramified
outside three: at every other finite prime the possible tame
value lies in $H^0(\mathbb F_q,M_3(-1))$, which is zero by
weights. Its effect is therefore in the actual global
restriction image already removed by $Q$; this is also the
localization comparison in A.33. Thus vanishing of
$G_\epsilon r$ modulo $\operatorname{im}H_\epsilon$ is a
necessary condition for a global class-three lift whose
restriction at seventeen is the actual cusp section.

For the actual, unassigned triple $\epsilon$, (A.39.13) says
that the kernel of this ambient map on the $45$-dimensional
space $D/3D$ has dimension $21$. The embedded subspace $R/3R$
in (A.39.4) has dimension $24$. Its restriction consequently
has rank at least $24-21=3$. This is the asserted obstruction.
It is obtained from actual integral cochains in $R$; no
mod-three local condition has been substituted for the exact
kernel in (A.39.1).

Finally every lift in (A.39.2) would supply the forbidden
locally cuspidal class-three lift. The full group maps to
$N/\gamma_4N$. The metabelian quotient does also, because
$N''\subset\gamma_4N$, and its local section is the actual
$O$ cusp packet by A.19. For the marked quotient, the image
of $I_{Q_c}$ in $V$ is the single line generated by $e_{Q_c}$.
Hence $[I_{Q_c},I_{Q_c}]\subset\gamma_3N$ and
$\gamma_3(I_{Q_c})\subset\gamma_4N$, proving the actual map
to the same class-three quotient. The marked cusp is $Q_c$,
whereas the ramified cusp here is $O$; B.14 identifies its
local marked section with the unique actual $O$ cusp section.
The local tame exponent is one, and its unramified cusp
parameter is zero because $17-1$ is a three-adic unit.
Thus all three types of lift are excluded whenever the
obstruction is nonzero. A basis of $R$ spans its reduction,
so an obstruction of rank at least three is nonzero on at
least three basis members. This proves the finite-family
assertion and the stated global exclusions.
\end{proof}

\begin{statement}{A.40}{Proved}{The excluded family also fails at the next nilpotent stage}
Every member of the family in A.39 with nonzero stated
obstruction has no global lift to $N/\gamma_5N$. Thus the
twenty-four-member finite family there contains at least three
class-two sections which lift globally to $N/\gamma_4N$ and
have no global lift to $N/\gamma_5N$.
The explicitly evaluated section in B.20 is one such member.
\end{statement}
\begin{proof}
First conjugate a proposed lift, through the actual surjection
to $N/\gamma_3N$, so that its class-two projection is the
specified cocycle. The coprime involution argument in the
kernel then normalizes its value at $c$ to one. Every global
class-three lift differs from the lift constructed in A.39
by an $M_3$ cocycle. A.39 proves that this difference is
unramified outside three, even when additional ramification
was initially allowed. Hence its actual tame inertia at
seventeen is still exactly $c_O$ in $N/\gamma_4N$.

Write $F_0$ for the reference rational Frobenius and $F'$ for
the Frobenius of the proposed lift. In class three their
difference is an element $z\in M_3$. The relative local group
has the exact relation
\[
       \operatorname{Frob}_{17}=c\,\varphi_{17}^{1/2}.
 \tag{A.40.1}
\]
Its involution commutes with $\varphi_{17}$, so normalization
at $c$ makes $z$ $c$-fixed. The integral A.7 matrix gives
$\varphi_{17}=1$ on $V/3$. Thus its square root is also one
on every graded quotient modulo three, and the cocycle norm
operator $1+\rho(\varphi_{17}^{1/2})$ is $2I$ modulo three.
The rational-Frobenius difference $z$ is consequently one
half of the norm-Frobenius difference modulo three. Nonzero
obstruction in A.39 implies
\[
                 \bar z\ne0,\qquad C\bar z=\bar z
                       \quad\text{in }M_3/3M_3,
 \tag{A.40.2}
\]
after every permitted global class-three correction.

Let $M_4=\gamma_4N/\gamma_5N$. In class four write the inertia
and Frobenius of the hypothetical lift in the exact order
\[
             \tau'=c_Oa_4,\qquad F'=z b_4F_0,
                       \qquad a_4,b_4\in M_4.
 \tag{A.40.3}
\]
A choice of lift of $z$ from $M_3$ only changes $b_4$.
The reference cusp relation is $F_0c_OF_0^{-1}=c_O^{17}$.
The group $M_4$ is geometrically central at this stage, and
collecting the relation $F'\tau'F'^{-1}=(\tau')^{17}$ gives
\[
             (F_{17}-17)a_4=-17[z,e_O]\quad\text{in }M_4.
 \tag{A.40.4}
\]
Indeed conjugating $c_O^{17}$ by the chosen lift of $z$
contributes $17[z,e_O]$; $b_4$ commutes with inertia, while
the two contributions from $a_4$ are $F_{17}a_4$ and
$17a_4$. This derives the sign in the stated ordering and
retains the actual reference relation.

On $M_4/3M_4$ the action of rational Frobenius is $C$ by
(A.40.1). Since $17\equiv-1\pmod3$, equation (A.40.4)
becomes
\[
                    (C+1)\bar a_4=[\bar z,e_O].
 \tag{A.40.5}
\]
The left side is $c$-invariant. The right side is
$c$-anti-invariant, because $C\bar z=\bar z$ and
$Ce_O=-e_O$. It remains to check its nonvanishing in the
actual ordinary lower-central graded quotient.

The ordinary free Lie algebra on $V/3$ embeds in its tensor
algebra: the Hall--Lyndon expansions have distinct unit
leading words, giving an integral comparison before reduction.
The element $e_O$ is a primitive basis letter. A homogeneous
tensor $P$ of degree three commuting with that letter must
be a scalar multiple of $e_O^3$. To see this directly, the
equality $e_OP=Pe_O$ forces every word of $P$ to start with
$e_O$; cancel that initial letter and repeat. The ordinary
degree-three free Lie subspace has zero component in the
multidegree consisting of that letter alone, since every
bracket on one letter vanishes. Thus
\[
       [\bar z,e_O]=0\text{ in }M_4/3M_4
                     \quad\Longrightarrow\quad\bar z=0.
 \tag{A.40.6}
\]
By (A.40.2) the right side of (A.40.5) is nonzero. The plus
and minus eigenspaces of an involution over $\mathbb F_3$
intersect trivially, so (A.40.5) is impossible. This excludes
every global class-four lift of that class-two target, while
A.39 already supplies its global class-three lift.
\end{proof}

\begin{statement}{A.41}{Proved}{The exact completed Fox image and fibers on the cusp family}
For the family $R$ in A.39 put
\[
 Q=N/N'',\qquad M=N'/N'',\qquad
 \Lambda=\mathbb Z_3[[V]],\qquad
 \mathfrak a=\ker(\Lambda\to\mathbb Z_3),\qquad J=\mathfrak aM.
 \tag{A.41.1}
\]
Give $M$ the arithmetic action twisted by the actual reference
$O$ section with tangent parameter $17$. Thus $m,n,c$ have
their original actions, and tame inertia acts by conjugation
by $c_O$. Define
\[
\begin{split}
 D_M&=M^{c=1}\oplus M^{nc=-1},\qquad
 D_J=J^{c=1}\oplus J^{nc=-1},\\
 E_W:D_M&\longrightarrow M^{c=1},\qquad
 E_W(a,b)=\beta_{a,b}(W),
\end{split}
 \tag{A.41.2}
\]
where $\beta_{a,b}$ is the continuous crossed cochain on the
free arithmetic source with values $(a,b,0,0)$ at
$(m,n,t,c)$. These are continuous linear maps of compact
$\mathbb Z_3$-modules. The natural quotient $M\to M_2$ gives
$\pi:D_M\twoheadrightarrow D$, with kernel $D_J$, and the
complete metabelian image on this family is exactly
\[
                  R_{\rm met}=\pi(\ker E_W)\subset R.
 \tag{A.41.3}
\]
Equivalently, it is the kernel of the explicit continuous map
\[
\begin{split}
 \mathfrak o_M:R&\longrightarrow
                   J^{c=1}/E_W(D_J),\\
 r&\longmapsto[\,E_W(\widetilde r)\,],
\end{split}
 \tag{A.41.4}
\]
where $\pi(\widetilde r)=r$. Every nonempty fiber of
$\operatorname{Sec}(Q)\to\operatorname{Sec}(N/\gamma_3N)$
over this family is an affine copy of
\[
             \ker(E_W:D_J\longrightarrow J^{c=1}).
 \tag{A.41.5}
\]
There is no remaining conjugacy quotient after the actual
local normalization in the proof. Formulas (A.41.2)--(A.41.5)
retain the completed arithmetic Fox operator. They do not
assign a numerical value to an unevaluated coefficient of
$W$ or of its action.
\end{statement}
\begin{proof}
We construct the completed module and its maps. Let
$x_1,\ldots,x_{10}$ be the actual geometric Schreier basis
in A.38, and denote their images in $V$ by $T_i$. The
completed Fox map gives the exact sequence
\[
 0\longrightarrow M\longrightarrow\Lambda^{10}
    \xrightarrow{(T_1-1,\ldots,T_{10}-1)}\mathfrak a
      \longrightarrow0.
 \tag{A.41.6}
\]
One direct proof uses the actual finite abelian covers with
geometric subgroups $H_k=N'N^{3^k}$. Their based cellular
chain maps send the $i$th lifted edge to $T_i-1$. The kernel
is $H_k^{\rm ab}$. Inverse limits of these compact chain
groups are exact, and the descending-subgroup comparison
proved in A.19 identifies
$\varprojlim H_k^{\rm ab}$ with $N'/N''$. This proves
(A.41.6) and its based comparison. The action of geometric
conjugation by $x_i$ on $M$ is multiplication by $T_i$.
Consequently
\[
  J=\mathfrak aM=\gamma_3N/N'',\qquad
  M/J=N'/\gamma_3N=\bigwedge^2V=M_2.
 \tag{A.41.7}
\]
These equalities identify the quotient map used in the
statement, including its integral lattice.

The algebra $\Lambda$ is
$\mathbb Z_3[[T_1-1,\ldots,T_{10}-1]]$, an integral domain.
In (A.41.6) the reference tame action $T$ is multiplication
by the actual deck element $T_O$, the image of $c_O$.
Therefore
\[
                   1-T^{289}:M\longrightarrow M
                         \quad\hbox{is injective}.
 \tag{A.41.8}
\]
Indeed it is multiplication by the nonzero power series
$1-T_O^{289}$ on a submodule of the torsion-free module
$\Lambda^{10}$. This also proves that the tame fixed
subgroup of $M$ is zero.

The arithmetic action on $M$ is semilinear over its actual
action on $\Lambda$: $g(\lambda z)=g(\lambda)g(z)$.
Its matrices are the Fox matrices of the actual based
geometric automorphisms. Thus the continuous group-ring
homomorphism from the free arithmetic source to
$\operatorname{End}_{\mathbb Z_3}(M)$ and the completed Fox
derivatives specify (A.41.2) explicitly as
\[
 E_W(a,b)=\rho\!\left(\frac{\partial W}{\partial m}\right)a
          +\rho\!\left(\frac{\partial W}{\partial n}\right)b.
 \tag{A.41.9}
\]
The derivative with respect to $t$ has zero cochain value.
Here $\rho$ is the reference twisted action just defined;
neither its semilinear part nor the inertia dependence of
$W$ is removed. Since the word $W$ is $c$-fixed and the
cochain is normalized at $c$, its value is $c$-fixed.
This proves the target in (A.41.2).

Averaging for the order-two actions lifts both eigenspaces
through the surjection $M\to M_2$. Hence $\pi:D_M\to D$ is
onto with kernel $D_J$. Reduction of (A.41.9) is exactly
$E_\varphi$, because the tame action on $M_2$ is trivial and
killing $t$ sends $W$ to the same actual $\varphi$. For
$r\in R$, $E_W(\widetilde r)$ therefore lies in $J^{c=1}$.
Two choices of $\widetilde r$ differ by $D_J$, proving
(A.41.4) is well-defined. The domains are compact and all
maps are continuous, so $E_W(D_J)$ is closed. Choosing
equivariant lifts of a finite basis of $D$ makes (A.41.4)
an explicitly continuous linear map.

We next prove that these kernels classify the actual section
classes. A cocycle $\beta$ with the values in (A.41.2),
relative to the full reference cusp section, is a cocycle
on the free arithmetic source. Its value on the relator is
\[
 \beta(WtW^{-1}t^{-289})=(1-T^{289})\beta(W).
 \tag{A.41.10}
\]
By (A.41.8) it descends to the proved arithmetic presentation
exactly when $E_W(a,b)=0$. Its values at $t,c,W$ then vanish,
so its local section at seventeen is the literal reference
section. Its class-two image lies in $R$ by the preceding
comparison. This constructs actual global metabelian sections
from every member of $\ker E_W$.

Conversely, take any global metabelian section whose class-two
class is the target indexed by $r\in R$. Lift a global
class-two conjugator through $Q\to N/\gamma_3N$ so that its
projection equals the exact target cocycle; then normalize
the value at $c$ in the kernel. At every prime outside three
and seventeen its abelian inertia is zero. A.19 implies that
its entire metabelian inertia there is zero. Thus it factors
through the arithmetic group of A.28 even if arbitrary
additional ramification was initially allowed.

At seventeen A.19 places it in the actual $O$ packet. The
tame exponent is one, and the unramified cusp parameter
vanishes since $17-1$ is a three-adic unit. Hence a local
geometric conjugator identifies it with the reference cusp
section. Its image in $N/\gamma_3N$ centralizes that local
section. This centralizer is trivial: its image in $V$ is
Frobenius-invariant and therefore zero; its remaining image
in $M_2$ is likewise zero by (A.39.3). Consequently the
conjugator lies in $J$. The normalization at $c$ makes it
$c$-fixed. Conjugating globally by this same element preserves
the exact class-two target and makes the whole local section
equal to the reference.

The resulting difference from the reference is an $M$-valued
cocycle, because its abelian projection is the same. Its
values at $t,c,W$ are zero, so it gives a member of
$\ker E_W$ reducing to $r$. The normalization is unique:
the ratio of two such conjugators is tame-invariant in $M$,
and (A.41.8) makes that ratio zero. More generally, a global
conjugator between two normalized representatives has
invariant image in $V$, hence lies in $M$, and its local
inertia commutation again forces it to be zero. This proves
the claimed bijection with section classes, including the
absence of a remaining quotient.

The image formula (A.41.3), the obstruction formula
(A.41.4), and the affine fiber (A.41.5) now follow from the
exact sequence of these actual linear maps. In particular
$R_{\rm met}$ is a closed $\mathbb Z_3$-submodule of the
finite free lattice $R$. Its restriction has the explicit
nonzero obstruction of A.39; its integral image has not been
replaced by that first finite obstruction.

There is also an exact finite-quotient interpretation. Use
the cofinal, arithmetic-stable quotients
\[
                  M_{a,b}=M/(3^aM+\mathfrak a^bM),
                             \qquad a,b\ge1.
 \tag{A.41.11}
\]
They are finite, since (A.41.6) is a finitely generated
module over the noetherian ring $\Lambda$. The formal-series
topology and a continuous surjection $\Lambda^r\to M$
show that these quotients are cofinal: a surjective
homomorphism of compact Hausdorff groups is open, and
the images of the powers of $(3,\mathfrak a)$ form the
corresponding neighborhood basis. The identical argument
applies to the finitely generated module $J$. For each pair,
take the actual finite images of the arithmetic action and
of the completed Fox derivatives in (A.41.9). The finite
system asks for the two equivariant cochain values reducing
to $r\bmod3^a$ and satisfying $E_W=0$ in $M_{a,b}$.
Membership in (A.41.3) is equivalent to solvability of every
one of these finite linear systems. To prove the reverse
direction, lift their solution conditions to closed subsets
of the compact module $D_M$. Every finite collection is
satisfied at a common refinement $(a,b)$, so these subsets
have the finite intersection property. Their intersection
gives an actual solution in $D_M$. This establishes the
integral criterion without an assumption that the transition
maps on solution sets are onto. The finite inputs are the
actual images of $W$ and of the action; no calculation of
them is inferred merely from writing the operator.
\end{proof}

\begin{statement}{A.42}{Proved}{A complete affine Fox criterion for every one-prime inert metabelian target}
Let $q\ne3$ be inert in $K_0$, and let $s$ be any class-two
section unramified outside $\{3,q,\infty\}$. Use the actual
presentation of A.28, written in the free basis $m,n,t$ with
$c(m)=m$, $c(n)=n^{-1}$, $c(t)=t^{-1}$, and put $Q_q=q^2$.
Retain its actual $c$-fixed Frobenius word $W$. Normalize
$s(c)=1$, and choose an equivariant lift $\widetilde s$ of
its three generator values to $Q=N/N''$ on this free source.
Such a lift always exists before imposing the single relator.

Let $\rho_\alpha$ be the action on $J=\gamma_3N/N''$ obtained
by conjugation by $\widetilde s$, where $\alpha$ is the
prescribed abelian section. Set
\[
 U=\rho_\alpha(m),\quad B=\rho_\alpha(n),\quad
 T=\rho_\alpha(t),\quad C=\rho_\alpha(c),\quad
 \Phi=\rho_\alpha(W).
 \tag{A.42.1}
\]
This action depends only on $\alpha$ and descends to the
actual arithmetic group. In particular
$\Phi T\Phi^{-1}=T^{Q_q}$. Define
\[
\begin{split}
 \mathcal D_\alpha&=
 J^{C=1}\oplus J^{BC=-1}\oplus J^{TC=-1},\\
 \mathcal J_\alpha&=J^{T^{Q_q}C=-1},\\
 \mathcal E_W(a,b,d)&=
 \rho_\alpha\!\left(\frac{\partial W}{\partial m}\right)a+
 \rho_\alpha\!\left(\frac{\partial W}{\partial n}\right)b+
 \rho_\alpha\!\left(\frac{\partial W}{\partial t}\right)d,\\
 \mathcal L_s(a,b,d)&=(1-T^{Q_q})\mathcal E_W(a,b,d)
                +\left(\Phi-\sum_{i=0}^{Q_q-1}T^i\right)d,\\
 e_s&=\widetilde s(WtW^{-1}t^{-Q_q})\in\mathcal J_\alpha.
\end{split}
 \tag{A.42.2}
\]
The full global metabelian image consists exactly of the
targets for which
\[
             [e_s]=0\quad\text{in}\quad
                  \mathcal J_\alpha/\mathcal L_s(\mathcal D_\alpha).
 \tag{A.42.3}
\]
For every target in that image, its fiber is an affine torsor
under the explicitly specified compact group
\[
 \frac{\ker\mathcal L_s}
 {\{((U-1)h,(B-1)h,(T-1)h):h\in J^{C=1}\}}.
 \tag{A.42.4}
\]
All source actions, the affine error, and the finite images of
$W$ in these formulas are the actual based ones. The criterion
is an exact completed linear system, including its affine
origin and every conjugation. Its numerical evaluation for
arbitrary input is not asserted.
\end{statement}
\begin{proof}
Use the completed module and actual quotient maps of
(A.41.6)--(A.41.7). The kernel of
$Q\to A=N/\gamma_3N$ is exactly the abelian compact module
$J$. Every section can be normalized at $c$ by the coprime
involution argument. A generator value at $m$ is fixed by
$c$, and it lifts through $Q^c\to A^c$. A generator value
at $n$ is a cocycle for the involution $nc$; it lifts by
coprime complement conjugacy. The same assertion holds for
the involution $tc$ and the tame generator. These three
choices are independent in the free source with operators.
They give the stated equivariant $\widetilde s$ before any
relator equation is required.

Two such lifts of a generator differ by $J\subset M$.
Conjugation by $M$ is trivial on the abelian module $M$,
and any two generator values with the same image in $V$
differ by $M$. Thus the induced action on $M$, and on $J$,
depends only on the abelian generator value. This proves its dependence
only on $\alpha$. The relator of A.28 maps under
$\widetilde s$ to an element of $J$, since the class-two
target is an actual section. This element acts trivially
on $J$. Thus $\rho_\alpha$ kills the relator and descends
to the actual arithmetic group. Its involutions satisfy
$C^2=(BC)^2=(TC)^2=1$, and its local relation is exactly
$\Phi T\Phi^{-1}=T^{Q_q}$.

All other equivariant free-source lifts with the same exact
class-two target are uniquely of the form
\[
                    \widetilde s'(g)=\beta(g)\widetilde s(g),
 \tag{A.42.5}
\]
where $\beta$ is a $J$-valued crossed cochain for
$\rho_\alpha$, normalized at $c$. The source involution
relations give
$C\beta(m)=\beta(m)$,
$BC\beta(n)=-\beta(n)$, and
$TC\beta(t)=-\beta(t)$. Hence their generator values are
precisely the domain $\mathcal D_\alpha$ in (A.42.2).
The completed Fox formula gives
$\beta(W)=\mathcal E_W(a,b,d)$, with no discarded tame term.

Write $r=WtW^{-1}t^{-Q_q}$. Direct crossed-cochain collection
gives
\[
 \beta(r)=(1-T^{Q_q})\beta(W)
             +\left(\Phi-\sum_{i=0}^{Q_q-1}T^i\right)\beta(t).
 \tag{A.42.6}
\]
For example the first three factors give
$(1-\Phi T\Phi^{-1})\beta(W)+\Phi\beta(t)$;
the last factor subtracts the displayed geometric sum.
This proves the operator $\mathcal L_s$. The source identity
$c(r)=t^{-Q_q}r^{-1}t^{Q_q}$ implies
\[
                  C e_s=-T^{-Q_q}e_s.
 \tag{A.42.7}
\]
It gives the identical relation for $\beta(r)$, proving
the target $\mathcal J_\alpha$ of the operator.

The error of (A.42.5) on the relator is the sum
$e_s+\mathcal L_s(a,b,d)$ in the actual abelian kernel $J$.
The proved presentation therefore shows that it descends to
a global metabelian section exactly when this sum is zero.
This proves (A.42.3). The image of $\mathcal L_s$ is closed,
because its domain is compact. A different equivariant free
lift changes $e_s$ by $\mathcal L_s$ of the corresponding
cochain and leaves $\rho_\alpha$ unchanged. Thus it translates
the affine solution space by that same cochain and preserves
the obstruction and the resulting sections. This gives the
explicit comparison between all choices of affine origin.

It remains to identify the quotient on the solution space.
The geometric centralizer of the class-two target in $A$ is
trivial. Its image in $V$ would be Galois-invariant, hence
zero by Frobenius weights; its remaining central image in
$M_2$ is likewise zero. Thus any geometric conjugator between
two lifts with the same exact target lies in $J$. To preserve
normalization at $c$ it must lie in $J^{C=1}$. Conjugation
by $h^{-1}$ changes (A.42.5) by the coboundary
\[
        ((U-1)h,(B-1)h,(T-1)h).
 \tag{A.42.8}
\]
Its image under $\mathcal L_s$ is
$(\rho_\alpha(r)-1)h=0$. These are all remaining
conjugations, proving (A.42.4). Their image is compact and
closed, so this is a Hausdorff compact quotient.

The criterion also applies to lifts initially allowed
ramification at other primes. The abelian target is
unramified there, and A.19 makes the entire local
metabelian section unramified. Such a lift consequently
factors through the same arithmetic group. To change the
chosen class-two representative, lift its geometric
conjugator through $Q\to A$ and conjugate every free lift
and every actual action by that element. On $M$ this is
the explicit deck multiplication by its image in $V$.
For normalized representatives the conjugator in $A$
is $c$-fixed; choose its lift through the actual
surjection $Q^c\to A^c$, preserving the normalization
literally.
It intertwines the Fox operators, their errors, and their
solution quotients. Thus (A.42.3)--(A.42.4) characterize
section classes and not just the chosen representatives.

For finite evaluation, reduce these actual maps on a
cofinal system of arithmetic-stable finite quotients of
$J$, for example $J/(3^aJ+\mathfrak a^bJ)$. The generator
actions and the completed Fox derivatives reduce to finite
operators, and $e_s$ reduces to its actual finite relator
error. Vanishing of (A.42.3) is equivalent to solvability
of all these finite affine linear systems. The proof is
the compact finite-intersection argument in A.41, applied
to the closed solution conditions in $\mathcal D_\alpha$.
It retains compatibility across every precision and every
augmentation degree. These finite inputs must be evaluated
from the actual arithmetic word and action; the operator
description does not supply numerical coefficients that
have not been computed.

On the special family $R$, the reference cusp section is
available. The unique local conjugation of A.41 sets the
tame correction to zero and the Frobenius correction to
zero. The injective operator $1-T^{289}$ then reduces
(A.42.6) exactly to $E_W=0$, and the same conjugation
removes the quotient (A.42.8). This is the explicit map
from the general affine criterion to the image and fibers
in A.41.
\end{proof}

\begin{statement}{A.43}{Proved}{Finite arithmetic image and fiber tests with arbitrary ramification support}
Let $s$ be an exact continuous class-two section, normalized
by $s(c)=1$, and let $\alpha$ be its abelian projection.
The following construction applies to each of
\[
       G=N,\qquad G=N/N'',\qquad G=N/\gamma_3(I_{Q_c}).
 \tag{A.43.1}
\]
It gives an exact image criterion and every fiber using
explicit finite arithmetic fields, finite groups, and finite
cocycle sets. The target may have infinite ramification
support.

Put $H_0=N$, $H_{n+1}=\Phi(H_n)=H_n^3[H_n,H_n]$, with
closures understood. The finite tests below are indexed by
$n\ge1$; set
\[
\begin{split}
 P_n&=N/H_n,\qquad A_n=N/(H_n\gamma_3N),\\
 G_n&=P_n,\quad P_n/P_n'',\quad
                 P_n/\gamma_3(I_n),
\end{split}
 \tag{A.43.2}
\]
in the three cases, where $I_n$ is the normal closure of
the actual image of $c_{Q_c}$ in $P_n$. Write $s_n$ for the
exact target cocycle in $A_n$, and let
\[
 S_n=\{3,\infty\}\cup
       \{q\ne3:\alpha\bmod3^n\text{ is ramified at }q\}.
 \tag{A.43.3}
\]
These are explicitly determined finite increasing sets at
each precision. All finite geometric groups and their based
arithmetic actions are constructed by the iterated elementary
abelian covers in the proof.

Let $F_n$ be the finite Galois field splitting the action on
$G_n$. Let $K_n$ be its compositum with the finite Galois
field cut out by the target homomorphism
$g\mapsto(s_n(g),\rho_{A_n}(g))$. If $s_n$ is ramified
outside $S_n$, the solution set below is empty immediately.
Otherwise put
\[
 d_n=|G_n|[F_n:\mathbb Q],\qquad
 B_n=\prod_{p\in S_n\setminus\{\infty\}}
               p^{d_n(1+\lfloor\log_p d_n\rfloor)}.
 \tag{A.43.4}
\]
Set $\Omega_0=\mathbb Q$. Construct $\Omega_n$ as the compositum of $\Omega_{n-1}$
and all Galois extensions $E/\mathbb Q$ containing $K_n$,
unramified outside $S_n$, and of degree at most $d_n$.
The degree and discriminant bound $B_n$ make this a
terminating finite field enumeration. Put
$\Delta_n=\operatorname{Gal}(\Omega_n/\mathbb Q)$.

Define $X_n(s,G)$ by finite enumeration of all functions
$z:\Delta_n\to G_n$ satisfying
\[
\begin{split}
 z(gh)&=z(g)\,g(z(h)),\qquad z(c)=1,\\
 \pi_n z&=s_n,
\end{split}
 \tag{A.43.5}
\]
and the following finite additional conditions: for every
$k\le n$, its projection to $G_k$ is trivial on the inertia
groups at every prime of $S_n\setminus S_k$. The target and
coefficient actions in (A.43.5) are inflated through their
specified finite fields. A failed immediate test makes all
subsequent solution sets empty.

There are explicitly induced transition maps between these
finite sets, and
\[
 [s]\in\operatorname{im}
    \bigl(\operatorname{Sec}(G)\to\operatorname{Sec}(N/\gamma_3N)\bigr)
       \quad\Longleftrightarrow\quad
                 X_n(s,G)\ne\varnothing\text{ for every }n.
 \tag{A.43.6}
\]
Let $J_n=\ker(G_n\to A_n)$, with its actual $c$-action.
The entire fiber over $[s]$ is
\[
                      \varprojlim_n\bigl(X_n(s,G)/J_n^c\bigr),
 \tag{A.43.7}
\]
where $h\in J_n^c$ acts by
$z(g)\mapsto h^{-1}z(g)g(h)$. Thus the finite quotients in
the fiber use the kernels $J_n$, with no replacement by a
finite-level centralizer of the target.

Each stage is a terminating arithmetic construction from
the finite target data at that precision. Formula (A.43.6)
is an infinite sequence of finite tests; no finite stopping
rule for every positive instance is asserted.
\end{statement}
\begin{proof}
First we construct the finite geometric inputs. The subgroups
$H_n$ are characteristic, open, and cofinal in $N$. Indeed
the iterated Frattini series of every finite three-group
terminates, so every finite quotient of $N$ kills some $H_n$.
Induction also gives $x^{3^n}\in H_n$ for every $x\in N$.
The image of $H_n$ in $V$ is exactly $3^nV$: cubes
multiply its image by three and commutators have zero
image. Thus $A_n^{\mathrm{ab}}=V/3^nV$, and $S_n$ is
computed from the same finite target table as $s_n$.
Consequently every group in (A.43.2) has exponent dividing
$3^n$. Closedness of the geometric kernels and compactness
give
\[
              G=\varprojlim G_n,
       \qquad N/\gamma_3N=\varprojlim A_n.
 \tag{A.43.8}
\]
Images of closed commutator subgroups under a surjection
are the corresponding closed commutator subgroups. This
proves the displayed formulas $P_n/P_n''$ and
$P_n/\gamma_3(I_n)$ for the actual quotient maps. The marked
normal subgroup is arithmetic-stable since $Q_c$ is rational.

There is an algebraic construction of every $P_n$ with its
based arithmetic action. At the stage corresponding to
$H_n$, write the actual open cover as $Y=C\setminus D$,
where $C$ is smooth proper and $D$ is its reduced boundary.
Choose an auxiliary interior point $b\in Y(\overline{\mathbb Q})$,
construct the generalized Jacobian $J_{C,D}$, and use
\[
       j_b:Y\longrightarrow J_{C,D},\qquad
       P\longmapsto[\mathcal O_C(P-b),\text{its boundary frames}].
\]
The meromorphic section associated with $P-b$ is nonzero
on $D$ and specifies these frames. Pull back the finite
\'etale map $[3]:J_{C,D}\to J_{C,D}$ along $j_b$.
The induced map on geometric abelian three-adic homology
is an isomorphism: on the proper quotient this is the
usual Abel--Jacobi isomorphism, and on the boundary kernel
it sends each inertia generator to the corresponding
toric generator, with the same diagonal relation. The
two exact sequences and this map prove the assertion
integrally. For the proper comparison and generalized
cover construction see \cite[Chapter~III, Propositions~9.1
and~9.7]{AMilneAV}. It follows that this pullback is
connected and has kernel $H_n^3[H_n,H_n]=H_{n+1}$.

All the constructions just used are finite algebraic
operations: smooth normalization, divisor and framed-line
bundle charts, the Jacobian group law, multiplication by
three, and finite fiber products. They supply equations
for the maximal elementary abelian three-cover without
requiring an algorithm for relations among an infinite
group of divisor classes. Choosing its lift of the
prescribed tangential Puiseux point retains the based
cover. In particular no divisor $P-d$ with $d\in D$ has
been used in the generalized Jacobian. Changing the
auxiliary point to $b'$ changes $j_b$ by the translation
$[b-b']$. If $3a=[b-b']$, the actual map between the two
pullbacks is $(P,z)\mapsto(P,z+a)$. The prescribed Puiseux
lift fixes its remaining three-torsion deck ambiguity.

The arithmetic descent is specified by an actual subgroup.
Under the given tangential splitting, the stage-$n$ cover
corresponds to $H_n\rtimes G_{\mathbb Q}$ inside
$N\rtimes G_{\mathbb Q}$; in the full arithmetic
fundamental group take the inverse image of this subgroup.
Its next pointed cover uses
$\Phi(H_n)\rtimes G_{\mathbb Q}$ in the same way.
The geometric isomorphism with the $[3]$-pullback is the
one matching its prescribed tangential lift. Thus the
splitting specifies the genuine framed descent, including
the conjugations in its Galois maps. Computing a finite
normal compositum of the equations and these based maps
realizes that descent algebraically. This is also the
iterated Kummer-cover construction in C.16.
Only finitely many Puiseux
coefficients are needed to distinguish the finitely many
branches; Galois permutation of those coefficients and
the finite equations computes the based action on $P_n$.
Finite group quotient operations then compute $G_n,A_n$,
their cusp images, and the field $F_n$. Thus the coefficient
action is constructed from the curve.

We next prove the uniform ramification bound that makes
the arithmetic sets finite even for infinite support.
At every $q\ne3$, each section of any group in (A.43.1)
has either trivial tame inertia or inertia conjugate to
\[
                         c_d^{a_q},\qquad a_q\in\mathbb Z_3,
 \tag{A.43.9}
\]
at a rational local cusp. For $N$ this is B.7--B.8, and
for $N/N''$ it is A.19. For the marked quotient the two
cases are explicitly B.12 and B.14: at the marked cusp
B.12 gives exactly the inertia $c_d^{a_q}$, including when
there is an additional central Frobenius parameter; at an
unmarked cusp B.14 gives the actual cusp section. The final
zero-inertia assertion of B.12 proves that zero abelian
inertia makes the entire marked inertia trivial.

Every cusp image $e_d$ is primitive in $V$. For the eight
non-origin cusps this follows from the toric basis, and
$e_O$ is minus their sum. Thus
$a_qe_d=0$ in $V/3^n$ implies $3^n\mid a_q$.
The exponent bound on $G_n$ then kills (A.43.9). Wild
inertia at $q\ne3$ is pro-$q$ and has trivial image in a
three-group. Hence every actual section above the target
reduces to an unramified section outside $S_n$ at level
$G_n$. This proves the bound using only the abelian target
at that precision. It also proves the immediate rejection
when the finite class-two target itself ramifies outside
$S_n$.

For precision, the sets in (A.43.3) are the ramification
sets of actual finite cocycles. At $q\ne3$ the coefficient
action is unramified, so their inertia restrictions and
their vanishing are independent of a coboundary change.
The good compactification of the given curve makes every
$F_n$ unramified outside three. The fields $K_n$ are finite
and have degree at most $|A_n|[F_n:\mathbb Q]\le d_n$:
use the target homomorphism with the action quotient
$\operatorname{Gal}(F_n/\mathbb Q)$ in its second component.
After the immediate check they too are unramified outside
$S_n$.

We give the field bound in (A.43.4). For a number field
$E/\mathbb Q$ of degree at most $d_n$, unramified outside
$S_n$, pass in each completion first to its maximal
unramified subextension. For its remaining ramification
index $e$, the derivative of an Eisenstein polynomial for
a uniformizer has terms with pairwise distinct valuations
modulo $e$. The different exponent is therefore at most
$e-1+e v_p(e)$. Multiplying by residue degrees and summing
gives
\[
 v_p\operatorname{disc}(E)
           \le d_n(1+\lfloor\log_p d_n\rfloor).
 \tag{A.43.10}
\]
Taking the product proves the bound $B_n$. Apply the
finite polynomial enumeration of C.14 and its checked
Hermite bound, now at every prime of $S_n$: enumerate
monic integral polynomials of degree at most $d_n$ whose
roots have the generous bound $2(B_n+1)$, factor them,
and perform the exact degree, Galois, containment and
ramification tests. This lists every required field with
its embeddings in the fixed algebraic closure. Taking
their finite compositum with $\Omega_{n-1}$ constructs
$\Omega_n$. Decomposition and inertia subgroups are then
computed by finite number-field operations. No feasibility
estimate for this deliberately large enumeration is used.

Every continuous $G_n$-valued cocycle $z$ unramified
outside $S_n$ defines the actual finite homomorphism
\[
 g\longmapsto(z(g),\rho_{G_n}(g))
      \quad\hbox{into}\quad
      G_n\rtimes\operatorname{Gal}(F_n/\mathbb Q).
 \tag{A.43.11}
\]
Its kernel cuts out a Galois extension of degree at most
$d_n$. If its projection is $s_n$, this field contains
$K_n$. It is one of the fields used in $\Omega_n$, so
every candidate cocycle factors through $\Delta_n$.
Conversely, every function satisfying (A.43.5) inflates
to an actual global cocycle with the stated ramification.
This proves that the finite enumeration misses no candidate.
Normalization at complex conjugation and all local inertia
checks are finite group checks on these same fields.

The additional lower-level ramification conditions are
needed for the transition maps. For $z\in X_n$, its
projection to $G_k$ is unramified outside $S_k$. The
homomorphism (A.43.11) at level $k$ therefore cuts out
one of the fields included in $\Omega_k$. It follows
that the projected cocycle factors through $\Delta_k$.
It satisfies all the conditions defining $X_k$, so
projection gives the asserted transition map. This also
describes its finite computation: deflate the projected
function through the known restriction
$\Delta_n\twoheadrightarrow\Delta_k$. The transition
maps need not be onto.

Suppose first that a class above $[s]$ exists. Lift a
geometric conjugator through $G\to N/\gamma_3N$ to make
its projection the exact chosen target cocycle. Its
value at $c$ is then in the pro-three kernel of this
map. Coprime complement conjugacy in that kernel
normalizes it to one while preserving the exact target.
Its projections satisfy every ramification bound proved
above, hence give a compatible point of $\varprojlim X_n$.

Conversely, nonempty finite sets $X_n$ with the specified
transition maps have a nonempty inverse limit: in the
compact product of the finite sets, every finite list
of compatibility equations is satisfied by projection
from a common higher level. A compatible sequence of
the inflated cocycles defines a continuous cocycle into
$G=\varprojlim G_n$, because the inverse-limit topology
tests continuity in each finite quotient. Its cocycle
equation, normalization at $c$, and projection to the
exact target follow at every level. It is the required
global section. This proves (A.43.6), including the
case of an infinite union of the sets $S_n$.

We finally prove the fiber formula rather than replacing
its conjugation groups by finite approximations without
comparison. Put $J=\ker(G\to N/\gamma_3N)$. Its image in
$G_n$ is exactly $J_n$: if a lift of an element of $J_n$
has class-two image in the image of $H_n$, multiply it
by the inverse of an element of the image of $H_n$ with that same
class-two image. The resulting lift lies in $J$.
Thus $J\twoheadrightarrow J_n$, and the same argument
gives $J_n\twoheadrightarrow J_k$ for $n\ge k$.
The kernels are pro-three and $c$-stable. Coprime
complement conjugacy gives the actual surjections
\[
               J^c\twoheadrightarrow J_n^c,
       \qquad J_n^c\twoheadrightarrow J_k^c.
 \tag{A.43.12}
\]

The global geometric centralizer of the target in
$N/\gamma_3N$ is trivial. Projection to $V$ and then
to $M_2$ reduces this assertion to their zero invariant
groups, proved by Frobenius weights. Hence the global
conjugators between two lifts of the exact target lie
in $J$. The normalizations at $c$ restrict them exactly
to $J^c$. The fiber is therefore
$(\varprojlim X_n)/J^c$.

This quotient is canonically the right side of (A.43.7).
For surjectivity, choose representatives of a compatible
sequence of finite orbits recursively. A conjugator
needed to match representatives at level $n$ lifts to
level $n+1$ by (A.43.12), so they can be chosen compatible.
For injectivity, if two compatible sequences are
conjugate at every level, the corresponding sets of
conjugators in the compact group $J^c$ are nonempty,
closed, and nested, again by (A.43.12). Their intersection
gives one global conjugator. This proves (A.43.7) and
retains every fiber identification.

At each precision the target input used by the procedure
is a finite Galois presentation of its projected cocycle
and its table. All subsequent field and finite-group
operations terminate. The infinite assertion in
(A.43.6) remains an infinite assertion, in agreement
with the decision obstruction in A.12. The construction
provides finite certificates of exclusion when a stage
is empty and exact compatible finite descriptions of
all nonempty fibers.
\end{proof}

\begin{statement}{A.44}{Proved}{Inherited comparison and every marked coefficient layer}
Let $F$ be the free pro-three geometric group of
$\mathbb P^1\setminus\{0,1,\infty\}$, with the based
generators $x,y$ of A.38. The actual Fermat covering and
its specified tangential comparison identify
\[
 N=\ker\bigl(F\longrightarrow(\mathbb Z/3)^2\bigr),
 \qquad c_{Q_c}=y^3.
 \tag{A.44.1}
\]
Put $J_y=\overline{\langle\!\langle y\rangle\!\rangle}_F$
and $I=I_{Q_c}$. The inherited quotients and their actual
comparison maps are
\[
\begin{aligned}
 N/N''&\twoheadrightarrow H_{\rm met}
       :=\operatorname{im}(N\to F/F''),\\
 &\ker=(N\cap F'')/N'',\\
 N/\gamma_3I&\twoheadrightarrow H_{\rm depth}
       :=\operatorname{im}(N\to F/\gamma_3J_y),\\
 &\ker=(N\cap\gamma_3J_y)/\gamma_3I.
\end{aligned}
 \tag{A.44.2}
\]
These are continuous genuine $G_{\mathbb Q}$-equivariant
surjections, with the origins of A.38. The second is
precisely the pro-three instance of LM7 (E25.1).
At the prounipotent level over $\mathbb Q_3$ the image in each
specified quotient of the base group is the whole base
quotient. On abelian Lie algebras the map from the
intrinsic group has kernel of dimension eight.

Set $Q=N/\gamma_3I$ and
$M_r^{\rm mark}=\gamma_rQ/\gamma_{r+1}Q$.
There is an actual quotient of graded $\mathbb Z_3$-Lie
algebras
\[
 \operatorname{gr}_{\gamma}Q
   =\mathcal L(V)/\gamma_3\mathfrak j,
 \qquad \mathfrak j=(e_{Q_c})_{\mathcal L(V)}.
 \tag{A.44.3}
\]
Every graded piece is finite free. Choose the free factor
$B$ of rank nine obtained by omitting the generator
$y^3$ from the ten literal Schreier generators of A.38,
and put $V_B=B^{\rm ab}$. As a choice of integral bases,
\[
\begin{aligned}
 M_r^{\rm mark}\simeq{}&
 \mathcal L_r(V_B)\ \oplus\
 V_B^{\otimes(r-1)}e_{Q_c}\\
 &{}\oplus\left(\bigwedge\nolimits^2
       \bigl(\mathrm T(V_B)e_{Q_c}\bigr)\right)_r.
\end{aligned}
 \tag{A.44.4}
\]
Here a tensor word of length $a$ followed by $e_{Q_c}$
has ordinary degree $a+1$. The arithmetic action is the
one induced by (A.44.3) from the actual action on $V$;
(A.44.4) does not impose a block diagonal arithmetic
action. For $r\ge2$ its rank is
\[
 \frac1r\sum_{d\mid r}\mu(d)9^{r/d}
 +9^{r-1}
 +\frac{(r-1)9^{r-2}
          -\mathbf1_{2\mid r}\,9^{r/2-1}}2.
 \tag{A.44.5}
\]
In particular the ranks in degrees one through four are
$10,45,330,2466$.

For $r\ge3$, put $A_{r,q}=M_r^{\rm mark}(-1)$ and
$D_{r,q}=\operatorname{Frob}_q-1$ on that lattice.
Every $D_{r,q}$ is injective, and, with the restricted
families of F.1, the complete coefficient formulas are
\[
\begin{split}
 H^2_{\rm cts}(G_{\mathbb Q},M_r^{\rm mark})
    &=\mathcal R(A_r)/D_r\mathcal R(A_r),\\
 \varprojlim{}^1_n H^1(G_{\mathbb Q},M_r^{\rm mark}/3^n)
    &=\frac{\mathcal R(A_r)\cap
                   \prod_{q\ne3}D_{r,q}A_{r,q}}
                 {D_r\mathcal R(A_r)}\ne0,\\
 H^1_{\rm cts}(G_{\mathbb Q},M_r^{\rm mark})
    &=M_r^{\rm mark}/(u-1)(M_r^{\rm mark})^c.
\end{split}
 \tag{A.44.6}
\]
The quotient functor (A.44.3), applied to the matrices
of A.7, gives all action matrices and hence an integral
Smith calculation of the last group at every degree.
Degrees one and two are $V$ and $\bigwedge^2V$; their
derived limits are those of F.5 and are also nonzero.
Together with A.21--A.22, these formulas specify the
cofinal finite-rank towers and the derived-limit term
of every layer required in Task A.6. Their exact images
and all fibers are the finite arithmetic sets of A.43.
\end{statement}
\begin{proof}
The morphism of curves is $t=-z^3/x^3$.
The two cube roots are $-z/x$ and $-y/x$, so its deck
quotient is the indicated $(\mathbb Z/3)^2$. The based
pullback $t^{1/3^n}\mapsto-s^{3/3^n}$ of A.38 identifies
the arithmetic splittings, not just the outer actions.
The Schreier loop $y^3$ is the marked cusp loop, with
the same chosen path. Thus $I\subset J_y$, and
$N''\subset F''$. Passing to these closed subgroups
proves the two kernels in (A.44.2). Compactness makes
their images closed, so the displayed quotient maps
are continuous surjections. Their arithmetic
equivariance follows from this based covering map.
They therefore also induce the specified pushforward
maps on section classes.

For the rational prounipotent assertion, use the
continuous morphism from the $\mathbb Q_3$-prounipotent
completion of $N$ to that of $F$. Its image contains
$x^3$ and $y^3$. In every characteristic-zero nilpotent
quotient, their logarithms are $3\log x$ and
$3\log y$, so the image Lie algebra contains both
$\log x$ and $\log y$. More explicitly, the continuous
free-Lie homomorphism
\[
 \log x\longmapsto\tfrac13\log_N(x^3),\qquad
 \log y\longmapsto\tfrac13\log_N(y^3)
\]
is a section of this morphism on the completed Lie
algebras. It maps the second derived ideal into the
second derived ideal. It also maps the ideal generated
by $\log y$ into the intrinsic ideal generated by
$\log_N(c_{Q_c})$, since $c_{Q_c}=y^3$.
It therefore descends through both specified quotients
and proves their inverse-limit surjectivity directly.
Their depth and derived
ideals have zero image in the two-dimensional
abelianization. The source has ten-dimensional
abelianization, and the displayed logarithms prove
surjectivity there, giving kernel dimension eight.
This argument concerns the specified prounipotent
image of the covering map; it does not substitute
completion for either integral kernel in (A.44.2).

We now prove the integral layer statement. The actual
Schreier basis gives $N=B*_3\langle c_{Q_c}\rangle$.
The kernel $I$ of its retraction to $B$ is the free
pro-three group on the continuous family
$b c_{Q_c}b^{-1}$ indexed by $b\in B$. One can verify
this directly: the free pro-three group on that
profinite family, with $B$ acting by left translation,
forms a semidirect product with $B$. Sending its
named generators to $b c_{Q_c}b^{-1}$ and sending
$B$ identically to $B$ defines a homomorphism to $N$.
The inverse sends $c_{Q_c}$ to the generator indexed
by $1$ and sends the nine generators of $B$ to
themselves. The two compositions are the identity
on their topological generating families.

Let
\[
\begin{aligned}
 \Lambda_B&=\mathbb Z_3[[B]],&
 \mathfrak a_B&=\ker(\Lambda_B\to\mathbb Z_3),\\
 M&=\Lambda_B e,&
 Z&=\widehat{\bigwedge}^{\,2}_{\mathbb Z_3}M.
\end{aligned}
 \tag{A.44.7}
\]
The last completion is over finite coefficient and
finite label quotients. The free class-two quotient
of $I$ is the group on $M\oplus Z$ with law
\[
 (m,z)(m',z')=
       (m+m',z+z'+\tfrac12m\wedge m').
 \tag{A.44.8}
\]
Its named cusp conjugate is $(be,0)$, and $B$ acts
by left multiplication on $M$ and its induced
exterior-square action on $Z$. The coefficient
$1/2$ is an integral three-adic unit. The class-two
collection formula first proves this model for
finite generating sets and then, through finite
partitions of $B$, for the completed family.
The two universal maps just constructed give the
actual continuous isomorphism
\[
       Q\simeq B\ltimes(M\oplus Z),\qquad
       \bar I:=I/\gamma_3I\simeq M\oplus Z.
 \tag{A.44.9}
\]
No arithmetic stability of the chosen free factor
$B$ is required for this geometric isomorphism.

For $r\ge1$ define
\[
\begin{aligned}
 M_r^*&=\mathfrak a_B^{r-1}M,\\
 Z_r^*&=\overline{\sum_{\substack{a,b\ge0\\a+b\ge r-2}}
               \mathfrak a_B^aM\wedge\mathfrak a_B^bM}.
\end{aligned}
 \tag{A.44.10}
\]
and set $G_r=\gamma_rB\ltimes(M_r^*\oplus Z_r^*)$.
We prove $G_r=\gamma_rQ$, including the closures.
The integral Magnus map identifies
$\Lambda_B$ with the noncommutative formal-series
ring on nine variables, and maps
$\gamma_rB$ into $1+\mathfrak a_B^r$.
Thus $\gamma_rB$ raises the augmentation degree
on $M$ by at least $r$, and raises total degree
on a wedge by at least $r$. Formula (A.44.8)
gives $[G_r,G_s]\subset G_{r+s}$. Since $G_1=Q$,
this proves $\gamma_rQ\subset G_r$.

For the reverse inclusion, the projection of a
commutator with $b\in B$ to $M$ is multiplication
by $b-1$. Iterating such commutators on the cusp
conjugates shows that
\[
       \gamma_{a+1}Q\cap\bar I\longrightarrow
                     \mathfrak a_B^aM
 \quad\hbox{is onto}.
 \tag{A.44.11}
\]
Indeed these commutators give every monomial in
the nine augmentation generators times every
$be$; their closed linear span is the indicated
ideal. The image is closed because the source
is compact. Bracketing a lift of
$m\in\mathfrak a_B^aM$ with a lift of
$m'\in\mathfrak a_B^bM$ gives exactly
$(0,m\wedge m')$ by (A.44.8). Hence
$Z_r^*\subset\gamma_rQ$. A lift in (A.44.11)
of any $m\in M_r^*$ has central component in
$Z_r^*$ by the already proved inclusion.
Remove that component to obtain $(m,0)$ in
$\gamma_rQ$. Finally the embedded
$\gamma_rB$ belongs to $\gamma_rQ$.
These facts prove equality with $G_r$.

Magnus monomials now give
$\operatorname{gr}_{\mathfrak a_B}M=
\mathrm T(V_B)e$. On $Z$ their unordered distinct
pairs give an integral basis at every total degree.
The completed product introduces no relations:
at any fixed degree there are only finitely many
such words and pairs, and the coefficient
projections are the actual continuous inverse-limit
maps. This proves (A.44.4) and finite freeness.

For completeness the comparison with the free
Lie quotient is explicit. The free Lie algebra
$\mathcal L(V_B)$ acts on
$\mathcal L(\mathrm T(V_B)e)$ by extending left
concatenation as derivations. There are inverse
Lie homomorphisms
\[
 \mathcal L(V_B\oplus\mathbb Z_3e)
   \ \rightleftarrows\
 \mathcal L(V_B)\ltimes
       \mathcal L(\mathrm T(V_B)e).
 \tag{A.44.12}
\]
One sends $V_B,e$ to the named generators. The
other sends $v_1\cdots v_a e$ to
$\operatorname{ad}(v_1)\cdots
\operatorname{ad}(v_a)e$ and is the identity
on $V_B$. The universal property of the free
Lie algebra extends these maps, and their
composition is the identity on $V_B,e$ and
all their iterated actions. Quotienting the
second free Lie factor by its third lower
central term leaves
$\mathrm T(V_B)e\oplus
\bigwedge^2(\mathrm T(V_B)e)$.
Under the actual maps in (A.44.11), this is
exactly the graded group just computed.
It proves (A.44.3) integrally. Equivariance
uses the actual action on $V$ and the
arithmetic-stable cusp line
$\mathbb Z_3e_{Q_c}$, so it also proves the
stated prescription for arithmetic matrices.

The first summand of (A.44.4) has the Witt
rank in (A.44.5), and its second has rank
$9^{r-1}$. For the third, there are
$(r-1)9^{r-2}$ ordered pairs of tensor words
of total ordinary degree $r$. An equal pair
exists only for even $r$, with
$9^{r/2-1}$ choices. Subtract these pairs
and divide by two to count unordered distinct
pairs. This proves (A.44.5).
The program \texttt{a\_marked\_layers.py} in
\texttt{code/} independently enumerates all
ordinary Lyndon words on ten letters through
degree five, retains marked-letter counts
zero, one and two, and checks all three
summands and the ranks $45,330,2466$.
Its exact output and JSON certificate are saved.

Finally, over $K_0$ the finite coefficient
representations of these graded modules have
three-group image, since they are functorial
quotients of tensor powers of $V$. The
composition-series argument of A.1 gives
$H^i(G_{\mathbb Q,S},M_r^{\rm mark}/3^n)=0$
for $i\ge2$ and $S=\{3,\infty\}$.
The normalized free-source calculation of
A.10 gives the complex
$[(M_r^{\rm mark})^c\xrightarrow{u-1}
M_r^{\rm mark}]$.
In degree $r\ge3$, all Frobenius weights
after twisting by $-1$ have absolute value
at least $q^{r/2-1}>1$; hence $D_{r,q}$
is injective. F.1--F.3 therefore give every
formula in (A.44.6). The nonzero free
summand $\mathcal L_r(V_B)$ makes each
coefficient lattice nonzero. Chebotarev at
splitting primes of the reduction of
$M_r^{\rm mark}(-1)$ modulo $3^n$
makes all entries of $D_{r,q}$
divisible by arbitrarily large powers of
three. F.3 then proves the strict nonvanishing
of its derived limit. This also proves the
same nonvanishing after every finite
number-field extension, as in F.5.

The finite-rank quotients are the actual
$N/(\gamma_3I\,\gamma_rN)$ of A.21,
and their coefficient kernels are precisely
the free lattices just computed. The
metabelian tower uses the actual quotients
$N/(N''\gamma_rN)$ and the lattices of
A.22. A.21 proves their cofinality and the
comparison on section classes. At every
layer the residue families in (A.22.2) or
(A.44.6) retain the integrality and
restricted-convergence requirements of
F.1. The finite sets of A.43 additionally
retain all inherited ramification conditions;
their compact inverse limit gives the exact
global image and fiber without omitting any
of these layer obstructions.
\end{proof}
